% Local macros (provide-command idempotent; required for cross-chapter
% references to thm:koszul-reflection in theorem_A_infinity_2.tex).
\providecommand{\Kosz}{\mathsf{Kosz}}
\providecommand{\co}{\mathrm{co}}

\bigskip
\noindent\textbf{Reconstruction and quadratic recognition.}\enspace
Let \(A\) be a complete augmented chiral algebra and let
\[
B_X(A)=T^c(s^{-1}\bar A),\qquad
\varepsilon_A\colon\Omega_X B_X(A)\longrightarrow A
\]
be its ordered bar coalgebra and the algebra-side counit.  In the
enhanced associative Ran ambient, the bar--cobar equivalence is the
universal reconstruction theorem of Francis--Gaitsgory
\cite[Proposition~4.1.2 and the proof of Theorem~5.1.1]{Francis2012};
its restriction to the factorization and completed
Fulton--MacPherson models uses the comparison package
\(H_{\mathrm{CL}}\) of
Definition~\ref{def:chiral-comparison-lemma-package}.

A quadratic presentation \(A=T(V)/(R)\) supplies a second coalgebra,
\[
A^{\mathrm i}=C(s^{-1}V,s^{-2}R)\subset T^c(s^{-1}V),
\]
and a canonical comparison
\[
q_A\colon A^{\mathrm i}\longrightarrow B_X(A).
\]
Theorem~B recognizes the locus on which \(q_A\) is a
quasi-isomorphism, equivalently the locus on which
\(\Omega_X(A^{\mathrm i})\to A\) is a quasi-isomorphism.  This is the
quadratic Koszul theorem of Loday--Vallette
\cite[Theorems~2.3.2 and~3.4.6]{LV12}, transported to the chiral
factorization surface by \(H_{\mathrm{CL}}\).

For a fixed chiral CDG-coalgebra \(C\), the co--contra equivalence
\[
D^{\mathrm{co}}(C\text{-}\mathrm{CoFact})
\simeq
D^{\mathrm{ctr}}(C\text{-}\mathrm{ContraFact})
\]
is the separate second-kind statement of
Theorem~\ref{thm:off-koszul-ran-inversion}.  An algebra comparison on
that surface is governed by the acyclic twisting package
\(\mathsf{Tw}^{\mathrm{ch}}_{\mathrm{acyc}}(C,A,\tau)\).
These three theorem lanes have distinct sources, targets, and
hypothesis packages.

\section{Non-quadratic chiral algebras}
\label{sec:filtered-vs-curved-comprehensive}

\label{sec:bar-convergence-overview}% see thm:bar-convergence below for the canonical theorem
Convention~\ref{conv:regime-tags} records four independent features of
a bar model: a quadratic presentation, a curved modular enhancement, a
chosen completion topology, and an infinite generator tower.  Theorem~A
supplies universal reconstruction in the enhanced associative Ran
ambient.  Theorem~B studies the quadratic comparison $q_A$.  A curved
enhancement uses the fixed-coalgebra second-kind theorem.  A completed
or infinite-type realization requires finite-window comparison maps,
Mittag--Leffler control, and continuity of every structure map.  The
examples of these features are given in
\S\ref{subsec:three-regimes} above.

\subsection{Convergence criteria}

\begin{theorem}[Finite-window convergence of the bar construction;
\ClaimStatusProvedHere]\label{thm:bar-convergence}
Let $\mathcal{A}$ be a chiral algebra and write
$C=\bar{B}(\mathcal{A})=\bigoplus_{n\ge0}C^n$ with finite-window
subcomplexes
\[
F_{\le p}C:=\bigoplus_{0\le n\le p}C^n
\]
and completion
$\widehat C=\varprojlim_p C/F_{>p}C\simeq\prod_{n\ge0}C^n$.
Assume:
\begin{enumerate}[label=\textup{(\roman*)}]
\item the bar differential is locally finite and preserves finite
 windows, so $F_{\le p}C$ is a subcomplex for every~$p$;
\item each finite window is finite-dimensional in every cohomological
 degree;
\item for every cohomological degree~$m$, the inverse system
 $\{H^m(C/F_{>p}C)\}_p$ is Mittag--Leffler;
\item the finite-support comparison
 \[
 H^m(C)\longrightarrow \varprojlim_p H^m(C/F_{>p}C)
 \]
 is an isomorphism for every~$m$.
\end{enumerate}
Then the natural map $C\to\widehat C$ is a quasi-isomorphism.
Condition~\textup{(iv)} is automatic when each cohomological degree has
a uniform bar-window bound. If either the Mittag--Leffler condition or
the finite-support comparison fails, the completed bar complex
$\widehat{\bar B}(\mathcal A)$ is the correct object; the raw
direct-sum bar complex does not carry the ordinary comparison.
\end{theorem}

\begin{proof}
The locally finite differential makes the finite-window quotients
$C/F_{>p}C$ into an inverse system of complexes and identifies
$\widehat C$ with its inverse limit. The Milnor exact sequence gives
\[
0\to\varprojlim\nolimits^1_p H^{m-1}(C/F_{>p}C)
\to H^m(\widehat C)
\to\varprojlim_p H^m(C/F_{>p}C)\to0.
\]
The first term vanishes by the Mittag--Leffler hypothesis. Hence
$H^m(\widehat C)\cong\varprojlim_p H^m(C/F_{>p}C)$, and condition~(iv)
identifies this group with $H^m(C)$. This is precisely the assertion
that $C\to\widehat C$ is a quasi-isomorphism.

In the quadratic finite-type case, every element has finite bar length
and each cohomological degree lies in a finite bar window; condition~(iv)
is automatic. In class-$\mathsf M$ examples the raw direct sum fails
this finite-window test, while the pro-object
$\varprojlim_p C/F_{>p}C$ is built to satisfy the inverse-limit
hypotheses.
\end{proof}

\begin{remark}[Finite-window input]
Theorem~\ref{thm:conilpotency-convergence} supplies local
conilpotency. The filtered-to-curved comparison
(Theorem~\ref{thm:filtered-to-curved}) and the completed Koszul
framework of Gui--Li--Zeng~\cite{GLZ22} supply the finite-window
models in non-quadratic regimes. None of these inputs replaces the
Mittag--Leffler and finite-support conditions in
Theorem~\ref{thm:bar-convergence}.
\end{remark}

\subsection{Collision codimension and completion}

\begin{remark}[Finite-window test before Verdier duality]\label{rem:finite-window-before-verdier}
Quadratic finite-type algebras live on the raw conilpotent bar
coalgebra. Curved-central algebras require a square-zero total model or
the second-kind coderived category. Filtered algebras require the
finite-window completion. Infinite-generator algebras require a strong
completion tower, as in Theorem~\ref{thm:completed-bar-cobar-strong}.
Geometrically, the filtration records the codimension of collision
loci: quadratic algebras see pairwise collisions, curved-central
algebras see central contributions from higher collisions, and filtered
algebras retain non-descendant composite fields in higher windows. The
completed bar complex is the formal neighbourhood of the diagonal in
configuration space.
\end{remark}

\section{Chiral coalgebra homological algebra}
\label{sec:chiral-coalgebra-homalg}
\index{chiral coalgebra!homological algebra|textbf}
\index{Positselski!chiral lift|textbf}

The bar complex $\bar{B}^{\mathrm{ch}}(\cA)$ is a chiral coalgebra
(Theorem~\ref{thm:bar-chiral}). At genus~$0$, it is a
DG chiral coalgebra ($\dzero^2 = 0$); at genus~$g \geq 1$ with
$\kappa(\cA) \neq 0$, it is a \emph{curved} DG chiral coalgebra
($\dfib^{\,2} = \mcurv{g} \neq 0$;
Convention~\ref{conv:higher-genus-differentials}). In the curved case,
ordinary homological algebra is not the ambient: under the finite-window
and second-kind hypotheses,
Proposition~\ref{prop:curved-bar-acyclicity} makes the bar-cobar
comparison cone coacyclic. It does not make the curved bar coalgebra
itself zero.

The chiral analogues of Positselski's coalgebra homological
algebra~\cite{Positselski11} replace the ground field~$k$ by
$\mathcal{D}_X$ and the tensor product by $\chirtensor$; all
exactness properties are inherited from configuration space geometry.

\begin{remark}[Coderived scope and the genus parameter]
\label{rem:coderived-scope}
At genus zero, the universal bar construction has square-zero
differential.  The reconstruction map
\(\Omega_XB_X(A)\to A\) belongs to Theorem~A, while the quadratic
comparison \(A^{\mathrm i}\to B_X(A)\) belongs to Theorem~B.

A genus-\(g\) modular enhancement may produce a chiral
CDG-coalgebra \(C_g=(C,d,h_g)\) with curvature \(h_g\).
Its intrinsic second-kind categories are
\[
D^{\mathrm{co}}(C_g\text{-}\mathrm{CoFact}),
\qquad
D^{\mathrm{ctr}}(C_g\text{-}\mathrm{ContraFact}).
\]
Theorem~\ref{thm:off-koszul-ran-inversion} compares these two
categories for the same fixed coalgebra under
\(\mathsf{Pos}^{\mathrm{ch}}_{\mathrm{co-ctr}}(C_g)\).
A twisting morphism \(\tau_g\colon C_g\to A_g\) produces an
algebra/module comparison after the independent package
\(\mathsf{Tw}^{\mathrm{ch}}_{\mathrm{acyc}}(C_g,A_g,\tau_g)\)
has been supplied.  Thus curvature, quadratic recognition, and
co--contra correspondence occupy three typed theorem surfaces.
\end{remark}

\subsection{Chiral CDG-coalgebras}
\label{subsec:chiral-CDG-coalgebra}

\begin{definition}[Chiral CDG-coalgebra]\label{def:chiral-CDG-coalgebra}
\index{CDG-coalgebra!chiral}
A \emph{chiral CDG-coalgebra} over a smooth curve $X$ is a triple
$(C, d, h)$ where:
\begin{enumerate}
\item $C$ is a graded chiral coalgebra on $X$: a graded
 $\mathcal{D}_X$-module equipped with a comultiplication
 $\Delta\colon C \to C \chirtensor C$
 (defined via configuration space splitting maps,
 cf.\ Theorem~\ref{thm:bar-chiral}),
 a counit $\varepsilon\colon C \to \omega_X$, coassociativity,
 and the counit axioms.
\item $d\colon C \to C$ is a homogeneous endomorphism of degree~$1$
 that is an odd coderivation: the composition
 $\Delta \circ d = (d \otimes \mathrm{id} +
 \mathrm{id} \otimes d) \circ \Delta$
 holds, and $\varepsilon \circ d = 0$ on $C^{>-2}$
 (the counit annihilates $d$ except possibly in the lowest degree).
\item $h\colon C \to \omega_X$ is a homogeneous linear functional
 of degree~$2$ (the \emph{curvature}), satisfying:
 \begin{enumerate}
 \item $d^2(c) = h \ast c - c \ast h$ for all $c \in C$,
 where $h \ast c$ denotes the coaction of the curvature
 (defined via the comodule structure of $C$ over itself);
 \item $h \circ d(c) = 0$ for all $c \in C$
 (the curvature is a cocycle for the coderivation~$d$).
 \end{enumerate}
\end{enumerate}
When $h = 0$, we recover an ordinary DG chiral coalgebra.
\end{definition}

\begin{example}[The bar complex as chiral CDG-coalgebra]
\label{ex:bar-as-CDG}
\index{bar complex!as CDG-coalgebra}
For a Koszul chiral algebra $\cA$ on $X$:
\begin{itemize}
\item At genus~$0$: $\bar{B}^{(0)}(\cA) = (\bar{B}^{\mathrm{ch}}(\cA), d_{\bar{B}}, 0)$
 is a DG chiral coalgebra ($h = 0$).
\item At genus~$g \geq 1$: $\bar{B}^{(g)}(\cA) = (\bar{B}^{\mathrm{ch}}(\cA), d_{\bar{B}}^{(g)}, h_g)$
 is a CDG chiral coalgebra with curvature
 $h_g = m_0^{(g)} = \kappa(\cA) \cdot \lambda_g$
 (Theorem~\ref{thm:genus-universality}), where $\kappa(\cA)$ is
 the obstruction coefficient and $\lambda_g$ is the Faber--Pandharipande
 class on $\overline{\mathcal{M}}_g$.
\end{itemize}
The CDG axioms are checked as follows. Axiom (2a): the identity
$d^2 = [m_0, -]$ is the curved $A_\infty$ relation
(Corollary~\ref{cor:critical-level-universality}).
Axiom (2b): $h \circ d = 0$ follows from $m_0$ being a scalar
(central element), so $d(m_0) = 0$.
\end{example}

\begin{remark}[Chiral CDG-coalgebra morphisms]
\label{rem:CDG-morphisms}
A morphism of chiral CDG-coalgebras $(C, d_C, h_C)$
$\to$ $(D, d_D, h_D)$ is a pair $(f, a)$ where
$f\colon C \to D$ is a morphism of graded chiral
coalgebras and $a\colon C \to \omega_X$
is a homogeneous linear function of degree~$1$ such
that
\begin{align*}
d_D(f(c)) &= f(d_C(c)) + f(a \ast c) - (-1)^{|c|}f(c \ast a), \\
h_D(f(c)) &= h_C(c) + a(d_C(c)) + a^2(c)
\end{align*}
for all $c \in C$.
Composition of morphisms is defined by
$(g, b) \circ (f, a) = (g \circ f,\allowbreak b \circ f + a)$.
This follows Positselski~\cite[\S4.1]{Positselski11}.
\end{remark}

\subsection{Chiral CDG-comodules and CDG-contramodules}
\label{subsec:chiral-CDG-comod-contra}

\begin{definition}[Chiral CDG-comodule; \ClaimStatusDefinitional]\label{def:chiral-CDG-comodule}
\index{CDG-comodule!chiral}
Let $(C, d, h)$ be a chiral CDG-coalgebra. A \emph{left chiral
CDG-comodule} over $C$ is a graded left $C$-comodule $M$
(i.e., a graded $\mathcal{D}_X$-module with coaction
$\lambda\colon M \to C \chirtensor M$ satisfying coassociativity
and counit axioms) endowed with a homogeneous endomorphism
$d_M\colon M \to M$ of degree~$1$ such that:
\begin{enumerate}
\item The coaction commutes with differentials:
 $\lambda \circ d_M = (d \otimes \mathrm{id} + \mathrm{id} \otimes d_M) \circ \lambda$.
\item The curvature relation:
 $d_M^2(x) = h \ast x$ for all $x \in M$,
 where $h \ast x$ denotes the action of the curvature
 via the coaction map.
\end{enumerate}

A \emph{right chiral CDG-comodule} is defined analogously with
coaction $\rho\colon N \to N \chirtensor C$ and
$d_N^2(y) = -y \ast h$.

Morphisms of chiral CDG-comodules are homogeneous $\mathcal{D}_X$-linear
maps of degree~$0$ that are compatible with the coaction
and commute with the differential (in the usual graded sense).
\end{definition}

\begin{definition}[Chiral CDG-contramodule; \ClaimStatusDefinitional]\label{def:chiral-CDG-contramodule}
\index{CDG-contramodule!chiral}
Let $(C, d, h)$ be a chiral CDG-coalgebra. A \emph{left chiral
CDG-contramodule} over $C$ is a graded left $C$-contramodule $P$
(i.e., a graded $\mathcal{D}_X$-module with contraaction
$\pi_P\colon \mathrm{Hom}^{\mathrm{ch}}(C, P) \to P$
satisfying the contraassociativity and contraunit axioms)
endowed with a
homogeneous endomorphism $d_P\colon P \to P$ of degree~$1$ such that:
\begin{enumerate}
\item The contraaction commutes with differentials on
 $\mathrm{Hom}^{\mathrm{ch}}(C, P)$ and $P$.
\item The curvature relation:
 $d_P^2(p) = h \ast p$ for all $p \in P$,
 where $h \ast p = \pi_P(h \otimes p)$ denotes the contraaction
 of the curvature.
\end{enumerate}
Morphisms of chiral CDG-contramodules are defined analogously.
\end{definition}

\begin{remark}[Contramodules in the chiral setting]
\label{rem:chiral-contramodules}
\index{contramodule!chiral}
The chiral contraaction $\pi_P\colon \mathrm{Hom}^{\mathrm{ch}}(C, P) \to P$
replaces Positselski's $\mathrm{Hom}_k(C, P) \to P$, where
$\mathrm{Hom}^{\mathrm{ch}}(C, P) := \mathrm{Hom}_{\mathcal{D}_X}(C, P)$
is the internal Hom in the chiral category. For a conilpotent chiral
coalgebra $C$ with finite-dimensional graded pieces (as holds for
$\bar{B}^{\mathrm{ch}}(\cA)$ when $\cA$ is finitely generated over $\mathcal{D}_X$),
$\mathrm{Hom}^{\mathrm{ch}}(C, P) \cong C^* \chirtensor P$
where $C^* = \mathrm{Hom}_{\mathcal{D}_X}(C, \omega_X)$ is the
graded dual. In this finite-type case, contramodules over $C$ are
equivalent to complete modules over the dual algebra $C^*$, recovering
the correspondence between conilpotent comodules and complete modules
already used in Theorem~\ref{thm:e1-module-koszul-duality}.
\end{remark}

\begin{remark}[No cohomology for CDG-comodules]
\label{rem:no-CDG-cohomology}
Since $d^2 \neq 0$ in a CDG-comodule (when $h \neq 0$), there is no
notion of cohomology, so the ordinary
derived category is inadequate at higher genus: one cannot define
acyclic objects via vanishing cohomology. The correct
substitute is the coderived category (Definition~\ref{def:chiral-coderived}).
\end{remark}

\subsection{Chiral cotensor and cohomomorphism complexes}
\label{subsec:chiral-cotensor}

\begin{definition}[Chiral cotensor product]\label{def:chiral-cotensor}
\index{cotensor product!chiral}
Let $(C, d, h)$ be a chiral CDG-coalgebra, $N$ a right chiral
CDG-comodule, and $M$ a left chiral CDG-comodule. The
\emph{chiral cotensor product} $N \square_C^{\mathrm{ch}} M$ is
the equalizer:
\[
N \square_C^{\mathrm{ch}} M =
\ker\bigl(\rho \chirtensor \mathrm{id} - \mathrm{id} \chirtensor \lambda\colon
N \chirtensor M \rightrightarrows N \chirtensor C \chirtensor M\bigr)
\]
in the category of graded $\mathcal{D}_X$-modules.
The differential on $N \chirtensor M$ restricts to
$N \square_C^{\mathrm{ch}} M$ (this uses the compatibility of coactions
with differentials), and the restricted differential has
$d^2 = h_N \ast (-) - (-) \ast h_M$ on the cotensor product, making
$N \square_C^{\mathrm{ch}} M$ a complex (i.e., $d^2 = 0$) precisely
when $h_N = h_M$ (the curvatures match), which holds when $N$ and $M$
are CDG-comodules over the same CDG-coalgebra~$C$.
\end{definition}

\begin{definition}[Chiral cohomomorphisms]\label{def:chiral-cohom}
\index{cohomomorphism!chiral}
For left chiral CDG-comodules $L$ and $M$ over $C$, the
\emph{chiral cohomomorphism complex} is:
\[
\mathrm{Cohom}_C^{\mathrm{ch}}(L, M) =
\{f \in \mathrm{Hom}_{\mathcal{D}_X}(L, M) :
\lambda_M \circ f = (\mathrm{id}_C \chirtensor f) \circ \lambda_L\}
\]
with differential $(d f)(x) = d_M(f(x)) - (-1)^{|f|} f(d_L(x))$.
This differential has a zero square (using the CDG-comodule axioms),
making $\mathrm{Cohom}_C^{\mathrm{ch}}(L, M)$ a genuine complex.
\end{definition}

\begin{definition}[Chiral contratensor product]\label{def:chiral-contratensor}
\index{contratensor product!chiral}
For a right chiral CDG-comodule $N$ and a left chiral
CDG-contramodule $P$ over $C$, the \emph{chiral contratensor
product} $N \odot_C^{\mathrm{ch}} P$ is the coequalizer:
\[
N \odot_C^{\mathrm{ch}} P =
\mathrm{coker}\bigl(
N \chirtensor \mathrm{Hom}^{\mathrm{ch}}(C, P) \rightrightarrows N \chirtensor P
\bigr)
\]
where the two maps are
$\rho \chirtensor \pi_P$ and $\mathrm{id}_N \chirtensor \pi_P$.
The induced differential on $N \odot_C^{\mathrm{ch}} P$ has zero
square by the same curvature cancellation as for the cotensor product.
\end{definition}

\begin{lemma}[Adjunction; \ClaimStatusProvedHere]\label{lem:chiral-co-contra-adjunction}
\index{adjunction!chiral cotensor--contratensor}
There is a canonical isomorphism of complexes:
\[
\mathrm{Hom}_{\mathcal{D}_X}(N \square_C^{\mathrm{ch}} M,\, V)
\;\cong\;
\mathrm{Cohom}_C^{\mathrm{ch}}\bigl(M,\,
\mathrm{Hom}_{\mathcal{D}_X}(N, V)\bigr)
\]
for any graded $\mathcal{D}_X$-module $V$, and dually:
\[
\mathrm{Hom}_{\mathcal{D}_X}(C \chirtensor V,\, P)
\;\cong\;
\mathrm{Hom}_{\mathcal{D}_X}(V,\,
\mathrm{Hom}^{\mathrm{ch}}(C, P))
\]
establishing $\Phi_C \dashv \Psi_C$
(see Definition~\textup{\ref{def:chiral-Phi-Psi}} below).
\end{lemma}

\begin{proof}
Both isomorphisms are instances of the tensor-hom adjunction in the
chiral category. For the first: an element of the left side is a
$\mathcal{D}_X$-linear map $f\colon N \square_C^{\mathrm{ch}} M \to V$.
By the universal property of the equalizer, this is equivalent to a
$\mathcal{D}_X$-linear map $\tilde{f}\colon N \chirtensor M \to V$
that equalizes the two coaction maps. By the chiral tensor-hom
adjunction (which holds for holonomic $\mathcal{D}_X$-modules
by~\cite[Proposition~3.4.6]{BD04}), $\tilde{f}$ corresponds to a map
$M \to \mathrm{Hom}_{\mathcal{D}_X}(N, V)$ that is
$C$-colinear, i.e., an element of
$\mathrm{Cohom}_C^{\mathrm{ch}}(M, \mathrm{Hom}_{\mathcal{D}_X}(N, V))$.
The second isomorphism is analogous.
\end{proof}

\subsection{Coderived and contraderived categories}
\label{subsec:chiral-coderived-contraderived}

\begin{definition}[Coderived and contraderived categories]
\label{def:chiral-coderived}
\index{coderived category!chiral|textbf}
\index{contraderived category!chiral|textbf}
Let $(C, d, h)$ be a chiral CDG-coalgebra.

\emph{Exact triples.}
An exact triple of left chiral CDG-comodules is a short sequence
$0 \to L \to M \to N \to 0$ that is exact in the abelian category
of graded $C$-comodules (with closed, degree-zero morphisms).
The analogous notion applies to CDG-contramodules.

\emph{Coacyclic objects.}
A left chiral CDG-comodule $M$ is \emph{coacyclic} if it belongs to
the minimal thick subcategory of $\mathrm{Hot}(C\text{-}\mathrm{comod}^{\mathrm{ch}})$
that contains the totalizations of all exact triples and is closed
under infinite direct sums. Denote by
$\mathrm{Acycl}^{\mathrm{co}}(C\text{-}\mathrm{comod}^{\mathrm{ch}})$
the thick subcategory of coacyclic CDG-comodules.

\emph{Contraacyclic objects.}
A left chiral CDG-contramodule $P$ is \emph{contraacyclic} if it
belongs to the minimal thick subcategory of
$\mathrm{Hot}(C\text{-}\mathrm{contra}^{\mathrm{ch}})$
that contains the totalizations of all exact triples and is closed
under infinite direct products. Denote by
$\mathrm{Acycl}^{\mathrm{ctr}}(C\text{-}\mathrm{contra}^{\mathrm{ch}})$
the thick subcategory of contraacyclic CDG-contramodules.

\emph{The exotic derived categories.}
\begin{align}
D^{\mathrm{co}}(C\text{-}\mathrm{comod}^{\mathrm{ch}})
&:= \mathrm{Hot}(C\text{-}\mathrm{comod}^{\mathrm{ch}})
\,/\, \mathrm{Acycl}^{\mathrm{co}}(C\text{-}\mathrm{comod}^{\mathrm{ch}})
\label{eq:chiral-coderived}\\
D^{\mathrm{ctr}}(C\text{-}\mathrm{contra}^{\mathrm{ch}})
&:= \mathrm{Hot}(C\text{-}\mathrm{contra}^{\mathrm{ch}})
\,/\, \mathrm{Acycl}^{\mathrm{ctr}}(C\text{-}\mathrm{contra}^{\mathrm{ch}})
\label{eq:chiral-contraderived}
\end{align}
\end{definition}

\begin{remark}[Comparison with ordinary derived category]
\label{rem:exotic-vs-ordinary}
When $h = 0$ (genus~$0$), CDG-comodules are ordinary DG-comodules
and one can also form the standard derived category
$D(C\text{-}\mathrm{comod}^{\mathrm{ch}})$
by localizing at acyclic objects (those with $H^*(M) = 0$). In general,
$\mathrm{Acycl}^{\mathrm{co}} \supset \mathrm{Acycl}$,
so $D^{\mathrm{co}}$ is a quotient of $D$. When $C$ has
finite homological dimension (e.g., when $\cA$ is Koszul with finite
global dimension), the two categories coincide
(Positselski~\cite[\S4.5]{Positselski11}). When $h \neq 0$, the standard
cohomological derived category is no longer the right invariant
for the curved theory; on the higher-genus bar-complex surfaces
treated here the relevant bar-cobar comparison cones are coacyclic in
the second-kind sense. This does not make the curved bar coalgebra
itself zero; $D^{\mathrm{co}}$ retains nontrivial comodule information.
\end{remark}

\subsection{Injective comodules and projective contramodules}
\label{subsec:chiral-inj-proj}

\begin{definition}[Cofree chiral CDG-comodule]
\label{def:cofree-chiral-comod}
\index{cofree comodule!chiral}
For a graded $\mathcal{D}_X$-module $V$, the \emph{cofree graded
chiral $C$-comodule} on $V$ is:
\[
J(V) := C \chirtensor V
\]
with coaction $\Delta \chirtensor \mathrm{id}_V\colon
C \chirtensor V \to C \chirtensor C \chirtensor V$.
The CDG-comodule differential is
$d_{J(V)}(c \otimes v) = d(c) \otimes v$,
and the curvature relation $d_{J(V)}^2 = h \ast (-)$ is inherited
from~$C$.
\end{definition}

\begin{proposition}[Injective and projective resolutions;
\ClaimStatusProvedHere]\label{prop:chiral-inj-proj-resolutions}
\index{injective resolution!chiral CDG-comodule}
\index{projective resolution!chiral CDG-contramodule}
Let $(C, d, h)$ be a conilpotent chiral CDG-coalgebra with
finite-dimensional graded pieces.
\begin{enumerate}[label=\textup{(\alph*)}]
\item For any left chiral CDG-comodule $M$, the cofree chiral
 comodule $J(M) = C \chirtensor M$ is \emph{injective} in the
 abelian category of graded $C^{\#}$-comodules (where $C^{\#}$
 denotes $C$ with forgotten differential and curvature).
\item For any left chiral CDG-contramodule $P$, the free chiral
 contramodule $F(P) = \mathrm{Hom}^{\mathrm{ch}}(C, P)$ is
 \emph{projective} in the abelian category of graded
 $C^{\#}$-contramodules.
\item The homotopy categories admit semiorthogonal decompositions:
 the coderived category $D^{\mathrm{co}}(C\text{-}\mathrm{comod}^{\mathrm{ch}})$
 is equivalent to the homotopy category of CDG-comodules
 with injective underlying graded comodules, and
 $D^{\mathrm{ctr}}(C\text{-}\mathrm{contra}^{\mathrm{ch}})$
 is equivalent to the homotopy category of CDG-contramodules
 with projective underlying graded contramodules.
\end{enumerate}
\end{proposition}

\begin{proof}
\emph{Part (a).}
Let $0 \to L \to M \to N \to 0$ be an exact sequence of graded
$C^{\#}$-comodules. We must show that
$0 \to \mathrm{Hom}_{C^{\#}}(N, J(V)) \to \mathrm{Hom}_{C^{\#}}(M, J(V))
\to \mathrm{Hom}_{C^{\#}}(L, J(V)) \to 0$
is exact. By the tensor-hom adjunction in the chiral category:
$\mathrm{Hom}_{C^{\#}}(M, C \chirtensor V)
\cong \mathrm{Hom}_{\mathcal{D}_X}(M, V)$,
and the functor $\mathrm{Hom}_{\mathcal{D}_X}(-, V)$ is exact
on holonomic $\mathcal{D}_X$-modules (since holonomic modules have
finite length).

\emph{Part (b).}
Dual to (a): for an exact sequence
$0 \to P \to Q \to R \to 0$ of graded $C^{\#}$-contramodules,
$\mathrm{Hom}_{C^{\#}}(\mathrm{Hom}^{\mathrm{ch}}(C, V), -)$
is exact because
$\mathrm{Hom}_{C^{\#}}(\mathrm{Hom}^{\mathrm{ch}}(C, V), P)
\cong \mathrm{Hom}_{\mathcal{D}_X}(V, P)$.

\emph{Part (c).}
For any CDG-comodule $M$, the cobar resolution provides a
functorial injective resolution: form
$J = \bigoplus_{n=0}^{\infty} C^{\chirtensor (n+1)} \chirtensor M$
with the standard bar-type differential (the chiral analogue of
Positselski~\cite[\S4.4, proof of Theorem]{Positselski11}).
The closed morphism $M \to J$ has the property that
$\mathrm{cone}(M \to J)$ is coacyclic (each quotient
$\tau_{\leq n}J / \tau_{\leq n-1}J$ is a cofree CDG-comodule, and
the cone is built from these by exact triangles and countable direct
sums). The projectivity statement for contramodules is dual:
take the bar resolution
$\cdots \to \mathrm{Hom}^{\mathrm{ch}}(C^{\chirtensor 2}, P)
\to \mathrm{Hom}^{\mathrm{ch}}(C, P) \to P$
and observe that the cone is contraacyclic (built from free
contramodules by exact triangles and countable direct products).

The semiorthogonal decomposition follows formally: the localization
functor from the full homotopy category to $D^{\mathrm{co}}$
restricts to an equivalence on the full subcategory of objects with
injective underlying graded comodule (since coacyclic objects with
injective underlying comodule are contractible, by the same argument as
Positselski~\cite[\S4.4]{Positselski11}, using the injective
resolution of the previous paragraph). The dual statement holds for
$D^{\mathrm{ctr}}$.
\end{proof}

\begin{proposition}[Explicit CDG Hom-complex; \ClaimStatusProvedHere]
\label{prop:cdg-hom-complex}
\index{CDG-comodule!Hom-complex|textbf}
Let $(C, d_C, h)$ be a chiral CDG-coalgebra and let $M$, $N$ be
left chiral CDG-comodules over~$C$. The \emph{CDG Hom-complex}
$\underline{\mathrm{Hom}}_C^{\mathrm{ch}}(M, N)$ is the graded
$\mathcal{D}_X$-module of graded $C^{\#}$-comodule morphisms,
with differential
\begin{equation}\label{eq:cdg-hom-differential}
d_{\mathrm{Hom}}(f) \;=\; d_N \circ f
\;-\; (-1)^{|f|}\, f \circ d_M.
\end{equation}
This differential satisfies $d_{\mathrm{Hom}}^2(f) = [h_N, f] - [h_M, f]$
where $h_M$, $h_N$ denote the curvature endomorphisms
$h \ast (-)$ on $M$ and $N$ respectively. In particular:
\begin{enumerate}[label=\textup{(\alph*)}]
\item When $h = 0$ \textup{(}genus~$0$\textup{)},
 $d_{\mathrm{Hom}}^2 = 0$ and the Hom-complex is a genuine
 DG-module.
\item When $h \neq 0$, the Hom-complex is a \emph{CDG-module}
 with curvature $h_{\mathrm{Hom}}(f) = h_N \circ f - f \circ h_M$.
 For endomorphisms ($M = N$),
 $d_{\mathrm{Hom}}^2(f) = [h_M, f]$, so $d_{\mathrm{Hom}}$ is
 a genuine differential on the \emph{center}
 $Z(\mathrm{End}_C(M)) = \{f : h_M \circ f = f \circ h_M\}$.
\end{enumerate}
\end{proposition}

\begin{proof}
Direct computation. For $f \in \underline{\mathrm{Hom}}_C^{\mathrm{ch}}(M, N)^k$:
\begin{align*}
d_{\mathrm{Hom}}^2(f)
&= d_N \circ \bigl(d_N \circ f - (-1)^k f \circ d_M\bigr)
 - (-1)^{k+1}\bigl(d_N \circ f - (-1)^k f \circ d_M\bigr) \circ d_M \\
&= d_N^2 \circ f + (-1)^{k+1}\, d_N \circ f \circ d_M
 - (-1)^{k+1}\, d_N \circ f \circ d_M - f \circ d_M^2 \\
&= d_N^2 \circ f - f \circ d_M^2 \\
&= (h_N \ast -) \circ f - f \circ (h_M \ast -),
\end{align*}
where the last step uses the CDG relation $d^2 = h \ast (-)$ on
$M$ and~$N$. When $M = N$, this is $[h_M, f]$, which vanishes
on central endomorphisms.
\end{proof}

\begin{corollary}[Contractibility of coacyclic injectives; \ClaimStatusProvedHere]
\label{cor:coacyclic-injective-contractible}
\index{coacyclic!contractibility criterion|textbf}
Let $(C, d, h)$ be a conilpotent chiral CDG-coalgebra with
finite-dimensional graded pieces, and let $M$ be a chiral
CDG-comodule with injective underlying graded $C^{\#}$-comodule.
Then $M$ is coacyclic if and only if it admits a
\emph{contracting homotopy}: a degree~$-1$ graded
$C^{\#}$-comodule endomorphism $s \colon M \to M$ satisfying
\begin{equation}\label{eq:contracting-homotopy-cdg}
d_M \circ s + s \circ d_M \;=\; \mathrm{id}_M.
\end{equation}
\end{corollary}

\begin{proof}
If such an $s$ exists, then $\mathrm{id}_M = d_{\mathrm{Hom}}(s)$
in the CDG Hom-complex of
Proposition~\ref{prop:cdg-hom-complex} (with $M = N$), so $M$ is
null-homotopic and \emph{a fortiori} coacyclic.

Conversely, suppose $M$ is coacyclic with injective underlying
graded comodule. By the semiorthogonal decomposition in
Proposition~\ref{prop:chiral-inj-proj-resolutions}(c),
$D^{\mathrm{co}}(C\text{-}\mathrm{comod}^{\mathrm{ch}})$
is equivalent to
$\mathrm{Hot}(C\text{-}\mathrm{comod}^{\mathrm{ch}}_{\mathrm{inj}})$.
A coacyclic object that already has injective underlying comodule
represents the zero object in this homotopy category, hence is
null-homotopic: there exists a degree~$-1$ closed morphism
$s$ with $d_{\mathrm{Hom}}(s) = \mathrm{id}_M$.
Expanding via~\eqref{eq:cdg-hom-differential} gives~\eqref{eq:contracting-homotopy-cdg}.
\end{proof}

\subsection{The correspondence functors}
\label{subsec:chiral-correspondence-functors}

\begin{definition}[Functors \texorpdfstring{$\Phi_C^{\mathrm{ch}}$}{Phi_C^ch} and \texorpdfstring{$\Psi_C^{\mathrm{ch}}$}{Psi_C^ch}; \ClaimStatusDefinitional]
\label{def:chiral-Phi-Psi}
\index{Positselski!functors PhiPsi@$\Phi_C$, $\Psi_C$}
Let $(C, d, h)$ be a chiral CDG-coalgebra.
\begin{itemize}
\item For a left chiral CDG-contramodule $P$, define
 $\Phi_C^{\mathrm{ch}}(P)$ to be the left chiral CDG-comodule
 whose underlying graded $\mathcal{D}_X$-module is the
 chiral contratensor product $C \odot_C^{\mathrm{ch}} P$.
 Concretely, $\Phi_C^{\mathrm{ch}}(P)$ is a quotient of
 $C \chirtensor P$; the comodule structure comes from the
 comultiplication on the left factor~$C$, and the differential is
 induced by the differentials on $C$ and $P$ via the standard
 formula for the tensor product of graded $\mathcal{D}_X$-modules
 with differentials.
\item For a left chiral CDG-comodule $M$, define
 $\Psi_C^{\mathrm{ch}}(M)$ to be the left chiral CDG-contramodule
 whose underlying graded $\mathcal{D}_X$-module is
 $\mathrm{Hom}_C^{\mathrm{ch}}(C, M) = \mathrm{Cohom}_C^{\mathrm{ch}}(C, M)$.
 The contramodule structure comes from the comultiplication
 on $C$, and the differential is induced from $d_C$ and $d_M$.
\end{itemize}
These define DG-functors:
\begin{align*}
\Phi_C^{\mathrm{ch}}\colon \mathrm{DG}(C\text{-}\mathrm{contra}^{\mathrm{ch}})
&\longrightarrow \mathrm{DG}(C\text{-}\mathrm{comod}^{\mathrm{ch}}), \\
\Psi_C^{\mathrm{ch}}\colon \mathrm{DG}(C\text{-}\mathrm{comod}^{\mathrm{ch}})
&\longrightarrow \mathrm{DG}(C\text{-}\mathrm{contra}^{\mathrm{ch}}).
\end{align*}
\end{definition}

\begin{lemma}[Key properties of \texorpdfstring{$\Phi_C^{\mathrm{ch}}$}{Phi_C^ch} and
\texorpdfstring{$\Psi_C^{\mathrm{ch}}$}{Psi_C^ch}; \ClaimStatusProvedHere]
\label{lem:Phi-Psi-properties}
\begin{enumerate}[label=\textup{(\alph*)}]
\item $\Phi_C^{\mathrm{ch}}$ is left adjoint to $\Psi_C^{\mathrm{ch}}$
 on the homotopy categories.
\item $\Phi_C^{\mathrm{ch}}$ maps the free graded contramodule
 $\mathrm{Hom}^{\mathrm{ch}}(C, V)$ to the cofree graded comodule
 $C \chirtensor V$, for any graded $\mathcal{D}_X$-module $V$.
\item $\Psi_C^{\mathrm{ch}}$ maps the cofree graded comodule
 $C \chirtensor V$ to the free graded contramodule
 $\mathrm{Hom}^{\mathrm{ch}}(C, V)$.
\item $\Phi_C^{\mathrm{ch}}$ maps projective graded contramodules
 to injective graded comodules, and $\Psi_C^{\mathrm{ch}}$ maps
 injective graded comodules to projective graded contramodules.
\end{enumerate}
\end{lemma}

\begin{proof}
\emph{Part (a).}
The adjunction $\Phi_C^{\mathrm{ch}} \dashv \Psi_C^{\mathrm{ch}}$
follows from the standard isomorphism
$\mathrm{Hom}_{\mathcal{D}_X}(C \chirtensor V, P)
\cong \mathrm{Hom}_{\mathcal{D}_X}(V, \mathrm{Hom}^{\mathrm{ch}}(C, P))$
(the chiral tensor-hom adjunction), which makes
$\Phi_C^{\mathrm{ch}}$ left adjoint to $\Psi_C^{\mathrm{ch}}$ as
DG-functors (Lemma~\ref{lem:chiral-co-contra-adjunction}).

\emph{Parts (b) and (c).}
The free graded contramodule on $V$ is
$\mathrm{Hom}^{\mathrm{ch}}(C, V) = \mathrm{Hom}_{\mathcal{D}_X}(C, V)$.
Then:
\[
\Phi_C^{\mathrm{ch}}(\mathrm{Hom}^{\mathrm{ch}}(C, V))
= C \odot_C^{\mathrm{ch}} \mathrm{Hom}^{\mathrm{ch}}(C, V).
\]
By the contramodule-comodule correspondence at the graded level
(cf.\ the example in
Positselski~\cite[\S5.1]{Positselski11}), the
contratensor product of $C$ with the free contramodule on $V$ is
canonically isomorphic to the cofree comodule $C \chirtensor V$.
In the chiral setting, this isomorphism follows from the
tensor-hom adjunction for $\mathcal{D}_X$-modules:
\begin{align*}
C \odot_C^{\mathrm{ch}} \mathrm{Hom}^{\mathrm{ch}}(C, V)
&\cong \mathrm{coker}\bigl(
C \chirtensor \mathrm{Hom}^{\mathrm{ch}}(C, \mathrm{Hom}^{\mathrm{ch}}(C, V))
\rightrightarrows C \chirtensor \mathrm{Hom}^{\mathrm{ch}}(C, V)
\bigr) \\
&\cong C \chirtensor V,
\end{align*}
where the last step uses the counit of the adjunction
$\mathrm{ev}\colon \mathrm{Hom}^{\mathrm{ch}}(C, V) \chirtensor C
\to V$ and the coassociativity of $C$. Part~(c) is the dual computation.

\emph{Part (d).}
Combine (b)--(c) with
Proposition~\ref{prop:chiral-inj-proj-resolutions}(a)--(b):
free graded contramodules are projective and cofree graded comodules
are injective, so $\Phi_C^{\mathrm{ch}}$ sends projective objects
to injective objects and vice versa.
\end{proof}

\subsection{The chiral comodule-contramodule correspondence}
\label{subsec:chiral-co-contra-correspondence}

\begin{theorem}[Chiral comodule-contramodule correspondence under
product exactness; \ClaimStatusConditional]
\label{thm:chiral-co-contra-correspondence}
\phantomsection\label{thm:chiral-positselski-5-2}
\index{comodule-contramodule correspondence!chiral|textbf}
\index{Positselski!chiral correspondence|textbf}
Let $(C, d, h)$ be a conilpotent chiral CDG-coalgebra on a smooth
curve $X$ with finite-dimensional graded pieces. Assume also that
the chiral adjunction
$\Phi_C^{\mathrm{ch}}\dashv\Psi_C^{\mathrm{ch}}$ preserves the
second-kind acyclic generators in the required directions:
$\Phi_C^{\mathrm{ch}}$ sends product-generated contraacyclic objects
to coacyclic comodules, and $\Psi_C^{\mathrm{ch}}$ sends
sum-generated coacyclic objects to contraacyclic contramodules. Then
the derived functors
\begin{align*}
\mathbb{L}\Phi_C^{\mathrm{ch}}\colon
D^{\mathrm{ctr}}(C\text{-}\mathrm{contra}^{\mathrm{ch}})
&\longrightarrow
D^{\mathrm{co}}(C\text{-}\mathrm{comod}^{\mathrm{ch}}), \\
\mathbb{R}\Psi_C^{\mathrm{ch}}\colon
D^{\mathrm{co}}(C\text{-}\mathrm{comod}^{\mathrm{ch}})
&\longrightarrow
D^{\mathrm{ctr}}(C\text{-}\mathrm{contra}^{\mathrm{ch}})
\end{align*}
are mutually inverse equivalences of triangulated categories.
\end{theorem}

\begin{proof}
The proof follows Positselski~\cite[\S5.2]{Positselski11}, with the
ordinary tensor product replaced by the chiral tensor product
$\chirtensor$ and exactness properties established via the geometry
of $\mathcal{D}_X$-modules on configuration spaces. We verify
each step.

\emph{Step~1: Derived functors.}
By Proposition~\ref{prop:chiral-inj-proj-resolutions}(c),
the contraderived and coderived categories
\begin{align*}
D^{\mathrm{ctr}}(C\text{-}\mathrm{contra}^{\mathrm{ch}})
&\simeq K_{\mathrm{proj}}(C\text{-}\mathrm{contra}^{\mathrm{ch}}),\\
D^{\mathrm{co}}(C\text{-}\mathrm{comod}^{\mathrm{ch}})
&\simeq K_{\mathrm{inj}}(C\text{-}\mathrm{comod}^{\mathrm{ch}})
\end{align*}
are equivalent to the homotopy categories of
CDG-contra\-modules (resp.\ CDG-comodules) with
projective (resp.\ injective) underlying graded module.
The left derived functor~$\mathbb{L}\Phi_C^{\mathrm{ch}}$
restricts to projective contramodules;
$\mathbb{R}\Psi_C^{\mathrm{ch}}$ restricts to injective comodules.

\emph{Step~2: Mutual inverse on generators.}
By Lemma~\ref{lem:Phi-Psi-properties}(b)--(c), $\Phi_C^{\mathrm{ch}}$
sends the free graded contramodule $\mathrm{Hom}^{\mathrm{ch}}(C, V)$
to the cofree graded comodule $C \chirtensor V$, and
$\Psi_C^{\mathrm{ch}}$ sends $C \chirtensor V$ back to
$\mathrm{Hom}^{\mathrm{ch}}(C, V)$. These are the generating objects
for the respective categories (any CDG-comodule has an injective
resolution by cofree comodules, and any CDG-contramodule has a
projective resolution by free contramodules, by
Proposition~\ref{prop:chiral-inj-proj-resolutions}).

\emph{Step~3: Equivalence on full triangulated subcategories.}
We claim that the restrictions of $\Phi_C^{\mathrm{ch}}$ and
$\Psi_C^{\mathrm{ch}}$ to the full triangulated subcategories
\begin{align*}
\mathrm{Hot}(C\text{-}\mathrm{contra}^{\mathrm{ch}}_{\mathrm{proj}})
&\subset \mathrm{Hot}(C\text{-}\mathrm{contra}^{\mathrm{ch}}), \\
\mathrm{Hot}(C\text{-}\mathrm{comod}^{\mathrm{ch}}_{\mathrm{inj}})
&\subset \mathrm{Hot}(C\text{-}\mathrm{comod}^{\mathrm{ch}})
\end{align*}
are mutually inverse equivalences.

To see this, observe that the unit and counit natural transformations
\[
\eta\colon \mathrm{id} \to \Psi_C^{\mathrm{ch}} \circ \Phi_C^{\mathrm{ch}},
\qquad
\varepsilon\colon \Phi_C^{\mathrm{ch}} \circ \Psi_C^{\mathrm{ch}} \to \mathrm{id}
\]
are isomorphisms on free contramodules and cofree comodules
(Step~2). Since these objects generate the respective
homotopy categories (every object is built from generators by
shifts, cones, and countable (co)products),
and both $\Phi_C^{\mathrm{ch}}$ and $\Psi_C^{\mathrm{ch}}$ preserve
exact triangles and countable (co)products, the natural
transformations $\eta$ and $\varepsilon$ are isomorphisms
on all objects with projective (resp.\ injective) underlying graded
structure.

\emph{Step~4: Descent to exotic derived categories.}
To pass from the full homotopy categories to $D^{\mathrm{co}}$ and
$D^{\mathrm{ctr}}$, we verify that $\Phi_C^{\mathrm{ch}}$ maps
contraacyclic CDG-contramodules to coacyclic CDG-comodules, and
$\Psi_C^{\mathrm{ch}}$ maps coacyclic CDG-comodules to
contraacyclic CDG-contramodules.

For $\Phi_C^{\mathrm{ch}}$, exact triples are preserved by
Lemma~\ref{lem:Phi-Psi-properties}(d) after passing to the derived
functor on projective contramodule resolutions. The remaining
second-kind point is not formal from left-adjointness: contraacyclic
objects are generated using direct products. This is exactly the
product-exactness hypothesis in the statement. Under that hypothesis
$\Phi_C^{\mathrm{ch}}$ sends contraacyclic objects to coacyclic
objects. The dual sum/product assertion for $\Psi_C^{\mathrm{ch}}$
is the second half of the same hypothesis.

\emph{Step~5: Conclusion.}
Combining Steps~3 and~4: $\mathbb{L}\Phi_C^{\mathrm{ch}}$ and
$\mathbb{R}\Psi_C^{\mathrm{ch}}$ are well-defined triangulated
functors between $D^{\mathrm{ctr}}$ and $D^{\mathrm{co}}$, they
are mutually inverse on the generating subcategories (projective
contramodules $\leftrightarrow$ injective comodules), and both
categories are compactly generated by finite-dimensional objects
(using conilpotency and finite-type grading; the argument parallels
Positselski~\cite[\S5.5]{Positselski11}, with the key input that any
graded $C$-comodule is the union of its finite-dimensional
subcomodules, which holds for conilpotent chiral coalgebras because
the conformal weight grading bounds the coaction filtration).
By the uniqueness of extensions of equivalences from compact generators
to the full compactly generated category
(cf.~\cite[\S5.5]{Positselski11}), the equivalence extends to all objects.
\end{proof}

\begin{remark}[Hypotheses]\label{rem:co-contra-hypotheses}
The two raw hypotheses are conilpotency and finite-type grading. They
hold on the finite-type conilpotent bar-coalgebra surface. The
additional second-kind hypothesis is product/sum exactness of the
chiral co/contra adjunction; it is automatic only after it is checked
in the chosen finite or completed ambient. The hypotheses do not
hold for every standard chiral algebra as a raw direct-sum statement:
Virasoro and principal $\mathcal W$-algebras require the
weight-completed / pro-conilpotent replacement of
Theorem~\ref{thm:chiral-positselski-weight-completed}. Finite-type
grading follows from the finite-generation hypothesis on~$\cA$; raw
conilpotency is a separate hypothesis.
\end{remark}

\subsection{Application: the Positselski equivalence on the chiral bar-coalgebra surface}
\label{subsec:positselski-chiral-equivalence}

The chiral comodule-contramodule correspondence applies directly to the
bar coalgebra. Its intrinsic output is a coderived/contraderived
equivalence for $\barB^{\mathrm{ch}}(\cA)$ and its contramodules; only
on additional finite-type or genus-$0$ bar-dual surfaces may one
rewrite the target in terms of ordinary modules over an explicit dual
algebra.

\begin{theorem}[Positselski equivalence for the finite-type chiral bar
coalgebra; \ClaimStatusConditional]
\label{thm:positselski-chiral-finite-type}
\index{Positselski!chiral equivalence|textbf}
\index{comodule-contramodule correspondence!chiral Koszul|textbf}
For a Koszul chiral algebra $\cA$ on a smooth curve $X$ whose
chiral bar coalgebra $C = \bar{B}^{\mathrm{ch}}(\cA)$ lies in the
finite-type conilpotent CDG surface and satisfies the product/sum
exactness hypothesis of
Theorem~\textup{\ref{thm:chiral-co-contra-correspondence}}, there is
an equivalence of
triangulated categories:
\begin{equation}\label{eq:positselski-chiral-finite-type}
D^{\mathrm{co}}(C\text{-}\mathrm{comod}^{\mathrm{conil,\,ch}})
\;\simeq\;
D^{\mathrm{ctr}}(C\text{-}\mathrm{contra}^{\mathrm{ch}})
\end{equation}
between the coderived category of conilpotent chiral
CDG-comodules over the bar coalgebra~$C$ and the contraderived
category of chiral CDG-contramodules over the same
CDG-coalgebra~$C$.

When $C$ has finite-dimensional graded pieces, the right-hand side may
be rewritten as the contraderived category of complete modules over the
graded dual algebra
$C^*=\mathrm{Hom}_{\mathcal{D}_X}(C,\omega_X)$
\textup{(}Remark~\textup{\ref{rem:chiral-contramodules}}\textup{)}.
On the quadratic genus-$0$ bar-dual surface where
Theorem~\textup{\ref{thm:bar-computes-dual}} applies, $C^*$ identifies
with~$\cA^!$, and one recovers the genus-$0$ module Koszul duality of
Theorem~\textup{\ref{thm:e1-module-koszul-duality}}.
At genus~$g \geq 1$, the coderived/contraderived categories are the
correct ambient categories for this bar-coalgebra comparison.
For class-$\mathsf M$ raw direct-sum bar coalgebras this theorem is
not the applicable surface; the completed conditional comparison is
Theorem~\textup{\ref{thm:chiral-positselski-weight-completed}}.
\end{theorem}

\begin{proof}
Apply Theorem~\ref{thm:chiral-co-contra-correspondence} with
$C = \bar{B}^{\mathrm{ch}}(\cA)$.

\emph{Conclusion.}
The hypotheses of Theorem~\ref{thm:chiral-co-contra-correspondence}
are precisely the hypotheses imposed in the statement: $C$ is
conilpotent, finite type, and product/sum exact for the chiral
co/contra adjunction. The equivalence
$D^{\mathrm{co}}(C\text{-}\mathrm{comod}^{\mathrm{ch}})
\simeq D^{\mathrm{ctr}}(C\text{-}\mathrm{contra}^{\mathrm{ch}})$
of Theorem~\ref{thm:chiral-co-contra-correspondence} therefore yields
\eqref{eq:positselski-chiral-finite-type}.

When the finite-type hypothesis of
Remark~\ref{rem:chiral-contramodules} applies, chiral
CDG-contramodules over~$C$ are equivalently complete modules over the
graded dual algebra~$C^*$. The further genus-$0$ identification
$C^* \cong \cA^!$ uses the quadratic-regime bar-dual theorem
\textup{(}Theorem~\textup{\ref{thm:bar-computes-dual}}\textup{)}, so it
is recorded only on that bar-dual surface.
\end{proof}

\begin{theorem}[Flat finite-type reduction on the completed-dual side;
\ClaimStatusConditional]
\label{thm:completed-dual-finite-type-reduction}
\index{derived category!flat finite-type reduction|textbf}
\index{Positselski!derived reduction on completed dual|textbf}
Let $C=\barB^{\mathrm{ch}}(\cA)$ be the chiral bar coalgebra of a
Koszul chiral algebra~$\cA$, and assume we are on a flat finite-type
locus where Positselski's comparison theorem identifies the
contraderived category of chiral CDG-contramodules over~$C$ with the
ordinary bounded derived category of complete modules over the graded
dual algebra~$C^*$. Then there is a triangulated equivalence
\begin{equation}\label{eq:completed-dual-finite-type-reduction}
D^{\mathrm{ctr}}(C\text{-}\mathrm{contra}^{\mathrm{ch}})
\;\simeq\;
D^b(\mathrm{Mod}^{\mathrm{compl}}(C^*)).
\end{equation}
Consequently, on that same locus, Theorem~\textup{\ref{thm:positselski-chiral-finite-type}}
reduces to an ordinary derived comparison between the coderived
bar-coalgebra side and the bounded derived category of complete
$C^*$-modules. On the quadratic genus-$0$ bar-dual surface where
$C^*\cong\cA^!$, this recovers the usual module Koszul duality of
Theorem~\textup{\ref{thm:e1-module-koszul-duality}}.
\end{theorem}

\begin{proof}
\emph{Step~1.}
By Remark~\ref{rem:chiral-contramodules}, under the standing
finite-type hypothesis, chiral CDG-contramodules over~$C$ are
equivalent to complete modules over the graded dual algebra~$C^*$.

\emph{Step~2.}
On the additional flat finite-type locus assumed in the statement,
Positselski's comparison theorem
\cite[\S4.5, Theorem~4.5.1]{Positselski11} identifies the
contraderived category of these contramodules with the ordinary
bounded derived category of complete $C^*$-modules. This gives
\eqref{eq:completed-dual-finite-type-reduction}.

\emph{Step~3.}
Combining \eqref{eq:completed-dual-finite-type-reduction} with
Theorem~\ref{thm:positselski-chiral-finite-type} gives the stated ordinary
derived reduction of the intrinsic coderived/contraderived
equivalence. On the quadratic genus-$0$ bar-dual surface, the
identification $C^*\cong \cA^!$ is supplied by
Theorem~\ref{thm:bar-computes-dual}, recovering
Theorem~\ref{thm:e1-module-koszul-duality}.
\end{proof}

\begin{remark}[Relation to genus-0 module Koszul duality]
\label{rem:genus-0-recovery}
At genus~$0$, the curvature $h = 0$, so CDG-comodules reduce to
DG-comodules, the coderived category reduces to the ordinary derived
category, and on the quadratic bar-dual surface the completed-dual
algebra $C^*=\barB^{\mathrm{ch}}(\cA)^*$ identifies with~$\cA^!$.
Theorem~\ref{thm:completed-dual-finite-type-reduction} is precisely the
flat finite-type reduction that then recovers the ordinary
genus-$0$ module Koszul duality of
Theorem~\ref{thm:e1-module-koszul-duality}.
The coderived/contraderived framework generalizes to genus~$g \geq 1$, where the
curvature $m_0^{(g)} \neq 0$ prevents replacing the exotic
categories by ordinary derived categories without extra flatness input.
\end{remark}

\begin{remark}[The role of the configuration space geometry]
\label{rem:config-space-role}
The proof of the chiral comodule-contramodule correspondence
(Theorem~\ref{thm:chiral-co-contra-correspondence}) uses the
geometry of configuration spaces in three places:
\begin{enumerate}
\item \emph{Exactness of the chiral tensor product.}
 The tensor-hom adjunction for holonomic $\mathcal{D}_X$-modules
 on $\overline{C}_n(X)$ ensures that cofree comodules are injective
 and free contramodules are projective
 (Proposition~\ref{prop:chiral-inj-proj-resolutions}).
\item \emph{Conilpotency from conformal weight.}
 The conformal weight grading on $\bar{B}^{\mathrm{ch}}(\cA)$
 (inherited from $\cA$) ensures conilpotency
 (Theorem~\ref{thm:coalgebra-via-NAP}(4)), which is needed for
 the compact generation argument in Step~5 of the proof.
\item \emph{Finite-dimensionality from holonomicity.}
 The holonomicity of the bar complex on each configuration space
 stratum (Lemma~\ref{lem:bar-holonomicity}) ensures
 finite-dimensional graded pieces, which is
 needed for the duality between contramodules and complete modules
 (Remark~\ref{rem:chiral-contramodules}).
\end{enumerate}
These geometric inputs are specific to the chiral setting and have
no analogue in Positselski's ground-field framework.
\end{remark}

%================================================================
% SECTION: BAR-COBAR INVERSION - COMPLETE QUASI-ISOMORPHISM
%================================================================

\section{Bar--cobar reconstruction and quadratic recognition}
\label{sec:bar-cobar-inversion-quasi-iso}

The bar construction produces the universal reconstruction coalgebra
\(B_X(A)\).  A quadratic presentation produces the comparison
coalgebra \(A^{\mathrm i}\).  The first object belongs to Theorem~A;
the morphism \(A^{\mathrm i}\to B_X(A)\) belongs to Theorem~B.

\subsection{Universal reconstruction in the Ran ambient}
\label{sec:theorem-A-platonic}

\begin{theorem}[Associative Ran reconstruction and factorization
transport; \ClaimStatusConditional]
\label{thm:bar-cobar-platonic}
\textup{[Type signature: Open quadrant; associative Ran and
factorization presentations; Beilinson levels
\(1\leftrightarrow2\); pro-nilpotent Francis--Gaitsgory Ran ambient;
hypothesis packages \(H_{\mathrm{fact}}\) and
\(H_{\mathrm{conv}}\).]}
Let
\(\mathcal C_X=(D(\operatorname{Ran}X),\otimes^{\mathrm{ch}})\)
for a separated finite-type scheme \(X\) over a characteristic-zero
field.
\begin{enumerate}[label=\textup{(\roman*)}]
\item \emph{Ambient Ran reconstruction;
\(\ClaimStatusProvedElsewhere\).}
Enhanced bar and cobar for the reduced associative operad are mutually
inverse on the pro-nilpotent Ran ambient.  For an augmented algebra
\(A\) and a conilpotent coalgebra \(C\), the induced maps are
\[
 \varepsilon_A\colon\Omega^{\mathrm{enh}}B^{\mathrm{enh}}(A)
 \xrightarrow{\sim}A,
 \qquad
 \eta_C\colon C\xrightarrow{\sim}
 B^{\mathrm{enh}}\Omega^{\mathrm{enh}}(C).
\]
\item \emph{Factorization restriction;
\(\ClaimStatusConditional\).}
The package \(H_{\mathrm{fact}}\) makes enhanced bar and cobar preserve
the factorization subcategories and their collision extensions.
\item \emph{Chain-model realization;
\(\ClaimStatusConditional\).}
The package \(H_{\mathrm{conv}}\) identifies the selected
Fulton--MacPherson chain models with the ordinary associative
bar--cobar complexes on each stratum and transports the unit, counit,
and their homotopies through the Ran gluing and completed inverse
limits.
\end{enumerate}
\end{theorem}

\begin{proof}
Francis--Gaitsgory
\cite[Proposition~4.1.2]{Francis2012} prove enhanced operadic
bar--cobar equivalence in a pro-nilpotent stable symmetric-monoidal
\(\infty\)-category; the proof of
\cite[Theorem~5.1.1]{Francis2012} supplies pro-nilpotence for the
chiral Ran product.  On ordinary chain models the unit and counit are
the universal bar--cobar resolutions of
Loday--Vallette \cite[Corollary~2.3.4]{LV12}.  The remaining two
clauses are precisely the structure and compatibility maps listed in
\(H_{\mathrm{fact}}\) and \(H_{\mathrm{conv}}\).
\end{proof}

\begin{proposition}[Unit and counit of the twisting adjunction;
\ClaimStatusProvedHere]
\label{prop:unit-counit-normalization-bci}
The twisting-morphism bijection
\[
\operatorname{Hom}_{\mathrm{Alg}}(\Omega_XC,A)
\cong\operatorname{Tw}(C,A)
\cong\operatorname{Hom}_{\mathrm{CoAlg}}(C,B_XA)
\]
assigns to the identity of \(\Omega_XC\) the unit
\(\eta_C\colon C\to B_X\Omega_XC\), and to the identity of \(B_XA\)
the counit
\(\varepsilon_A\colon\Omega_XB_XA\to A\).  Their composites are
\[
\varepsilon_{\Omega_XC}\circ\Omega_X(\eta_C)
=\operatorname{id}_{\Omega_XC},
\qquad
B_X(\varepsilon_A)\circ\eta_{B_XA}
=\operatorname{id}_{B_XA}.
\]
\end{proposition}

\begin{proof}
The two identities are the triangular identities of the adjunction
\(\Omega_X\dashv B_X\), written under the universal
twisting-morphism bijection.
\end{proof}

\begin{remark}[Ordered and symmetric presentations]
\label{rem:bar-cobar-platonic-foundational-scope}
The ordered object
\(B_X^{\mathrm{ord}}(A)=T^c(s^{-1}\bar A)\) carries the
deconcatenation coproduct and the adjacent-collision differential.
The symmetric object \(B_X^\Sigma(A)\) is obtained by homotopy descent
along the ordered configuration torsors.  The package
\(H_{\mathrm{fact}}^R\) supplies the equivariant local system,
nearby-cycle extension, and comparison maps that transport
reconstruction from the ordered presentation to the symmetric one.
\end{remark}

\begin{remark}[Finite and completed chain presentations]
\label{rem:bar-cobar-platonic-class-M-ambient}
A finite-window presentation uses direct sums and the ordinary cobar
construction.  A completed presentation uses a strict inverse system
of finite windows, continuous cobar, and the Mittag--Leffler
comparison.  Theorem~\ref{thm:bar-convergence} identifies the
conditions under which the direct-sum and completed bar objects are
quasi-isomorphic.  The following package records the completed
reconstruction statement independently of quadratic Koszulness.
\end{remark}

\begin{definition}[Finite-window reconstruction package;
\ClaimStatusDefinitional]
\label{def:finite-window-essential-image-bci}
A complete augmented chiral algebra \(A\) carries the
\emph{finite-window reconstruction package} when it is equipped with:
\begin{enumerate}[label=\textup{(R\arabic*)}]
\item quotient algebras \(A_{\leq N}\) with strict surjective
transition maps and \(A\simeq\varprojlim_NA_{\leq N}\);
\item finite conilpotent bar coalgebras \(B_X(A_{\leq N})\) and
universal counits
\(\varepsilon_N\colon\Omega_XB_X(A_{\leq N})\to A_{\leq N}\);
\item compatible realizations of the ordinary universal
bar--cobar contraction in every finite window;
\item Mittag--Leffler cohomology towers for the counits and their
mapping cones;
\item continuous comparison maps
\[
B_X^{\mathrm{cont}}(A)\xrightarrow{\sim}
\varprojlim_NB_X(A_{\leq N}),
\qquad
\Omega_X^{\mathrm{cont}}B_X^{\mathrm{cont}}(A)
\xrightarrow{\sim}
\varprojlim_N\Omega_XB_X(A_{\leq N}).
\]
\end{enumerate}
\end{definition}

\begin{proposition}[Completed universal reconstruction;
\ClaimStatusConditional]
\label{prop:finite-window-ml-inversion-bci}
If \(A\) carries the finite-window reconstruction package, then the
continuous counit
\[
\widehat\varepsilon_A\colon
\Omega_X^{\mathrm{cont}}B_X^{\mathrm{cont}}(A)\longrightarrow A
\]
is a quasi-isomorphism of complete augmented factorization algebras.
\end{proposition}

\begin{proof}
Each finite-stage counit is a quasi-isomorphism by the ordinary
universal bar--cobar resolution
\cite[Corollary~2.3.4]{LV12}.  The cone tower is
Mittag--Leffler, so its derived inverse limit has zero cohomology.
The continuous comparison maps identify this inverse limit with the
cone of \(\widehat\varepsilon_A\).
\end{proof}

\subsection{The five typed objects}

\begin{proposition}[Five objects and their canonical morphisms;
\ClaimStatusProvedHere]
\label{prop:bar-cobar-object-firewall-bci}
Let \(A=T(V)/(R)\) be a complete augmented quadratically presented
chiral algebra.
\begin{enumerate}[label=\textup{(\roman*)}]
\item \(A\) is an augmented factorization algebra.
\item \(B_X(A)=T^c(s^{-1}\bar A)\) is its complete conilpotent
bar coalgebra.
\item \(A^{\mathrm i}=C(s^{-1}V,s^{-2}R)\) is the quadratic coalgebra of the
chosen presentation, with canonical comparison
\(q_A\colon A^{\mathrm i}\to B_X(A)\).
\item Under \(H_{\mathbb D}^{\mathrm{bar}}(A)\),
\[
A^!_\infty:=\mathbb D_{\operatorname{Ran}}B_X(A)
\]
is the post-Verdier complete algebra.  Under local finite duality,
\(A^!:=\mathbb D_{\operatorname{Ran}}A^{\mathrm i}\) is the strict
quadratic dual algebra.
\item
\(Z^{\mathrm{der}}_{\mathrm{ch}}(A)
=C^\bullet_{\mathrm{ch}}(A,A)\)
is the derived chiral centre.
\end{enumerate}
The canonical morphisms are the reconstruction counit
\(\varepsilon_A\colon\Omega_XB_X(A)\to A\), the quadratic comparison
\(q_A\), its Verdier transpose
\[
\mathbb D_{\operatorname{Ran}}(q_A)\colon
A^!_\infty\longrightarrow A^!,
\]
and the defining Hochschild map to the derived centre.
\end{proposition}

\begin{proof}
Each object is obtained by the displayed functor in its stated
ambient.  The four morphisms follow from the bar--cobar adjunction,
the inclusion of quadratic corelations, contravariance of Verdier
duality, and the Hochschild cochain construction.
\end{proof}

\begin{theorem}[Reconstruction--Verdier comparison;
\ClaimStatusProvedHere]
\label{thm:omega-bar-not-verdier-koszul-dual}
Assume the reconstruction package makes
\(\varepsilon_A\colon R_X(A):=\Omega_XB_X(A)\to A\) an equivalence,
and assume \(H_{\mathbb D}^{\mathrm{bar}}(A)\) defines
\(K_X(A):=\mathbb D_{\operatorname{Ran}}B_X(A)\).
Precomposition with \(\varepsilon_A\) gives an equivalence
\[
\operatorname{Map}(A,K_X(A))
\xrightarrow{\sim}
\operatorname{Map}(R_X(A),K_X(A)).
\]
Accordingly, a self-duality equivalence
\(\sigma_A\colon A\xrightarrow{\sim}K_X(A)\) determines the comparison
\(\sigma_A\circ\varepsilon_A\colon R_X(A)\xrightarrow{\sim}K_X(A)\),
and every equivalence \(R_X(A)\simeq K_X(A)\) arises uniquely in this
way up to the contractible mapping-space choice.
\end{theorem}

\begin{proof}
Precomposition with an equivalence induces an equivalence of mapping
spaces.  Apply this fact to \(\varepsilon_A\).
\end{proof}

\begin{corollary}[The homotopy Verdier--Koszul dual algebra;
\ClaimStatusConditional]
\label{cor:A-shriek-constructed}
Under \(H_{\mathbb D}^{\mathrm{bar}}(A)\), the formula
\[
A^!_\infty:=\mathbb D_{\operatorname{Ran}}B_X(A)
\]
defines the homotopy Verdier--Koszul dual algebra in the complete
post-Verdier factorization ambient.  Its multiplication is the Verdier
transpose of the bar coproduct.
\end{corollary}

\begin{corollary}[Quadratic comparison after Verdier duality;
\ClaimStatusConditional]
\label{cor:finite-type-verdier-koszul-dual-formality}
Assume \(H_{\mathrm{CL}}(A,A^{\mathrm i},\tau_{\mathrm i})\) and
\(H_{\mathbb D}^{\mathrm{bar}}(A)\), and assume the
quadratic pieces are locally finite and Verdier dualizable.  Then
\[
q_A\colon A^{\mathrm i}\xrightarrow{\sim}B_X(A)
\]
is a quasi-isomorphism of complete chiral coalgebras and
\[
\mathbb D_{\operatorname{Ran}}(q_A)\colon
A^!_\infty\xrightarrow{\sim}A^!
\]
is a quasi-isomorphism of complete chiral algebras.
\end{corollary}

\begin{proof}
The first map is the quadratic recognition theorem transported by
\(H_{\mathrm{CL}}\).  Verdier duality is exact on the supplied
dualizable finite windows and continuous on their strict inverse
limit, giving the second map.
\end{proof}

\begin{proposition}[Formality obstruction for the homotopy dual;
\ClaimStatusConditional]
\label{prop:verdier-koszul-formality-criterion-bci}
Assume \(A^!_\infty\) admits a finite-window strong deformation
retract onto
\(M=H^\bullet(A^!_\infty)\), compatible with the strict
Mittag--Leffler completion.  Homotopy transfer gives a minimal
\(A_\infty\)-structure
\(\{m_r^{\mathrm{tr}}\}_{r\geq2}\) on \(M\).  The object
\(A^!_\infty\) is formal precisely when the recursive Hochschild
classes
\[
\mathfrak o_r(A)=[m_r^{\mathrm{tr}}]
\in HH_{\mathrm{ch}}^{\,2-r}(M,M),
\qquad r\geq3,
\]
vanish after the lower arities have been removed by
\(A_\infty\)-gauge transformations.
\end{proposition}

\begin{proof}
The minimal-model obstruction theory identifies the first surviving
higher multiplication with the obstruction to extending a
strictification through that arity.  Induction on \(r\), followed by
the strict Mittag--Leffler limit, gives the criterion.
\end{proof}

\begin{proposition}[Standard-family quadratic comparison;
\ClaimStatusConditional]
\label{prop:standard-family-firewall-table-bci}
For a parameterized family \(A_U=T(V_U)/(R_U)\), a
\emph{Theorem~B comparison datum} consists of:
\begin{enumerate}[label=\textup{(\alph*)}]
\item the quadratic coalgebra
\(A_U^{\mathrm i}=C(s^{-1}V_U,s^{-2}R_U)\) and the map
\(q_{A_U}\colon A_U^{\mathrm i}\to B_X(A_U)\);
\item a proof of \(H_{\mathrm{CL}}\) on the stated parameter locus;
\item finite-window or completed convergence for the two sides;
\item the Verdier package \(H_{\mathbb D}^{\mathrm{bar}}(A_U)\) whenever a dual
algebra is asserted;
\item a separate Hochschild comparison whenever a centre statement
is asserted.
\end{enumerate}
A family carrying this datum inherits the five typed objects and
morphisms of
Proposition~\ref{prop:bar-cobar-object-firewall-bci}.  Its Theorem~B
status is the status of the comparison \(q_{A_U}\).
\end{proposition}

\paragraph{Boundary comparison problems.}
The following conjectures isolate four extensions of the proved
ambient and quadratic-comparison lanes.

\begin{conjecture}[$\Pi_1$: stratified factorization transport;
\ClaimStatusConjectured]
\label{conj:pi1-francis-gaitsgory}
\textup{[Type signature: Open quadrant; Whitney-stratified
factorization presentation; Beilinson levels
\(1\leftrightarrow2\); hypothesis package: stratified
\(H_{\mathrm{fact}}\).]}
For a Whitney-regular stratification of \(X\), the enhanced
associative Ran equivalence restricts to constructible
factorization algebras whose collision and nearby-cycle maps preserve
the chosen strata.
\end{conjecture}

\begin{conjecture}[$\Pi_2$: higher-\(E_n\) geometric transport;
\ClaimStatusConjectured]
\label{conj:pi2-en-bar-higher-dimension}
\textup{[Type signature: geometric quadrant; locally constant
\(E_n\)-factorization presentation; bar--cobar level; hypothesis
package: framed descent and completed \(E_n\)-Koszul duality.]}
The local \(E_n\) bar--cobar equivalence on framed disks extends,
under framed factorization descent and completion, to the selected
global factorization models on an \(n\)-manifold.
\end{conjecture}

\begin{conjecture}[$\Pi_3$: Lagrangian recognition of the quadratic
comparison; \ClaimStatusConjectured]
\label{conj:pi3-lagrangian-koszul-converse}
\textup{[Type signature: Open quadrant; derived deformation
presentation; Beilinson levels \(1\leftrightarrow2\); hypothesis
package: perfect shifted-symplectic pairing and a derived-intersection
model for \(\operatorname{Cone}(q_A)\).]}
When the derived intersection of the two quadratic deformation
branches represents the comparison cone of
\(q_A\colon A^{\mathrm i}\to B_X(A)\), transversality of the branches
is equivalent to quadratic Koszul recognition.
\end{conjecture}

\begin{conjecture}[$\Pi_4$: unbounded completed quadratic
recognition; \ClaimStatusConjectured]
\label{conj:pi4-unbounded-rank}
\textup{[Type signature: Open quadrant; strict inverse-limit
presentation; Beilinson levels \(1\leftrightarrow2\); hypothesis
package: uniform finite-window \(H_{\mathrm{CL}}\) and
Mittag--Leffler convergence.]}
For a strict tower of quadratically presented finite-window quotients,
windowwise quasi-isomorphisms
\[
A_{\leq N}^{\mathrm i}\xrightarrow{\sim}B_X(A_{\leq N})
\]
assemble to a quasi-isomorphism between the completed quadratic
coalgebra and the completed bar coalgebra.
\end{conjecture}

\subsection{Theorem~B: quadratic Koszul recognition}

\begin{theorem}[Quadratic Koszul recognition;
\ClaimStatusConditional]
\label{thm:bar-cobar-inversion-qi}
\textup{[Type signature: Open quadrant; completed chiral
Ran/factorization presentation; Beilinson levels
\(1\leftrightarrow2\); hypothesis package: a connected positive
quadratic presentation and
\(H_{\mathrm{CL}}(A,A^{\mathrm i},\tau_{\mathrm i})\).]}

Let \(A=T(V)/(R)\) be a complete augmented quadratically presented
chiral algebra.  Let
\[
A^{\mathrm i}=C(s^{-1}V,s^{-2}R)\subset T^c(s^{-1}V)
\]
be its quadratic chiral coalgebra, let
\(\tau_{\mathrm i}\colon A^{\mathrm i}\to A\) be the canonical
quadratic twisting morphism, and let
\[
q_A\colon A^{\mathrm i}\longrightarrow B_X(A)
\]
be the corresponding coalgebra map.  Assume
\(H_{\mathrm{CL}}(A,A^{\mathrm i},\tau_{\mathrm i})\).
Then the following conditions are equivalent:
\begin{enumerate}[label=\textup{(\roman*)}]
\item the associated-graded quadratic twisting morphism is Koszul in
the sense of Loday--Vallette;
\item \(q_A\) is a quasi-isomorphism of complete chiral coalgebras;
\item the adjoint morphism
\[
\Omega_X(A^{\mathrm i})\longrightarrow A
\]
is a quasi-isomorphism of complete augmented chiral algebras;
\item the two twisted tensor products
\[
K^L_{\tau_{\mathrm i}}(A,A^{\mathrm i}),
\qquad
K^R_{\tau_{\mathrm i}}(A^{\mathrm i},A)
\]
are acyclic;
\item the PBW/bar spectral sequence is supported on the quadratic
Koszul diagonal and converges strongly to the completed comparison.
\end{enumerate}
Any one of these equivalent conditions defines the quadratic Koszul
locus of the chosen presentation.
\end{theorem}

\begin{proof}
On every associated-graded finite window, the equivalence of
conditions~\textup{(i)}--\textup{(iv)} is the fundamental theorem of
twisting morphisms
\cite[Theorem~2.3.2]{LV12}; for the quadratic coalgebra
\(C(s^{-1}V,s^{-2}R)\), the comparison with the bar construction and the
quadratic cobar criterion are
\cite[Theorem~3.4.6]{LV12}.  The filtration in
\(H_{\mathrm{CL}}\) identifies condition~\textup{(v)} with this
associated-graded diagonal criterion.  Split finite windows, strong
convergence, and the Mittag--Leffler clause transport all five
conditions to the completed Ran/factorization objects.
\end{proof}

\begin{remark}[Three reconstruction problems]
Theorem~\ref{thm:bar-cobar-inversion-qi} concerns the map
\(A^{\mathrm i}\to B_X(A)\).  The universal counit
\(\Omega_XB_X(A)\to A\) belongs to
Theorem~\ref{thm:bar-cobar-platonic}.  For a fixed bar coalgebra
\(C=B_X(A)\), the equivalence between \(C\)-comodules and
\(C\)-contramodules belongs to
Theorem~\ref{thm:off-koszul-ran-inversion}.  A module comparison
induced by a general twisting morphism is governed by
\(\mathsf{Tw}^{\mathrm{ch}}_{\mathrm{acyc}}(C,A,\tau)\).
\end{remark}

\begin{definition}[Ambient reconstruction test;
\ClaimStatusDefinitional]
\label{def:counit-detected-weak-equivalence-bci}
Let \(A\) be a complete augmented chiral algebra equipped with either
a finite-window bar model or a strict completed bar tower.  The
\emph{ambient reconstruction test} is the assertion that the
continuous universal counit
\[
\varepsilon_A\colon
\Omega_X^{\mathrm{cont}}B_X^{\mathrm{cont}}(A)\longrightarrow A
\]
is a quasi-isomorphism in that chosen algebraic ambient.  We denote
this assertion by \(\varepsilon_A\in\mathsf W_{\mathrm{rec}}\).
Theorem~\ref{thm:bar-cobar-platonic} supplies the enhanced Ran
version; Proposition~\ref{prop:finite-window-ml-inversion-bci}
supplies its completed finite-window realization.
\end{definition}

\begin{proposition}[Quadratic Koszul criteria;
\ClaimStatusConditional]
\label{prop:koszul-locus-criteria-package-bci}
Let \(A=T(V)/(R)\) be connected, positively graded, complete, and
quadratically presented.  Put
\(A^{\mathrm i}=C(s^{-1}V,s^{-2}R)\), and assume
\(H_{\mathrm{CL}}(A,A^{\mathrm i},\tau_{\mathrm i})\).
Then the following conditions are equivalent:
\begin{enumerate}[label=\textup{(\roman*)}]
\item \(q_A\colon A^{\mathrm i}\to B_X(A)\) is a
quasi-isomorphism;
\item \(\Omega_X(A^{\mathrm i})\to A\) is a quasi-isomorphism;
\item the two twisted tensor products for
\(\tau_{\mathrm i}\) are acyclic;
\item the associated PBW/bar spectral sequence is supported on the
Koszul diagonal and converges strongly;
\item after the standard bar/internal reindexing,
\(\operatorname{Ext}^{p,q}_A(\mathbf1,\mathbf1)\) is supported on
\(p=q\).
\end{enumerate}
When an additional Fulton--MacPherson boundary comparison identifies
the filtered bar complex with its collision-stratum realization,
these conditions transport to the corresponding boundary
concentration statement.
\end{proposition}

\begin{proof}
Conditions~\textup{(i)}--\textup{(iii)} are
Loday--Vallette's fundamental twisting theorem and quadratic
recognition theorem
\cite[Theorems~2.3.2 and~3.4.6]{LV12}.
For a connected quadratic presentation, the associated graded of the
comparison \(q_A\) is the classical quadratic Koszul complex.
Therefore its acyclicity is equivalent to diagonal support of the
PBW/bar spectral sequence and to diagonal support of the bigraded
Ext algebra.  The split finite-window and strong-convergence clauses
of \(H_{\mathrm{CL}}\) lift the equivalence to the completed chiral
objects.  The final sentence is transport along the supplied
boundary comparison.
\end{proof}

\begin{example}[Admissible
\(\widehat{\mathfrak{sl}}_2\) as a quadratic-comparison boundary;
\ClaimStatusConjectured]
\label{ex:admissible-sl2-failure}
Let \(L_{-1/2}(\mathfrak{sl}_2)\) be the simple admissible quotient of
the universal affine algebra.  Its null-vector ideal changes the
quadratic presentation and hence changes
\[
q_A\colon A^{\mathrm i}\longrightarrow B_X(A).
\]
A proof of quadratic Koszulness at this point requires an
explicit finite-window calculation of the first off-diagonal
cohomology of \(\operatorname{Cone}(q_A)\), together with convergence
of the resulting PBW/bar filtration.  Conjecture~\ref{conj:admissible-2-koszul}
records the proposed outcome.  The universal counit
\(\Omega_XB_X(A)\to A\) remains the reconstruction map of
Theorem~\ref{thm:bar-cobar-platonic}.
\end{example}

\begin{conjecture}[Admissible two-diagonal bar support;
\ClaimStatusConjectured]
\label{conj:admissible-2-koszul}%
\index{admissible level!two-diagonal bar support}%
For \(A=L_{-1/2}(\mathfrak{sl}_2)\), the completed comparison cone
\[
\operatorname{Cone}\!\left(
q_A\colon A^{\mathrm i}\longrightarrow B_X(A)
\right)
\]
has finite-window cohomology supported on the first off-diagonal
adjacent to the quadratic Koszul diagonal.  Equivalently, after the
standard reindexing, the bar cohomology occupies two adjacent
diagonals.  A proof requires the explicit null-vector differential at
conformal weight~\(4\), the subsequent finite-window differentials,
and strong convergence of the completed PBW/bar spectral sequence.
\end{conjecture}

\begin{remark}[Formality and quadratic recognition]
\label{rem:bar-cobar-inversion-alt-formality}
The comparison \(q_A\colon A^{\mathrm i}\to B_X(A)\) may be studied
through the minimal \(A_\infty\)-model of the bar coalgebra.  A
finite-window strong deformation retract transfers the bar structure
to its cohomology.  On the quadratic Koszul locus, diagonal
concentration identifies that cohomology with \(A^{\mathrm i}\), and
the transferred higher operations vanish up to gauge.  This gives the
factorization
\[
A^{\mathrm i}\xrightarrow{\sim}B_X(A),
\qquad
\Omega_X(A^{\mathrm i})\xrightarrow{\sim}A.
\]
The obstruction classes are the Hochschild classes of the transferred
operations, as recorded in
Proposition~\ref{prop:verdier-koszul-formality-criterion-bci}.
\end{remark}

\begin{remark}[Quadratic Koszul locus and its first obstruction]
\label{rem:koszul-locus-off-locus-obstruction}
\index{Koszul locus!quadratic comparison}
For a connected quadratic presentation, the Koszul locus is the
vanishing locus of
\[
\operatorname{Cone}\!\left(
q_A\colon A^{\mathrm i}\longrightarrow B_X(A)
\right).
\]
The first obstruction is the first nonzero cohomology class of this
cone away from the Koszul diagonal.  In the PBW/bar spectral sequence
it appears as the first surviving off-diagonal class; in the minimal
model it appears as the first higher operation obstructing
strictification to the quadratic coalgebra.  Higher-genus curvature
belongs to the modular extension of the bar construction and carries
its own hypothesis package.  The universal reconstruction counit
remains governed by Theorem~\ref{thm:bar-cobar-platonic}.
\end{remark}

\begin{conjecture}[Derived quadratic Koszul wall;
\ClaimStatusConjectured]
\label{conj:koszul-wall-associated-variety}
\index{Koszul locus!derived quadratic wall}
Let \(\mathrm{QuadChAlg}^{\mathrm{aug}}\) denote a derived moduli
problem of complete augmented quadratically presented chiral
algebras for which \(A^{\mathrm i}\), \(B_X(A)\), and \(q_A\) form a
perfect family in every finite window.  Then the derived vanishing
locus
\[
\mathrm{Kos}^{\mathrm{quad}}
:=
\left\{A\;\middle|\;
q_A\colon A^{\mathrm i}\xrightarrow{\sim}B_X(A)\right\}
\]
is represented locally by the derived zero locus of the perfect
comparison cone.  Its natural strata are the Fitting strata of the
first off-diagonal cohomology sheaf of
\(\operatorname{Cone}(q_A)\).  On a representation-theoretic family,
a comparison with associated varieties is supplied by a separate
microlocal theorem identifying these Fitting strata with singular
support.
\end{conjecture}

\begin{remark}[Theorem~B as a specialization of the twisting theorem]
\label{rem:inversion-vs-fundamental}
Theorem~\ref{thm:fundamental-twisting-morphisms} accepts a general
twisting datum \((A,C,\tau)\).  Theorem~B is its quadratic
specialization
\[
(C,\tau)=(A^{\mathrm i},\tau_{\mathrm i}).
\]
Under \(H_{\mathrm{CL}}(A,A^{\mathrm i},\tau_{\mathrm i})\), the unit
is \(q_A\colon A^{\mathrm i}\to B_X(A)\), the counit is
\(\Omega_X(A^{\mathrm i})\to A\), and the twisted tensor products are
the two quadratic Koszul complexes.
\end{remark}

\begin{remark}[Twisted tensor products as the acyclicity criterion]
\label{rem:bar-cobar-twisted-tensor-anchor}
For the quadratic twisting morphism \(\tau_{\mathrm i}\), the
fundamental theorem of twisting morphisms
\cite[Theorem~2.3.2]{LV12} identifies the acyclicity of
\(K^L_{\tau_{\mathrm i}}\) and \(K^R_{\tau_{\mathrm i}}\) with the two
comparison quasi-isomorphisms of Theorem~B.  The package
\(H_{\mathrm{CL}}\) transports this ordinary statement through the
chiral strata, factorization gluing, and completed inverse limit.
\end{remark}

\begin{remark}[Modular extension of quadratic recognition]
\label{rem:higher-genus-mk3-conditionality}
Theorem~B is a level-\(1\leftrightarrow2\) genus-zero recognition
theorem.  A genus-\(g\) analogue requires a modular quadratic
coalgebra \(A_g^{\mathrm i}\), a twisting morphism
\(\tau_{\mathrm i,g}\), a square-zero total graph differential, and a
genuswise comparison package
\(H_{\mathrm{CL}}^{(g)}(A,A_g^{\mathrm i},\tau_{\mathrm i,g})\).
When these data are supplied, the same associated-graded argument
applies.  The modular construction and its convergence are independent
proof obligations.
\end{remark}

\begin{remark}[Homotopy structures carried by the comparison]
\label{rem:bar-cobar-inversion-linfty}
An operadic enhancement of the bar and cobar functors carries
\(q_A\) and \(\Omega_X(A^{\mathrm i})\to A\) to morphisms in the
corresponding homotopy-algebra categories.  Under the convolution
comparison of
Theorem~\ref{thm:operadic-homotopy-convolution}, a Theorem~B
quasi-isomorphism therefore transports the induced
\(L_\infty\)-deformation structures.  This is the homotopy-level
realization of the quadratic comparison.
\end{remark}

\begin{remark}[Chain and algebraic quasi-isomorphisms]
\label{rem:qi-vs-homology-iso}
A morphism of dg or \(A_\infty\) chiral algebras is a
quasi-isomorphism when its linear component is a quasi-isomorphism of
complexes.  The algebraic structure maps are part of the morphism.
An isomorphism between the cohomology vector spaces alone supplies
only the underlying graded comparison.  Homotopy transfer equips
cohomology with a minimal \(A_\infty\)-structure whose binary
operation is the induced product and whose first higher obstruction
is represented by \(m_3\).  The comparison
\(\Omega_X(A^{\mathrm i})\to A\) in Theorem~B is a morphism of
augmented chiral algebras and hence carries this coherent structure.
\end{remark}

\begin{conjecture}[Morita transport from a geometric bar equivalence;
\ClaimStatusConjectured]
\label{conj:bar-morita-koszul-conductor}
\index{derived Morita equivalence!geometric bar transport}
Let \(A\) and \(A'\) be complete augmented factorization algebras in
the factorization-reconstruction domain of
Theorem~\ref{thm:bar-cobar-platonic}.  Suppose an equivalence
\[
B_X(A)\xrightarrow{\sim}B_X(A')
\]
is compatible with factorization gluing, augmentations, and chosen
compact generators.  Then continuous cobar transports this
equivalence to an equivalence \(A\simeq A'\), and the induced
equivalence of module categories preserves the traces used to define
their scalar modular characteristics.  The compact-generator and
trace compatibility constitute the geometric Morita package of the
conjecture.
\end{conjecture}

\subsection{Proof strategy and filtration}

\begin{definition}[Bar-cobar filtration]\label{def:bar-cobar-filtration}
Define a decreasing filtration on $\Omega(\bar{B}(\mathcal{A}))$ by:
\[F^p\Omega(\bar{B}(\mathcal{A})) = \bigoplus_{n \geq p} \Omega^n(\bar{B}^n(\mathcal{A}))\]

This is the filtration by \emph{bar degree} (= cobar degree).

\emph{Geometric meaning.} $F^p$ consists of elements involving at least $p$ points
in configuration space. As $p \to \infty$, we are considering increasingly complicated
configurations.

\emph{Properties.}
\begin{enumerate}
\item $F^0 \supseteq F^1 \supseteq F^2 \supseteq \cdots$
\item $\bigcap_{p=0}^\infty F^p = 0$ (completeness)
\item The differential respects filtration: $d(F^p) \subseteq F^p$
\item The natural map factors through the filtration
\end{enumerate}
\end{definition}

\begin{lemma}[Associated graded; \ClaimStatusProvedHere]\label{lem:bar-cobar-associated-graded}
The associated graded of the bar-cobar filtration is:
\[\text{Gr}^p\Omega(\bar{B}(\mathcal{A})) = \Omega^p(\bar{B}^p(\mathcal{A}))\]

The induced differential on $\text{Gr}^\bullet$ is the \emph{bar complex differential} $d_{\bar{B}} = d_{\text{internal}} + d_{\text{bar}}$ (i.e., the full differential on $\bar{B}(\mathcal{A})$, lifted to each cobar tensor factor).
\end{lemma}

\begin{proof}
By definition of associated graded:
\[\text{Gr}^p = F^p / F^{p+1} = \Omega^p(\bar{B}^p(\mathcal{A}))\]

The full differential on $\Omega(\bar{B}(\mathcal{A}))$ decomposes as
$d = d_{\bar{B}} + d_{\text{cobar}}$, where
$d_{\bar{B}} = d_{\text{internal}} + d_{\text{bar}}$ is the lift of
the bar complex differential to each cobar tensor factor and preserves
cobar degree~$p$.
The comultiplication term $d_{\text{cobar}}$ raises cobar degree by~$1$,
so it sends $F^p$ into $F^{p+1}$ and vanishes on
$F^p/F^{p+1}=\text{Gr}^p$.
Therefore the induced differential on the associated graded is exactly
$d_{\bar{B}}$, which is the $d_0$ differential of the spectral
sequence.
\end{proof}

\subsection{Spectral sequence construction}

\begin{theorem}[Bar-cobar spectral sequence; \ClaimStatusConditional]\label{thm:bar-cobar-spectral-sequence}
\label{thm:koszul-spectral-sequence}
\textup{[Regime: square-zero total differential on the filtered
bar-cobar complex.]}
Assume the total differential on $\Omega(\bar{B}(\mathcal{A}))$
satisfies $d^2 = 0$ and preserves the filtration from
Definition~\ref{def:bar-cobar-filtration}. Then the filtration from
Definition~\ref{def:bar-cobar-filtration} induces a spectral
sequence:
\[E_0^{p,q} = \left(F^p / F^{p+1}\right)^{p+q} \implies H^{p+q}(\Omega(\bar{B}(\mathcal{A})))\]
converging to the cohomology of $\Omega(\bar{B}(\mathcal{A}))$. Here $p$ is the filtration degree (cobar degree) and $q$ is the complementary degree (from the internal and bar gradings), so the total degree is $p + q$.

\emph{Explicit description of pages.}
\begin{align*}
E_0^{p,q} &= \left(F^p\Omega^{p+q}(\bar{B}(\mathcal{A})) \big/ F^{p+1}\Omega^{p+q}(\bar{B}(\mathcal{A}))\right) \quad \text{(associated graded)} \\
E_1^{p,q} &= H^{p+q}\!\left(E_0^{p,\bullet},\, d_{\bar{B}}\right)
 \quad \text{(bar cohomology at fixed cobar degree $p$)} \\
E_2^{p,q} &= H^p\!\left(E_1^{\bullet,q},\, d_{\text{cobar}}\right)
 \quad \text{(cohomology w.r.t.\ comultiplication-induced differential)} \\
E_\infty^{p,q} &= \mathrm{Gr}^p H^{p+q}(\Omega(\bar{B}(\mathcal{A}))) \quad \text{(limiting page)}
\end{align*}
\end{theorem}

\begin{proof}
This is a standard spectral sequence associated to a filtered complex (Weibel~\cite{Weibel94}, Chapter~5). We verify the key properties:

\emph{Step 1: $E_0$ page.} This is just the raw complex with its bigrading:
\[E_0^{p,q} = F^p\Omega^{p+q}(\bar{B}(\mathcal{A})) / F^{p+1}\Omega^{p+q}(\bar{B}(\mathcal{A}))\]

By definition of the filtration, this isolates the piece in cobar degree~$p$ and total degree $p+q$.

\emph{Step 2: $d_0$ differential.} On the $E_0$ page:
\[d_0: E_0^{p,q} \to E_0^{p,q+1}\]
is the bar complex differential $d_{\bar{B}} = d_{\text{internal}} + d_{\text{bar}}$, lifted to each cobar tensor factor. This is the filtration-preserving component of $d$: both $d_{\text{internal}}$ (from the differential on~$\mathcal{A}$) and $d_{\text{bar}}$ (from the multiplication of~$\mathcal{A}$, encoding collisions) preserve cobar degree.

Taking cohomology gives the $E_1$ page: bar cohomology at each fixed cobar degree.

\emph{Step 3: $d_1$ differential.} On the $E_1$ page:
\[d_1: E_1^{p,q} \to E_1^{p+1,q}\]
is induced by $d_{\text{cobar}}$ (the comultiplication on $\bar{B}(\mathcal{A})$, which increases cobar degree by~$1$).

Taking cohomology gives the $E_2$ page.

\emph{Step 4: Higher differentials.} For $r \geq 2$:
\[d_r: E_r^{p,q} \to E_r^{p+r,q-r+1}\]

These differentials encode higher-order interactions between bar and cobar operations.

\emph{Step 5: Convergence.} The spectral sequence converges because:
For the application to Theorem~\ref{thm:bar-cobar-inversion-qi}, the completeness statement in Definition~\ref{def:bar-cobar-filtration}(2) is supplied by its conilpotency/completion hypothesis: if $\bar{B}(\mathcal{A})$ is conilpotent in the sense of Definition~\ref{def:conilpotent-complete}, every element has finite bar degree, while in the augmentation-ideal completion regime the filtration is complete by construction.
\begin{enumerate}
\item The filtration is complete: $\bigcap_p F^p = 0$
\item The filtration is exhaustive: $\bigcup_p F^p = \Omega(\bar{B}(\mathcal{A}))$
\item The complex is bounded in each column (fixed $p$)
\end{enumerate}

By standard spectral sequence theory (Weibel \cite{Weibel94}, Chapter 5), this ensures:
\[E_\infty^{p,q} = \text{Gr}^p H^{p+q}(\Omega(\bar{B}(\mathcal{A})))\]
\end{proof}

\begin{theorem}[Collapse at \texorpdfstring{$E_2$}{E2}; \ClaimStatusConditional]\label{thm:spectral-sequence-collapse}
\textup{[Regime: quadratic, Koszul locus, square-zero total differential
\textup{(}Convention~\textup{\ref{conv:regime-tags})}.]}

For the Heisenberg algebra, the spectral sequence collapsed trivially because $d_{\mathrm{bracket}} = 0$ (no simple poles): the bar cohomology was concentrated in degrees $0$ and $2$ without further work (\S\ref{sec:frame-bar-all}). For non-abelian algebras, the collapse is a consequence of the Koszul property.

For a Koszul chiral algebra $\mathcal{A}$ (Remark~\ref{rem:koszul-chiral}) on the
strict square-zero surface of Theorem~\ref{thm:bar-cobar-spectral-sequence}, the spectral sequence from
Theorem~\ref{thm:bar-cobar-spectral-sequence} collapses at the $E_2$ page:
\[E_2^{p,q} = E_\infty^{p,q}.\]
\end{theorem}

\begin{proof}
The Koszul property gives a quasi-isomorphism
$\iota\colon \mathcal{A}^! \hookrightarrow \bar{B}(\mathcal{A})$ from
the quadratic dual into the bar complex
\cite[Theorem~3.4.1]{LV12}; Chapter~\ref{chap:koszul-pairs} gives the
chiral version used here.
Because $\mathcal{A}^!$ is generated in weight~$1$ with quadratic
relations, the bar cohomology is concentrated on the diagonal in the
bigrading by cobar degree and internal degree.
Thus on the $E_1$ page, where the differential is
$d_0=d_{\bar B}$, one has $E_1^{p,q}=0$ for $q\neq 0$.
The $d_1$ differential comes from $d_{\text{cobar}}$, so the same
diagonal concentration persists to $E_2$.
For $r\ge 2$, the differential
$d_r\colon E_r^{p,q}\to E_r^{p+r,q-r+1}$ changes the second grading,
so with $E_2^{p,q}$ supported only at $q=0$, either its source or its
target vanishes.
Hence every $d_r$ for $r\ge 2$ is zero, and the spectral sequence
collapses at $E_2$.
\end{proof}

\begin{remark}[Square-zero boundary of the spectral-sequence proof]
\label{rem:spectral-sequence-square-zero-boundary}
Theorems~\ref{thm:bar-cobar-spectral-sequence}
and~\ref{thm:spectral-sequence-collapse} are statements about filtered
complexes with square-zero total differential. In this chapter they
apply to the strict genus-$0$ bar-cobar complex, and to any
completed/corrected model whose total differential squares to zero.
They do \emph{not} apply to the curved fiberwise higher-genus
differential $\dfib$ with
$\dfib^{\,2} = \kappa(\cA)\cdot\omega_g$. Off the Koszul locus, the
proved replacement is the coderived bar-degree filtration of
Proposition~\ref{prop:coderived-bar-degree-spectral-sequence},
together with the intrinsic coderived/contraderived bar-coalgebra
comparison of Theorem~\ref{thm:positselski-chiral-finite-type}. Ordinary
derived reduction remains available only on the flat finite-type
completed-dual loci of Theorem~\ref{thm:completed-dual-finite-type-reduction}
or when the curvature collapses.
\end{remark}

\subsection{Convergence at all genera}

\begin{theorem}[Formal assembly of genuswise quadratic comparisons;
\ClaimStatusConditional]
\label{thm:genus-graded-convergence}
\textup{[Type signature: Open modular extension; completed
genus-graded presentation; Beilinson levels
\(1\leftrightarrow2\); hypothesis package:
\(\{H_{\mathrm{CL}}^{(g)}\}_{g\geq0}\), graph-compatible
transition maps, and \(\hbar\)-adic completeness.]}
For every \(g\geq0\), suppose a modular quadratic coalgebra
\(A_g^{\mathrm i}\), a bar object \(B_g(A)\), and a comparison
\[
q_{A,g}\colon A_g^{\mathrm i}\xrightarrow{\sim}B_g(A)
\]
have been constructed, and suppose the comparisons commute with the
genus transition and clutching maps.  Then
\[
q_A(\hbar):=\prod_{g\geq0}
\hbar^{2g}\,q_{A,g}
\]
is a quasi-isomorphism in the \(\hbar\)-adically complete product
category.  Applying continuous cobar gives the corresponding formal
family of algebra comparisons
\[
\prod_{g\geq0}\hbar^{2g}\Omega_g(A_g^{\mathrm i})
\xrightarrow{\sim}
\prod_{g\geq0}\hbar^{2g}A_g .
\]
\end{theorem}

\begin{proof}
Modulo \(\hbar^{2N+2}\), the statement is the finite product of the
quasi-isomorphisms \(q_{A,g}\) for \(0\leq g\leq N\).
The transition maps are strict and the cohomology tower is
Mittag--Leffler, so the derived inverse limit of these finite products
has zero comparison cone.  Continuity of cobar is part of the
genuswise comparison package.
\end{proof}

\subsection{The counit of the adjunction}

\begin{proposition}[Coalgebra-side universal unit;
\ClaimStatusConditional]
\label{prop:counit-qi}
Let \(C\) be an uncurved complete conilpotent chiral coalgebra
satisfying
\(H_{\mathrm{CL}}(\Omega_XC,C,\tau_{\mathrm{univ}})\).
Then
\[
\eta_C\colon C\longrightarrow B_X\Omega_X(C)
\]
is a quasi-isomorphism.  For a curved twisting datum
\((A,C,\tau)\), the corresponding algebra and module comparisons are
governed by
\(\mathsf{Tw}^{\mathrm{ch}}_{\mathrm{acyc}}(C,A,\tau)\).
\end{proposition}

\begin{proof}
Loday--Vallette \cite[Corollary~2.3.4]{LV12} prove the ordinary
conilpotent unit theorem.  The chiral Comparison-Lemma package
transports it across the finite windows, collision strata, and
completed inverse limit.
\end{proof}

\subsection{Functoriality of the quasi-isomorphism}

\begin{theorem}[Naturality of universal reconstruction;
\ClaimStatusProvedHere]
\label{thm:bar-cobar-inversion-functorial}
For every morphism \(f\colon A\to A'\) of augmented chiral algebras,
the counit square commutes:
\begin{center}
\begin{tikzcd}
\Omega_XB_X(A) \ar[r, "\varepsilon_A"]
 \ar[d, "\Omega_XB_X(f)"']
& A \ar[d, "f"] \\
\Omega_XB_X(A') \ar[r, "\varepsilon_{A'}"']
& A'.
\end{tikzcd}
\end{center}
When \(A\) and \(A'\) carry reconstruction packages, this square is
a natural equivalence in the selected completed factorization
ambient.
\end{theorem}

\begin{proof}
The counit of an adjunction is a natural transformation
\(\Omega_XB_X\Rightarrow\operatorname{id}\).  The displayed square is
its naturality square.
\end{proof}

\subsection{Applications to module categories}

\begin{remark}[Module Koszul duality on the quadratic genus-zero surface]
\label{rem:derived-equiv-scope}
The intrinsic second-kind comparison in this chapter is the
bar-coalgebra statement of
Theorems~\ref{thm:positselski-chiral-finite-type} and
\ref{thm:completed-dual-finite-type-reduction}: it compares
comodules and contramodules over one fixed bar coalgebra.
Theorem~\ref{thm:e1-module-koszul-duality} supplies the
module-theoretic comparison on the quadratic genus-zero bar-dual
component.  A further module-category realization is therefore
recorded together with its acyclic twisting and completed-dual
hypotheses.
\end{remark}

\begin{remark}[Homotopy information carried by the counit]
\label{rem:qi-vs-homology-iso-counit}% renamed 2026-05-17 from rem:qi-vs-homology-iso (collision with L2299)
The bar-cobar counit used in
Theorem~\ref{thm:bar-cobar-platonic} is a morphism of chiral
algebras. Its quasi-isomorphism class preserves the OPE operations and
their coherent homotopies.  The induced cohomology isomorphism is its
underlying graded shadow; the algebra morphism also carries the
twisting and brace operations.
\end{remark}

%================================================================
% COMPLETED BAR-COBAR INVERSION
%================================================================

\subsection{Completed bar-cobar inversion}
\label{subsec:completed-bar-cobar-inversion}

Finite-stage universal reconstruction is supplied by
Theorem~\ref{thm:bar-cobar-platonic}.  Passing this reconstruction
through a separated complete inverse limit requires a strict tower and
Mittag--Leffler control.  The key structural input is the strong
completion-tower framework developed in
\S\ref{sec:curved-koszul-i-adic}
(Definition~\ref{def:strong-completion-tower},
Theorem~\ref{thm:completed-bar-cobar-strong}); what follows gives the
homological comparison lemma that underlies the inverse-limit passage
and develops the inversion-side consequences of the completion closure
$\CompCl(\Fft)$ (Definition~\ref{def:completion-closure}).

\begin{lemma}[Complete filtered comparison lemma; \ClaimStatusProvedHere]
\label{lem:complete-filtered-comparison}
\index{complete filtered comparison lemma|textbf}
\index{Mittag--Leffler!complete filtered comparison}
Let $f\colon C \to D$ be a morphism of separated complete filtered
cochain complexes, with descending filtrations
$F^\bullet C$ and $F^\bullet D$ satisfying:
\begin{enumerate}[label=\textup{(\roman*)}]
\item the quotient towers are surjective: for every~$N$, the
 transition maps $C/F^{N+2}C \twoheadrightarrow C/F^{N+1}C$ and
 $D/F^{N+2}D \twoheadrightarrow D/F^{N+1}D$ are surjections;
\item every quotient map is a quasi-isomorphism:
 \[
 f_N \colon C/F^{N+1}C \xrightarrow{\;\sim\;} D/F^{N+1}D
 \]
 for all~$N \ge 0$.
\end{enumerate}
Then $f$ is a quasi-isomorphism.
\end{lemma}

\begin{proof}
Consider the mapping cone $\operatorname{Cone}(f)$, which inherits a
separated complete filtration from~$C$ and~$D$. Each quotient
$\operatorname{Cone}(f)/F^{N+1} \cong \operatorname{Cone}(f_N)$ is
acyclic because $f_N$ is a quasi-isomorphism by hypothesis~(ii).
The Milnor exact sequence for the inverse system
$\{\operatorname{Cone}(f_N)\}_{N \ge 0}$ reads
\[
0 \to \varprojlim\nolimits^1_N
 H^{m-1}(\operatorname{Cone}(f_N))
\to H^m(\operatorname{Cone}(f))
\to \varprojlim_N H^m(\operatorname{Cone}(f_N)) \to 0.
\]
The right-hand term vanishes because every
$H^m(\operatorname{Cone}(f_N)) = 0$. For the left-hand term: the
quotient tower is surjective by hypothesis~(i), so the transition maps
on each cohomology group form a surjective inverse system of zero
groups, and the Mittag--Leffler condition is satisfied trivially.
Hence $\varprojlim^1 = 0$. It follows that
$H^m(\operatorname{Cone}(f)) = 0$ for all~$m$, so $f$ is a
quasi-isomorphism.
\end{proof}

\begin{remark}[Role in the completed bar-cobar theorem]
\label{rem:complete-filtered-comparison-role}
Lemma~\ref{lem:complete-filtered-comparison} is the homological
workhorse of the completed bar-cobar theorem
(Theorem~\ref{thm:completed-bar-cobar-strong}). In that application,
$C = \widehat\Omega^{\mathrm{ch}}(\widehat{\bar B}^{\mathrm{ch}}(\cA))$,
$D = \cA$, and the quotient quasi-isomorphisms $f_N$ are the
finite-stage counits
$\epsilon_N\colon \Omega^{\mathrm{ch}}(\bar B^{\mathrm{ch}}(\cA_{\le N}))
\xrightarrow{\sim} \cA_{\le N}$, which are quasi-isomorphisms by the
finite-type regime. The surjectivity of the quotient tower follows from
the strong completion-tower axiom: the projections
$\cA_{\le N+1} \twoheadrightarrow \cA_{\le N}$ are strict surjections
by construction. The lemma then gives the completed counit
$\widehat\epsilon$ is a quasi-isomorphism; this is Step~3 of the proof of
Theorem~\ref{thm:completed-bar-cobar-strong}.
\end{remark}

\begin{remark}[Completed inversion: the inversion-side summary]
\label{rem:completed-inversion-summary}
The completion closure theory
(Definition~\ref{def:completion-closure}--Theorem~\ref{thm:uniform-pbw-bridge},
developed in \S\ref{sec:curved-koszul-i-adic}) gives a complete
inversion package for the infinite-type regime:
\begin{enumerate}
\item \emph{Domain.}
 $\CompCl(\Fft)$ is the class of strong completion towers
 (Definition~\ref{def:strong-completion-tower}) with finite-type
 quotients.
\item \emph{Inversion.}
 Completed bar and cobar are quasi-inverse on $\CompCl(\Fft)$
 (Theorem~\ref{thm:completed-bar-cobar-strong}).
\item \emph{Homotopy category.}
 Localizing by quotientwise quasi-isomorphisms gives an equivalence
 of homotopy categories
 (Corollary~\ref{cor:completion-closure-equivalence}).
\item \emph{Stability under twisting.}
 If $\alpha \in F^1\cA^1$ is a continuous Maurer--Cartan element,
 then $\cA_\alpha \in \CompCl(\Fft)$ and inversion persists
 (Theorem~\ref{thm:mc-twisting-closure}).
\item \emph{Representability.}
 Algebra and coalgebra maps between completed objects are represented
 by continuous twisting morphisms
 (Theorem~\ref{thm:completed-twisting-representability}).
\item \emph{Verification criteria.}
 Coefficient stability
 (Theorem~\ref{thm:coefficient-stability-criterion}) reduces
 completed inversion to stabilization of finite matrix entries on
 windows; the uniform PBW bridge
 (Theorem~\ref{thm:uniform-pbw-bridge}) connects PBW concentration
 (MC1) directly to completed duality (MC4).
\end{enumerate}
Together with finite-stage reconstruction
(Theorem~\ref{thm:bar-cobar-platonic}), these results give completed
reconstruction.  Quadratic recognition remains the
comparison \(A^{\mathrm i}\to B_X(A)\) of
Theorem~\ref{thm:bar-cobar-inversion-qi}.
\end{remark}

\begin{corollary}[Completed derived equivalence; \ClaimStatusConditional]
\label{cor:completed-derived-equivalence}
\index{derived equivalence!completed}
Let $\cA \in \CompCl(\Fft)$ be a strong completion tower on a
smooth curve~$X$. Then the completed bar and cobar functors induce
an equivalence of triangulated categories
\[
\mathcal{D}^b(\operatorname{Mod}^{\mathrm{cts}}(\cA))
\;\simeq\;
\mathcal{D}^b(\operatorname{Comod}^{\mathrm{pronil}}
(\widehat{\bar B}^{\mathrm{ch}}(\cA))),
\]
where $\operatorname{Mod}^{\mathrm{cts}}$ denotes continuous modules
over the complete algebra and
$\operatorname{Comod}^{\mathrm{pronil}}$ denotes pronilpotent
comodules over the completed bar coalgebra. The twisting morphism
$\widehat\tau\colon \widehat{\bar B}^{\mathrm{ch}}(\cA) \to \cA$
mediates the equivalence via completed extension of scalars.
\end{corollary}

\begin{proof}
The completed counit and unit are quasi-isomorphisms by
Theorem~\ref{thm:completed-bar-cobar-strong}(3)--(4). The
passage to module/comodule categories follows from completed twisting
representability
(Theorem~\ref{thm:completed-twisting-representability}), which
ensures that the completed bar and cobar functors on
module categories are adjoint and that the unit/counit natural
transformations are quasi-isomorphisms quotientwise. Applying
Lemma~\ref{lem:complete-filtered-comparison} to each module/comodule
then gives the equivalence. The argument parallels
Theorem~\ref{thm:completed-dual-finite-type-reduction}, with
finite-generation replaced by continuity and conilpotence replaced
by pronilpotence.
\end{proof}

\begin{remark}[Four comparison ambients]
\label{rem:four-regime-inversion}
\index{bar-cobar inversion!four-regime table}
The following table summarizes the inversion theorem across all four
regimes (Convention~\ref{conv:regime-tags}):
\begin{center}
\renewcommand{\arraystretch}{1.3}
\begin{tabular}{@{}llll@{}}
\toprule
\textbf{Surface} & \textbf{Input} & \textbf{Ambient}
 & \textbf{Comparison} \\
\midrule
Quadratic &
 \(H_{\mathrm{CL}}(A,A^{\mathrm i},\tau_{\mathrm i})\) &
 complete coalgebras &
 Thm~\ref{thm:bar-cobar-inversion-qi} \\
Curved-central &
 \(\mathsf{Pos}^{\mathrm{ch}}_{\mathrm{co-ctr}}(C)\) &
 coderived/contraderived &
 Thm~\ref{thm:off-koszul-ran-inversion} \\
Filtered-complete &
 finite-window reconstruction package &
 complete factorization algebras &
 Prop~\ref{prop:finite-window-ml-inversion-bci} \\
Strict tower &
 strong completion tower &
 $D^b(\mathrm{Mod}^{\mathrm{cts}})$ &
 Thm~\ref{thm:completed-bar-cobar-strong} \\
\bottomrule
\end{tabular}
\end{center}
Each row states its own source, target, and hypothesis package.  The
complete filtered comparison lemma
(Lemma~\ref{lem:complete-filtered-comparison}) passes
quotientwise reconstruction maps to the strict inverse limit.
\end{remark}

%================================================================
% GEOMETRIC AND CATEGORICAL CHARACTERIZATIONS OF CHIRAL KOSZULNESS
%================================================================

\subsection{Geometric and categorical characterizations of chiral Koszulness}
\label{subsec:geometric-categorical-koszulness}

Theorem~\ref{thm:bar-cobar-inversion-qi} identifies quadratic chiral
Koszulness with the comparison
\(A^{\mathrm i}\to B_X(A)\).  The following subsections study
additional descent, factorization-homology, boundary, Hodge, and
shifted-symplectic criteria.  Every criterion carries the comparison
map and convergence package that connect it to this quadratic core.

\medskip

\subsubsection{Barr--Beck--Lurie characterization}

Lurie's $\infty$-categorical Barr--Beck theorem~\cite[Theorem~4.7.3.5]{HA}
characterizes (co)monadic adjunctions by conservativity and totalization
preservation.  For any proposed comparison functor in this chapter,
these two conditions form an abstract descent package.  A relation to
quadratic Koszulness requires an additional theorem identifying the
resulting descent comparison with
$q_A\colon A^{\mathrm i}\to B_X(A)$.

\begin{theorem}[Barr--Beck--Lurie comparison for a selected
factorization domain; \ClaimStatusConditional]
\label{thm:barr-beck-lurie-koszulness}
\textup{[Type signature: Open quadrant; factorization
\(\infty\)-categorical presentation; Beilinson levels
\(1\leftrightarrow2\); hypothesis package: a specified
bar--cobar adjunction, conservativity of \(B_X\), and preservation of
\(B_X\)-split totalizations.]}
Let
\[
\Omega_X\colon\mathcal C_{\mathrm{coalg}}
\rightleftarrows
\mathcal C_{\mathrm{alg}}\colon B_X
\]
be the selected factorization bar--cobar adjunction.  If \(B_X\) is
conservative on \(\mathcal C_{\mathrm{alg}}\) and preserves
totalizations of \(B_X\)-split cosimplicial objects, then the
Barr--Beck--Lurie comparison functor
\[
\mathcal C_{\mathrm{alg}}
\longrightarrow
\operatorname{Coalg}_{B_X\Omega_X}
(\mathcal C_{\mathrm{coalg}})
\]
is an equivalence.
\end{theorem}

\begin{proof}
This is the comonadic Barr--Beck theorem
\cite[Theorem~4.7.3.5]{HA}, applied to the displayed adjunction.
The two hypotheses in the statement are precisely its
conservativity and split-totalization conditions.
\end{proof}

\begin{remark}[Relation to quadratic recognition]
\label{rem:bbl-content}
The Barr--Beck theorem reconstructs objects from the comonad
\(B_X\Omega_X\) on a chosen categorical domain.  Theorem~B compares
the quadratic coalgebra \(A^{\mathrm i}\) with the full bar coalgebra
\(B_X(A)\).  A theorem connecting these two surfaces must identify
the comonadic resolution with the filtered quadratic comparison and
must verify the strong-convergence package
\(H_{\mathrm{CL}}(A,A^{\mathrm i},\tau_{\mathrm i})\).
\end{remark}

\medskip

\subsubsection{Factorization homology concentration}

Proposition~\ref{prop:bar-fh} identifies the bar complex with a
factorization-homology model only under constructibility, compactness,
and filtration hypotheses. On that comparison surface, Koszulness has
a homological test. Outside it, factorization homology is a comparison
object, not a definition of Koszulness and not a substitute for the
bar-cobar counit.

\begin{theorem}[Factorization-homology transport of the quadratic
comparison; \ClaimStatusConditional]
\label{thm:fh-concentration-koszulness}
\textup{[Type signature: Open quadrant; filtered
factorization-homology presentation; Beilinson levels
\(1\leftrightarrow2\); hypothesis package:
\(\mathsf{FH}_{\Sigma}(A)\) and a filtered detecting comparison.]}
Let \(\mathsf{FH}_{\Sigma}(A)\) be a selected chain model for
factorization homology and suppose a filtered quasi-isomorphism
\[
\chi_\Sigma\colon B_X(A)\xrightarrow{\sim}
\mathsf{FH}_{\Sigma}(A)
\]
is compatible with collision filtrations and completed inverse
limits.  Then
\[
q_A\colon A^{\mathrm i}\xrightarrow{\sim}B_X(A)
\quad\Longleftrightarrow\quad
\chi_\Sigma q_A\colon A^{\mathrm i}
\xrightarrow{\sim}\mathsf{FH}_{\Sigma}(A).
\]
Consequently every cohomological support statement proved for one
side transports to the other.
\end{theorem}

\begin{proof}
The map \(\chi_\Sigma\) is a quasi-isomorphism by hypothesis.
The equivalence of the two displayed conditions is the
two-out-of-three property.  A quasi-isomorphism identifies
cohomological support.
\end{proof}

\begin{remark}[Configuration filtration on the comparison model]
\label{rem:goodwillie-tower-fh}
Under the constructibility, compactness, filtration, and
Mittag--Leffler hypotheses of
Proposition~\ref{prop:bar-fh}, the selected factorization-homology
model has a filtration by configuration number.  Its associated
graded is assembled from the configuration terms
\[
\operatorname{Conf}_n(X)\otimes_{\Sigma_n}
(s^{-1}\bar A)^{\otimes n}.
\]
The filtered comparison \(\chi_\Sigma\) of
Theorem~\ref{thm:fh-concentration-koszulness} identifies this
spectral sequence with the filtered bar model.  Hence the quadratic
map \(q_A\) may be tested after transport to the configuration
filtration.
\end{remark}

\begin{remark}[Poincar\'e/Koszul duality]\label{rem:ayala-francis-poincare-koszul}
Ayala--Francis~\cite{AF15} supply the Poincar\'e/Koszul-duality
framework for comparing factorization homology and cohomology. In this
chapter it can be used only after the same finite-type Verdier,
bar-comparison, and detecting hypotheses have identified the
post-Verdier object with the strict model denoted~\(\cA^!\). Under
those hypotheses, and on the Koszul locus,
\(Q_g(\cA)=H^0(\int_{\Sigma_g}\cA)\); without them, conformal blocks are
not identified with the full chiral homology by this remark.
\end{remark}

\medskip

\subsubsection{Fulton--MacPherson boundary acyclicity}

The bar complex lives on $\overline{\Conf}_n(X)$; restriction to
boundary strata encodes OPE residues. Koszulness is equivalent
to acyclicity of all such restrictions.

\begin{theorem}[Fulton--MacPherson boundary detection;
\ClaimStatusConditional]
\label{thm:fm-boundary-acyclicity}
\textup{[Type signature: Open quadrant; constructible
Fulton--MacPherson presentation; Beilinson levels
\(1\leftrightarrow2\); hypothesis package: constructibility of
\(\operatorname{Cone}(q_A)\), conservative open-and-boundary
restriction, and \(t\)-exact filtered base change.]}
Let
\[
K_A:=\operatorname{Cone}\!\left(
q_A\colon A^{\mathrm i}\longrightarrow B_X(A)\right)
\]
be constructible for the Fulton--MacPherson stratification.  Assume
that restriction to the open configuration locus together with all
functors \(i_S^!\) for boundary strata is conservative on the
subcategory containing \(K_A\).  Then \(q_A\) is a quasi-isomorphism
if and only if
\[
j^*K_A\simeq0,
\qquad
i_S^!K_A\simeq0
\quad\text{for every boundary stratum }S.
\]
When the open comparison is the ordinary quadratic Koszul complex,
the first condition is supplied by the local quadratic calculation,
and the boundary conditions measure the collision-gluing extension.
\end{theorem}

\begin{proof}
A morphism is an equivalence precisely when its cone vanishes.
Conservativity of the stated family of restriction functors detects
the vanishing of \(K_A\).  The final sentence identifies the open
restriction through the supplied filtered base-change comparison.
\end{proof}

\medskip
\medskip

\subsubsection{D-module purity criterion}

For regular holonomic bar components, the classical
Beilinson--Ginzburg--Soergel purity theorem~\cite{BGS96} motivates the
following chiral comparison problem.

\begin{conjecture}[D-module purity criterion;
\ClaimStatusConjectured]\label{conj:d-module-purity-koszulness}
\index{Koszulness!D-module purity criterion}
\index{D-module!purity and Koszulness}
\index{characteristic variety!alignment with FM stratification}
Let $\cA$ be a chiral algebra on~$X$ satisfying the standing hypotheses.
Then $\cA$ is chirally Koszul if and only if the following two conditions
hold for all $n \geq 1$:
\begin{enumerate}[label=\textup{(\roman*)}]
\item \textup{(Purity)}\quad
 Each bar component $\barBch_n(\cA)$, viewed as a mixed Hodge module on
 $\overline{\Conf}_n(X)$, is \emph{pure} of weight~$n$.
\item \textup{(Characteristic variety alignment)}\quad
 The characteristic variety
 $\operatorname{Ch}(\barBch_n(\cA)) \subset T^*\overline{\Conf}_n(X)$
 is contained in the union of conormal bundles to boundary strata:
 \[
 \operatorname{Ch}\bigl(\barBch_n(\cA)\bigr) \;\subset\;
 \bigcup_{S \subset \partial \overline{\Conf}_n(X)}
 T^*_S \overline{\Conf}_n(X).
 \]
\end{enumerate}
\end{conjecture}

\begin{remark}[Purity comparison package]
\label{rem:d-module-purity-content-bci}
The BGS theorem identifies Koszul diagonal concentration with a purity
statement in its graded geometric setting
\cite{BGS96}.  A chiral analogue requires three additional maps:
a mixed-Hodge realization of every bar component, a
Fulton--MacPherson base-change theorem for \(i_S^!\), and an
identification of the weight filtration with the PBW/bar filtration.
With these maps, purity and characteristic-variety alignment transport
to the boundary detection criterion of
Theorem~\ref{thm:fm-boundary-acyclicity}.  Constructing this package
and proving its converse constitute
Conjecture~\ref{conj:d-module-purity-koszulness}.
\end{remark}

\medskip
\medskip

\subsubsection{Lagrangian criterion in the shifted-symplectic deformation space}

The ambient complementarity criterion
(Theorem~\ref{thm:ambient-complementarity};
Theorem~\ref{thm:ambient-complementarity-fmp}) equips the modular
deformation space with a $(-1)$-shifted symplectic structure.
Koszulness is conjectural Lagrangian transversality.

\begin{conjecture}[Conditional Lagrangian characterization of chiral Koszulness;
\ClaimStatusConjectured]%
\label{conj:lagrangian-koszulness}%
\index{Koszulness!Lagrangian characterization}
\index{Lagrangian!transversality and Koszulness}
\index{shifted-symplectic!Lagrangian criterion for Koszulness}
Suppose the full shifted-symplectic Lagrangian package is supplied:
\begin{enumerate}[label=\textup{(L\arabic*)}]
\item the ambient complementarity tangent complex
 \(T_{\mathrm{comp}}(\cA)\) is perfect;
\item the shifted pairing is non-degenerate, giving a
 \((-1)\)-shifted symplectic structure~\(\omega\) as in
 Theorem~\textup{\ref{thm:ambient-complementarity-fmp}};
\item the cyclic/BV compatibility needed to compare the bar chart with
 the ambient deformation complex holds;
\item the one-sided deformation maps
\[
 i_{\cA} \colon \mathcal{M}_{\cA} \longrightarrow \mathcal{M}_{\mathrm{comp}},
 \qquad
 i_{\cA^!} \colon \mathcal{M}_{\cA^!} \longrightarrow \mathcal{M}_{\mathrm{comp}}
\]
are isotropic embeddings;
\item the derived intersection
\(\mathcal{M}_\cA\times^h_{\mathcal{M}_{\mathrm{comp}}}
\mathcal{M}_{\cA^!}\) is represented by the completed twisted tensor
product \(K_\tau(\cA,\cA^!)\);
\item the relevant bar-chart acyclicity and completed bar-cobar
comparison hypotheses of Theorem~\textup{\ref{thm:bar-cobar-inversion-qi}}
hold.
\end{enumerate}
Then $\cA$ is chirally Koszul if and only if both
$i_{\cA}$ and $i_{\cA^!}$ are \emph{Lagrangian} \textup{(}isotropic of maximal
dimension\textup{)}. Equivalently: the bar-cobar adjunction defines a Lagrangian
correspondence in the $(-1)$-shifted symplectic deformation space.
\end{conjecture}

\begin{remark}[Lagrangian criterion]
The conjecture has the following exact content. In the forward
direction, on the Koszul locus, the bar-cobar adjunction gives a
free resolution and the complementarity splitting
$\mathbf{Q}_g(\cA) \oplus \mathbf{Q}_g(\cA^!) \simeq
H^*(\overline{\cM}_g, \cZ_\cA)$ (Theorem~C) identifies $\cM_\cA$
and $\cM_{\cA^!}$ as complementary subspaces. The
shifted-symplectic structure
(Theorem~\ref{thm:ambient-complementarity-fmp}) makes them
isotropic; complementarity gives maximal dimension, hence
Lagrangian.

The converse requires an additional identification: the derived
intersection
\[
\mathcal{M}_\cA\times^h_{\mathcal{M}_{\mathrm{comp}}}
\mathcal{M}_{\cA^!}
\]
must be represented by the twisted tensor product
$K_\tau(\cA,\cA^!)$. Once this comparison is present, transverse
Lagrangians in a $(-1)$-shifted symplectic space have derived
intersection of expected dimension~$0$, and the resulting acyclicity
of $K_\tau(\cA,\cA^!)$ is precisely Koszulness.
This is a Koszul-locus argument. Off the Koszul locus, the chapter
proves only the intrinsic coderived/contraderived bar-coalgebra
comparison of Theorem~\ref{thm:positselski-chiral-finite-type}, not an
ordinary twisted-tensor-product or free-resolution statement.
\end{remark}

\begin{remark}[PTVV interpretation]\label{rem:lagrangian-ptvv}
In the PTVV framework~\cite{PTVV13}, the derived intersection
$\mathcal{M}_\cA \times^h_{\mathcal{M}_{\mathrm{comp}}}
\mathcal{M}_{\cA^!}$ computes the twisted tensor product
$K_\tau(\cA, \cA^!)$ on the Koszul surface of
Theorem~\ref{thm:bar-cobar-inversion-qi}; its acyclicity
\textup{(}equivalent to Koszulness\textup{)} is
precisely the transversality of the two Lagrangians.
The proved linear shadow is Theorem~C:
$Q_g(\cA) \oplus Q_g(\cA^!) \cong
H^*(\overline{\mathcal{M}}_g, \mathcal{Z}(\cA))$.
The off-locus coderived surface is not identified here with an
ordinary twisted tensor product.
\end{remark}

\begin{remark}[Conditional hypotheses and their verification]
\label{rem:lagrangian-conditional}
Conjecture~\ref{conj:lagrangian-koszulness} is conditional on the
shifted-symplectic hypothesis package: perfectness of the deformation
complex, non-degeneracy of the shifted-symplectic pairing, finiteness
of the modular envelope, the cyclic/BV compatibility inputs, and the
bar-chart acyclicity inputs. The tangent-level
shifted-symplectic structure
(Theorem~\ref{thm:ambient-complementarity-fmp}) is proved
at the ambient cyclic-deformation level; the characterization of Koszulness as Lagrangian
transversality then follows from the PTVV derived-intersection
computation (Remark~\ref{rem:lagrangian-ptvv}) only after that package
is supplied. Global Lagrangian at the tangent level and globally
Lagrangian are equivalent when the derived intersection is discrete.

\emph{Verification for the standard landscape.}
Proposition~\ref{prop:lagrangian-perfectness} verifies
perfectness for the standard landscape at non-critical,
non-degenerate levels. The proof proceeds in four steps:
(1)~finite-dimensionality of the deformation complex at
each weight level (from~(P1));
(2)~chain-level nondegeneracy of the cyclic pairing
(the pairing on $\fg = \Defcyc^{\mathrm{mod}}(\cA)$ is
block-diagonal in the $(g,r)$-decomposition, and each block
is a convolution of nondegenerate pairings on finite-dimensional
spaces (the invariant form~(P2) and the cyclic
trace on the modular cooperad));
(3)~skew-adjointness of the differential with respect to the
pairing (a defining property of invariant pairings on dg~Lie
algebras);
(4)~the induced map $\fg|_n \to (\fg|_n)^\vee[-1]$
is a chain isomorphism (a chain map that is a graded isomorphism
is automatically a quasi-isomorphism), hence the pairing is
perfect.
Once perfectness holds, the PTVV principle gives
$\cM_{\mathrm{comp}}$ its $(-1)$-shifted symplectic structure;
finiteness of the modular envelope is immediate from~(P1). The
remaining cyclic/BV compatibility and bar-chart acyclicity inputs are
not consequences of perfectness.
Corollary~\ref{cor:lagrangian-conditional-standard-landscape}
therefore records the conditional Lagrangian upgrade on the
standard-landscape Koszul cases covered by
Theorem~\ref{thm:bar-cobar-inversion-qi} once the full
shifted-symplectic hypothesis package is supplied. The
filtered-complete off-locus regime remains on the intrinsic
coderived/contraderived bar-coalgebra surface of
Theorem~\ref{thm:positselski-chiral-finite-type} and is not part of
that upgrade.
\end{remark}

\begin{proposition}[Perfectness for the standard landscape]
\label{prop:lagrangian-perfectness}
\ClaimStatusProvedHere
\index{Lagrangian!perfectness verification}
\index{perfectness!cyclic pairing on complementarity datum}
Let $\cA$ satisfy the standing hypotheses and the following:
\begin{enumerate}[label=\textup{(P\arabic*)}]
\item\label{item:perf-fingen}
 \textup{(Finite weight spaces)}\quad
 The conformal weight spaces $\cA_n$ are finite-dimensional for
 all~$n$, and $\cA_n = 0$ for $n \ll 0$.
\item\label{item:perf-nondeg}
 \textup{(Nondegenerate invariant form)}\quad
 The invariant bilinear form
 $\langle{-},{-}\rangle \colon \cA \otimes \cA \to \omega_X$
 restricts to a nondegenerate pairing on each weight space~$\cA_n$.
\item\label{item:perf-dual}
 \textup{(Dual regularity)}\quad
 The Koszul dual $\cA^!$ also satisfies
 \textup{(P\ref{item:perf-fingen})} and
 \textup{(P\ref{item:perf-nondeg})}.
\end{enumerate}
Then the cyclic pairing on the ambient complementarity datum
\textup{(}Definition~\textup{\ref{def:ambient-complementarity-datum}(v))}
is perfect in the weight-graded pro-complete sense: at each conformal
weight level~$n$, the induced map
\[
(L_\cA \oplus K_\cA \oplus L_{\cA^!})|_n
\;\longrightarrow\;
(L_\cA \oplus K_\cA \oplus L_{\cA^!})^\vee[-1]\big|_n
\]
is an isomorphism of finite-dimensional vector spaces.
\end{proposition}

\begin{proof}
The total complex $L_\cA \oplus K_\cA \oplus L_{\cA^!}$ of the
ambient complementarity datum decomposes the modular cyclic
deformation complex $\fg := \Defcyc^{\mathrm{mod}}(\cA)$
(Definition~\ref{def:modular-cyclic-deformation-complex})
relative to the Koszul duality structure: $L_\cA$ governs
$\cA$-deformations, $L_{\cA^!}$ governs $\cA^!$-deformations,
$K_\cA$ governs the twisting interaction. The cyclic pairing
on $L_\cA \oplus K_\cA \oplus L_{\cA^!}$
(Definition~\ref{def:ambient-complementarity-datum}(v))
is the invariant pairing on~$\fg$ expressed in this
decomposition. We prove perfectness of this pairing
weight-by-weight.

Fix a conformal weight level~$n$.

\smallskip\noindent
\emph{Step~1: finite-dimensionality.}
By~(P\ref{item:perf-fingen}), each $\cA_m$ is finite-dimensional
and vanishes for $m \ll 0$. The complex~$\fg$ decomposes at weight~$n$ as
\[
\fg|_n
\;=\;
\bigoplus_{\substack{g,r \\ 2g-2+r>0}}
\operatorname{CoDer}^{\mathrm{cyc}}\bigl(
\barB^{(g,r)}(\cA)
\bigr)[1]\big|_n,
\]
a finite direct sum (only finitely many stable pairs $(g,r)$
contribute at weight~$n$) of finite-dimensional spaces (each
cyclic coderivation space is finite-dimensional because $\cA_m$ is,
and the bar complex at fixed $(g,r)$ is built from finite tensor
products of the~$\cA_m$). Hence $\dim \fg|_n < \infty$.

\smallskip\noindent
\emph{Step~2: chain-level nondegeneracy.}
The degree~$-1$ cyclic pairing~$\omega$ on~$\fg$
is induced from two inputs:
\begin{enumerate}[label=\textup{(\alph*)}]
\item the invariant form
 $\langle{-},{-}\rangle_\cA$ on~$\cA$, which is nondegenerate on
 each weight space by~(P\ref{item:perf-nondeg});
\item the cyclic trace $\operatorname{tr}_{(g,r)}$ on each component
 of the modular cooperad $\mathcal{P}^!_{(g,r)}$,
 which is nondegenerate by the cyclic
 modular operad axioms (Serre duality on
 $\overline{\mathcal{M}}_{g,r}$).
\end{enumerate}
At each $(g,r)$, the pairing on
$\operatorname{CoDer}^{\mathrm{cyc}}(\barB^{(g,r)}(\cA))[1]|_n
\cong \operatorname{Hom}(\mathcal{P}^!_{(g,r)},
\operatorname{End}^{\mathrm{cyc}}_\cA)|_n$
is the convolution of~(a) and~(b): explicitly,
$\omega(\alpha, \beta)
= \operatorname{tr}_{(g,r)}\bigl(
\langle \alpha(-), \beta(-) \rangle_\cA\bigr)$.
This is nondegenerate because it is the composite of the
canonical isomorphism
$\operatorname{Hom}(V, W) \cong V^\vee \otimes W$
with the tensor-product pairing, and the tensor product of two
nondegenerate finite-dimensional bilinear forms is nondegenerate.
Because the total pairing on~$\fg|_n$ is the direct sum over
$(g,r)$ of these component pairings (different $(g,r)$-components
are $\omega$-orthogonal by the grading), the total chain-level
pairing is nondegenerate.

\smallskip\noindent
\emph{Step~3: skew-adjointness of the differential.}
The differential $d = d_{\mathrm{int}} + d_{\mathrm{sew}}$
on~$\fg$ is skew-adjoint with respect to~$\omega$:
\[
\omega(d\alpha, \beta)
\;+\; (-1)^{|\alpha|}\,\omega(\alpha, d\beta)
\;=\; 0.
\]
This is the defining property of an invariant pairing on a
dg~Lie algebra and follows from
$\omega([\Theta, \alpha], \beta)
= \omega(\alpha, [\Theta, \beta])$
applied to the Maurer--Cartan element~$\Theta_\cA$ that
twists the differential.

\smallskip\noindent
\emph{Step~4: perfectness.}
The pairing~$\omega$ induces a chain map
$\varphi \colon \fg|_n \to (\fg|_n)^\vee[-1]$
defined by $\varphi(\alpha)(\beta) := \omega(\alpha, \beta)$.
That $\varphi$ is a chain map follows from the
skew-adjointness in Step~3: $\varphi \circ d = d^\vee \circ \varphi$.
By Step~2, $\varphi$ is an isomorphism of graded vector spaces
at each chain degree (nondegenerate finite-dimensional bilinear
form $\Leftrightarrow$ the induced linear map is bijective).
A chain map that is an isomorphism of underlying graded vector
spaces is a chain isomorphism, hence a fortiori a
quasi-isomorphism. Therefore the pairing on~$\fg|_n$ is perfect.

Taking the pro-complete limit over all weight levels:
$\omega$ is perfect at each~$n$, so the pro-complete
pairing on~$\fg = \varprojlim_n \fg/F_n$ is perfect.
\end{proof}

\begin{corollary}[Conditional Lagrangian criterion for the
standard landscape]
\label{cor:lagrangian-conditional-standard-landscape}
\ClaimStatusConditional
\index{Koszulness!Lagrangian criterion conditional}
\index{Lagrangian!conditional equivalence}
Under hypotheses \textup{(P1)--(P3)}, the remaining inputs of the
shifted-symplectic hypothesis package, and on the Koszul locus to
which Theorems~\ref{thm:koszul-reflection} and
\ref{thm:higher-genus-inversion} apply, the Lagrangian criterion
\textup{(}item~\textup{(xi)} of the meta-theorem,
Theorem~\textup{\ref{thm:koszul-equivalences-meta})} is an
equivalence: $\cA$ is chirally Koszul if and only if
$\mathcal{M}_\cA$ and $\mathcal{M}_{\cA^!}$ are transverse
Lagrangians in the $(-1)$-shifted symplectic deformation
space~$\mathcal{M}_{\mathrm{comp}}$.

In particular, Proposition~\ref{prop:lagrangian-perfectness} supplies
the perfectness part of the package for algebras in the standard
landscape at the non-critical, non-degenerate levels for which
Theorems~\ref{thm:koszul-reflection} and
\ref{thm:higher-genus-inversion} apply:
Heisenberg at $k \neq 0$, affine Kac--Moody $V_k(\fg)$ at
$k \neq -h^\vee$, Virasoro at generic~$c$, principal
$\mathcal{W}$-algebras at non-critical level, and the free-field
families $bc$, $\beta\gamma$. The C2 Lagrangian conclusion still
requires the remaining cyclic/BV compatibility and bar-chart
acyclicity inputs.
\end{corollary}

\begin{proof}
We verify the shifted-symplectic hypotheses and then prove the two
implications directly on the Koszul surface covered by
Theorems~\ref{thm:koszul-reflection} and
\ref{thm:higher-genus-inversion}.

\smallskip\noindent
\emph{(1) Formal moduli problem with $(-1)$-shifted symplectic
structure.}
The dg~Lie algebra $\fg := \Defcyc^{\mathrm{mod}}(\cA)$ is
complete and filtered (by genus). By the Lurie--Pridham
correspondence \cite{LurieDAGX, Pridham17}, $\fg$ integrates
to a formal moduli problem $\cM_{\mathrm{comp}}(\cA)$ with
tangent complex $\fg[1]$.
Proposition~\ref{prop:lagrangian-perfectness} gives perfectness
of the degree~$-1$ invariant pairing on~$\fg$
under~(P1)--(P3).
By the PTVV/Kontsevich--Pridham principle \cite{PTVV13,Pridham17},
a perfect invariant pairing of degree~$-1$ on a dg~Lie algebra
induces a $(-1)$-shifted symplectic structure on the
corresponding formal moduli problem. Hence
$\cM_{\mathrm{comp}}(\cA)$ carries a canonical $(-1)$-shifted
symplectic structure.

\smallskip\noindent
\emph{(2) Isotropy.}
The one-sided deformation maps $i_\cA$, $i_{\cA^!}$ are
isotropic embeddings: this is
Theorem~\ref{thm:ambient-complementarity-tangent}(ii), proved
unconditionally.

\smallskip\noindent
\emph{(3) Equivalence on the Koszul surface.}
\begin{itemize}
\item \emph{Forward} (Koszul $\Rightarrow$ Lagrangian).
 On the Koszul locus, the twisted tensor product
 $K_\tau(\cA, \cA^!)$ is acyclic. The complementarity
 splitting (Theorem~C) gives a direct-sum decomposition
 $T_{\mathrm{comp}} \simeq T_\cA \oplus T_{\cA^!}$:
 the one-sided tangent complexes are complementary subspaces.
 By~(2) they are isotropic; by complementarity they are
 half-dimensional. Isotropic + half-dimensional = Lagrangian.
\item \emph{Converse} (Lagrangian $\Rightarrow$ Koszul).
 Transverse Lagrangians in a $(-1)$-shifted symplectic
 formal moduli problem have discrete derived
 intersection~\cite{PTVV13}. The derived intersection
 $\cM_\cA \times^h_{\cM_{\mathrm{comp}}} \cM_{\cA^!}$
 computes $K_\tau(\cA, \cA^!)$
 (this is the formal content of the ambient complementarity
 construction: the derived fiber product of the two one-sided
 deformation problems is the twisted tensor product).
 Discreteness (virtual dimension~$0$) gives acyclicity of
 $K_\tau$, which is Koszulness
 (Definition~\ref{def:chiral-koszul-morphism}).
\end{itemize}
This step uses the ordinary twisted-tensor-product interpretation on
the Koszul surface. Off that surface, the file proves only the
intrinsic coderived/contraderived bar-coalgebra comparison of
Theorem~\ref{thm:positselski-chiral-finite-type}.

\smallskip\noindent
\emph{Standard-landscape verification.}
Hypotheses (P1)--(P3) are satisfied by every listed family.
(P\ref{item:perf-fingen}): all standard families have
positive-energy grading with finite-dimensional weight spaces.
(P\ref{item:perf-nondeg}): the invariant form is the Killing
form (nondegenerate at $k \neq -h^\vee$ for Kac--Moody), the
Shapovalov form (nondegenerate at generic~$c$ for Virasoro),
or the canonical pairing for free fields.
(P\ref{item:perf-dual}): Koszul duality preserves the positive-energy
grading and finite-dimensionality of weight spaces; the dual
pairing inherits nondegeneracy from the original.
\end{proof}

\begin{remark}[Scope and degeneration locus]
\label{rem:lagrangian-degeneration-locus}
The Lagrangian criterion degenerates precisely where the invariant
form degenerates: at the critical level $k = -h^\vee$ for
Kac--Moody, and at central charges where the Kac determinant
$\det G_h(c)$ vanishes (admissible levels for Virasoro).
At these points, the cyclic pairing develops a kernel and
perfectness fails, so the formal moduli problem
$\cM_{\mathrm{comp}}$ may not carry a well-defined
$(-1)$-shifted symplectic structure. This degeneration locus
is the same locus where the Shapovalov form degenerates
and null vectors appear: both phenomena reflect the
breakdown of the PBW structure.
\end{remark}

\begin{remark}[Calabi--Yau interpretation of perfectness]
\label{rem:cy-interpretation-perfectness}
\index{Calabi--Yau!Koszul duality exchange}
\index{perfectness!Calabi--Yau interpretation}
Holstein--Rivera~\cite{HR24} prove that Koszul duality exchanges
smooth and proper Calabi--Yau structures in the dg setting.
Hypothesis~(P\ref{item:perf-nondeg}) (nondegenerate invariant form)
gives the chiral algebra~$\cA$ a smooth CY structure; the
Holstein--Rivera exchange then implies that $\barB(\cA)$ carries a
proper CY structure, which is precisely the fiber-level perfectness
used in Step~4 of the proof of
Proposition~\ref{prop:lagrangian-perfectness}.
This does \emph{not} render the Lagrangian criterion~(K11) fully
unconditional: the family-level perfectness over
$\overline{\mathcal{M}}_g$ (the pro-complete limit in Step~4)
still requires~(P\ref{item:perf-fingen}) (finite weight spaces from
PBW) and~(P\ref{item:perf-dual}) (dual regularity via base change).
The CY exchange explains the \emph{mechanism} of perfectness at
each fiber; the \emph{uniformity} across the moduli family is a
separate input.
\end{remark}

\begin{remark}[Synthesis: six faces of chiral Koszulness]
\label{rem:six-faces-koszulness}
Theorems~\ref{thm:barr-beck-lurie-koszulness}--\ref{thm:fm-boundary-acyclicity},
Conjecture~\ref{conj:d-module-purity-koszulness},
and Conjecture~\ref{conj:lagrangian-koszulness}
assemble six comparison surfaces for chiral Koszulness. Let $\cA$ be a
chiral algebra satisfying the standing hypotheses.
The following conditions are related as indicated
\textup{(}cf.\ the meta-theorem, Theorem~\textup{\ref{thm:koszul-equivalences-meta}}\textup{)}:
\begin{enumerate}[label=\textup{(\alph*)}]
\item \textup{(Algebraic)}\quad
 $\cA$ is chirally Koszul:
 $\operatorname{Ext}^i_\cA(k,k)_j = 0$ for $i \neq j$
 \textup{(meta-theorem item~(i))}.
\item \textup{(Descent-theoretic)}\quad
 The Barr--Beck--Lurie comparison functor for
 $\Omega_\kappa \dashv B_\kappa$ is an equivalence
 \textup{(Theorem~\ref{thm:barr-beck-lurie-koszulness};
 meta-theorem item~(vi))}.
\item \textup{(Homological)}\quad
 Factorization homology on a surface satisfying the bar-comparison
 and collapse hypotheses of Theorem~\ref{thm:fh-concentration-koszulness}
 is concentrated in degree~$0$; in genus~$0$ the converse uses the
 PBW filtration and strong convergence
 \textup{(Theorem~\ref{thm:fh-concentration-koszulness};
 meta-theorem item~(vii))}.
\item \textup{(Geometric)}\quad
 FM boundary restrictions $i_S^!\, \barBch_n(\cA)$ are
 acyclic outside degree~$0$
 \textup{(Theorem~\ref{thm:fm-boundary-acyclicity};
 meta-theorem item~(x))}.
\item \textup{(Hodge-theoretic)}\quad
 Bar components $\barBch_n(\cA)$ are Hodge-pure with aligned
 characteristic variety
 \textup{(Conjecture~\ref{conj:d-module-purity-koszulness};
 meta-theorem item~(xii))}.
\item \textup{(Symplectic)}\quad
 $\mathcal{M}_\cA$ and $\mathcal{M}_{\cA^!}$ are transverse
 Lagrangians in $\mathcal{M}_{\mathrm{comp}}$
 \textup{(Conjecture~\ref{conj:lagrangian-koszulness};
 meta-theorem item~(xi);
 conditional on the shifted-symplectic hypothesis package;
 Proposition~\ref{prop:lagrangian-perfectness} supplies the
 perfectness input on the standard landscape)}.
\end{enumerate}
The algebraic, descent-theoretic, and FM-boundary criteria are proved
equivalent on their stated surfaces
\textup{(}Theorems~\ref{thm:barr-beck-lurie-koszulness}
and~\ref{thm:fm-boundary-acyclicity}\textup{)}. The factorization-homology
criterion is the conditional comparison of
Theorem~\ref{thm:fh-concentration-koszulness}; it is not a definition
of Koszulness outside those hypotheses.
The implication (e)$\Rightarrow$(d) is proved; the converse is open.
Item~(f) is conditional on the shifted-symplectic hypothesis package;
Proposition~\ref{prop:lagrangian-perfectness} verifies only its
perfectness input on the standard landscape.
The full equivalence of all six would parallel the classical theorem
that Koszulness, purity, formality, and derived equivalence coincide
for graded algebras arising from geometry~\cite{BGS96}.
\end{remark}

%================================================================
% RECOGNIZING KOSZUL DUALS IN PRACTICE
%================================================================

\section{Recognizing Koszul duals in practice}
\label{sec:recognizing-koszul-duals}

\begin{remark}[Identifying \texorpdfstring{$\mathcal{A}^!$}{A!} in practice]\label{rem:identify-koszul-wild}
Let $\widehat{\mathcal C}$ be a complete conilpotent factorization
coalgebra.  Its cobar algebra
\[
\mathcal A_{\mathrm{cand}}:=\Omega_X^{\mathrm{cont}}
(\widehat{\mathcal C})
\]
comes with the adjunction unit
\[
\eta_{\widehat{\mathcal C}}\colon
\widehat{\mathcal C}\longrightarrow
B_X^{\mathrm{cont}}(\mathcal A_{\mathrm{cand}}).
\]
An equivalence of this map identifies $\widehat{\mathcal C}$ with the
bar coalgebra of $\mathcal A_{\mathrm{cand}}$.

Assume the holonomic duality package for
$\widehat{\mathcal C}$.  Its Verdier dual
\[
\mathbb D_{\operatorname{Ran}}^{\mathrm{cont}}
(\widehat{\mathcal C})
\]
is an algebra whose product is the transpose of the coproduct.  Under
$\eta_{\widehat{\mathcal C}}$, this algebra identifies with
$K_X(\mathcal A_{\mathrm{cand}})$.  A family-specific Koszul-dual
identification is an algebra map
\[
\nu\colon
\mathbb D_{\operatorname{Ran}}^{\mathrm{cont}}
(\widehat{\mathcal C})\longrightarrow\mathcal A^!_{mathrm{cand}}
\]
which is proved to be a quasi-isomorphism in the chosen completion.
For a quadratic presentation, the canonical instance is the Verdier
transpose
\[
\mathbb D(q_{\mathcal A})\colon
K_X(\mathcal A)\longrightarrow
\mathbb D\mathcal A^{\mathrm i}.
\]
The derived centre is the independently defined Hochschild algebra
$C^\bullet_{\mathrm{ch}}(\mathcal A,\mathcal A)$.
\end{remark}

\section{The Ran-space perspective}
\label{sec:ran-space-perspective}
\index{Ran space!bar-cobar perspective}

The $\infty$-categorical framework of factorization
algebras on $\operatorname{Ran}(X)$ provides the conceptual home for
the bar-cobar adjunction.

\begin{remark}[From configuration spaces to Ran space]\label{rem:config-to-ran}
The configuration spaces $C_n(X)$ and their FM compactifications
$\overline{C}_n(X)$ are finite-dimensional approximations to the
\emph{Ran space} $\operatorname{Ran}(X) = \operatorname{colim}_n\,
X^n / \Sigma_n$, the space of all finite non-empty subsets of~$X$.
The bar complex $\bar{B}(\cA) = \bigoplus_n \bar{B}^n(\cA)$, with
$\bar{B}^n(\cA)$ living on $\overline{C}_n(X)$, assembles into a
single object on $\operatorname{Ran}(X)$. The bar differential,
given by residues at collision divisors
(Definition~\ref{def:bar-differential-complete}), becomes the
factorization structure map: when two points collide, the OPE
produces the ``fusion'' of the corresponding tensor factors.

In this language, the Arnold relations that ensure $d^2 = 0$ on the
ordered genus-$0$ square-zero bar complex
(Theorem~\ref{thm:bar-nilpotency-complete}) are the coherence conditions
for the factorization algebra axiom: the result of a three-point
collision does not depend on the order in which pairs collide,
because the relevant boundary strata of $\overline{C}_3(X)$ are
related by the Arnold relation
$\eta_{12} \wedge \eta_{23} + \eta_{23} \wedge \eta_{31}
+ \eta_{31} \wedge \eta_{12} = 0$.
\end{remark}

\begin{remark}[Verdier duality and the Koszul-dual component]\label{rem:verdier-engine}
\index{Verdier duality!Koszul duality engine}
The twisting adjunction is
\[
\Omega_X \colon
\operatorname{CoAlg}_{\mathrm{conil}}(\operatorname{Fact}(X))
\;\rightleftarrows\;
\operatorname{Alg}_{\mathrm{aug}}(\operatorname{Fact}(X))
 :\!\bar{B}_X
\]
with counit $\Omega_X\bar B_X(\cA)\to\cA$.  On the holonomic
dualizable domain, continuous Verdier duality sends the bar coalgebra
to the algebra
\[
K_X(\cA):=
\mathbb{D}_{\operatorname{Ran}}^{\mathrm{cont}}
\bar{B}_X^{\mathrm{cont}}(\cA).
\]
Its multiplication is
$\mathbb D(\Delta_B)\circ\chi^{-1}$, where $\chi$ is the Verdier
$!$-to-$*$ tensor exchange.  Thus $K_X$ has variance
\[
K_X\colon
(\operatorname{Alg}^{\mathrm{aug}}(\operatorname{Fact}(X)))^{\mathrm{op}}
\longrightarrow
\operatorname{Alg}^{*,\mathrm{aug}}(\operatorname{Fact}(X)).
\]

The natural Verdier--bar comparison is
\[
\gamma_\cA\colon
\Omega_X^{\mathrm{geom,cont}}(\mathbb D_X\cA)
\xrightarrow{\sim}K_X(\cA).
\]
Here $\mathbb D_X\cA$ is a coalgebra, so the displayed cobar has the
required source type.  The map is induced by biduality
$\cA\to\mathbb D_X\mathbb D_X\cA$ as in
Theorem~\ref{thm:bar-cobar-verdier}.

A chiral Koszul pair supplies the further algebra map
\[
\nu_{12}\colon K_X(\cA_1)\longrightarrow\cA_2
\]
of Theorem~\ref{thm:bar-cobar-isomorphism-main}.  A comparison with a
bar coalgebra $B_X(\cA^!)$ would additionally require a monoidal
exchange functor from post-Verdier $*$-algebras to pre-Verdier
$!$-coalgebras and a natural morphism from the image of~$K_X(\cA)$ to
$B_X(\cA^!)$.  The present theorem uses the typed maps~$\gamma_\cA$
and~$\nu_{12}$.
\end{remark}

\begin{remark}[Curvature and coderived categories]\label{rem:curvature-coderived}
\index{coderived category!bar-cobar}
\index{curvature!coderived category}
When curvature is present, i.e., when the genus-$1$ obstruction
$m_0^{(1)} = \kappa(\cA) \cdot \lambda_1 \neq 0$ (as it is for all
Koszul chiral algebras with $\kappa(\cA) \neq 0$; see
Theorem~\ref{thm:genus-universality}), the bar complex satisfies
$\dfib^{\,2} = \kappa(\cA) \cdot \omega_1 \cdot \mathrm{id}$, not $\dfib^{\,2} = 0$
(see Convention~\ref{conv:higher-genus-differentials} for the notation).
The ordinary derived category of $\cA$-modules is then the wrong
quotient: it erases the distinction between ``acyclic because of
$\dfib^{\,2} \neq 0$'' and ``acyclic because of genuine exactness.''

The correct ambient category is the \emph{coderived category}
in the sense of Positselski~\cite{Positselski11}. In this category,
a curved complex $(M, d_M)$ with $d_M^2 = \kappa \cdot \mathrm{id}_M$
is \emph{not} set to zero; it represents a genuine object encoding the
curvature. On the intrinsic bar-coalgebra surface, the
Positselski comparison yields an equivalence
\[
\mathcal{D}^{\mathrm{co}}(\barB_X(\cA)\text{-}\mathrm{comod})
\;\simeq\;
\mathcal{D}^{\mathrm{ctr}}(\barB_X(\cA)\text{-}\mathrm{contra}),
\]
and, in finite type, the contraderived side may be rewritten as
complete modules over the graded dual bar-dual algebra
$\barB_X(\cA)^*$. On the flat finite-type locus where the internal
Positselski comparison applies, this exotic comparison can be reduced
to ordinary derived categories on the completed-dual side; on the
quadratic genus-$0$ bar-dual surface it further reduces to
Theorem~\ref{thm:e1-module-koszul-duality}. Off these loci, the
coderived formalism captures the ``acyclic-but-not-contractible''
phenomena that higher-genus curvature creates.

This perspective clarifies the physical meaning of curvature: the
cosmological constant $\Lambda \sim \kappa(\cA)$ in the AdS$_3$/CFT$_2$
correspondence (Conjecture~\ref{conj:ads-cft-bar}) is the obstruction
to reducing the coderived category to the derived category. Anomaly
cancellation at $c = 26$ ($\kappa_{\mathrm{eff}} = 0$; note $\kappa(\mathrm{Vir}_{26}) = 13$) removes this scalar
curvature obstruction on the effective scalar surface, but by
itself it is not the full criterion under which
coderived and derived categories agree, nor does the manuscript
identify the full bosonic-string theory as ``on-shell'' in a
derived-categorical sense.
\end{remark}

At genus zero the bar differential is the signed residue differential
on Fulton--MacPherson compactifications, and the Arnold relation gives
$d^2=0$.  The enhanced Ran theorem supplies bar--cobar reconstruction;
the geometric realization uses its factorization and completion
packages.  Verdier duality supplies the algebra
\[
K_X(\cA)=\mathbb D_{\operatorname{Ran}}^{\mathrm{cont}}B_X(\cA)
\]
and the natural equivalence
$\Omega_X^{\mathrm{geom,cont}}(\mathbb D_X\cA)\simeq K_X(\cA)$.
For a chiral Koszul pair, the map~$\nu_{12}$ identifies this algebra
with the opposite member.  These constructions live on the fixed
curve~$X$; the relative theory adds base change and sewing.

\subsection{Bar-cobar as factorization homology}
\label{subsec:bar-cobar-fh}

The Ran-space remarks above assemble the bar complex into an
object on $\operatorname{Ran}(X)$. Under the locally constant,
constructible, and compactness hypotheses used below, this assembly
gives a comparison model for factorization homology; outside those
hypotheses it is not a definition-free identification.

\begin{proposition}[Bar construction as factorization homology;
\ClaimStatusConditional]
\label{prop:bar-fh}
\index{factorization homology!bar construction}
\index{bar construction!as factorization homology}
Let $\cA$ be an augmented $\Eone$-chiral algebra on a smooth
curve~$X$, and let $\overline{\cA} = \ker(\epsilon \colon \cA \to k)$
be the augmentation ideal. Assume the associated Ran object is
locally constant on configuration strata, constructible on
Fulton--MacPherson compactifications, and compact enough for the
Ran colimit to commute with the bar-length filtration. Then there is
a natural comparison equivalence
\[
\Bbar^{\mathrm{geom}}_X(\cA)
\;\simeq\;
\int_X \overline{\cA}^{\mathrm{Lie}}
\]
where $\overline{\cA}^{\mathrm{Lie}}$ denotes $\overline{\cA}$ viewed as
an $\Eone$-Lie algebra via the commutator bracket, and the right-hand
side is factorization homology in the sense of
\textup{\cite{AF15, BD04}}. More precisely:
\begin{enumerate}[label=(\roman*)]
\item \textbf{Finite approximations}: For each $n \geq 1$,
the degree-$n$ piece $\Bbar^n_X(\cA) =
(s^{-1}\overline{\cA})^{\otimes n}_{\overline{C}_n(X)}$
(the bar complex on~$n$ points, supported on the
Fulton--MacPherson compactification $\overline{C}_n(X)$) is the
restriction of the factorization homology $\int_X\overline{\cA}$
to $n$-point configurations.
\item \textbf{Assembly}: The bar differential
$d_{\mathrm{bar}} = \sum_{i < j} d_{ij}$ (residue along collision
divisors, Definition~\ref{def:bar-differential-complete}) is the
structure map of the factorization algebra:
the collision $z_i \to z_j$ implements the factorization product
$\int_{U_i}\overline{\cA} \otimes
\int_{U_j}\overline{\cA} \to
\int_{U_i \cup U_j}\overline{\cA}$
for shrinking discs $U_i \ni z_i$, $U_j \ni z_j$.
\item \textbf{Colimit}: The total bar complex
$\Bbar_X(\cA) = \bigoplus_{n \geq 0}\Bbar^n_X(\cA)$
is the colimit over $n$ of the finite approximations,
i.e., the factorization homology
$\int_X\overline{\cA} = \operatorname{colim}_{S \in
\mathrm{Fin}_{/X}} \bigotimes_{s \in S}\overline{\cA}_s$.
\end{enumerate}
\end{proposition}

\begin{proof}
The factorization homology of an $E_n$-algebra $A$ on a manifold
$M$ is defined as
\[
\textstyle\int_M A \;=\; \operatorname{colim}_{U_1 \sqcup \cdots
\sqcup U_k \hookrightarrow M} A(U_1) \otimes \cdots
\otimes A(U_k)\,,
\]
where the colimit is over rectilinear embeddings of disjoint discs
\cite[Definition~3.1]{AF15}. For $n = 1$ on a curve~$X$, this is
the colimit over configurations of disjoint discs
$D_1 \sqcup \cdots \sqcup D_k \hookrightarrow X$, with
$A(D_i) = \overline{\cA}_{z_i}$ (the fiber of $\overline{\cA}$
at the center $z_i$ of $D_i$).

Part (i): The term $\Bbar^n_X(\cA)$ places one copy of
$s^{-1}\overline{\cA}$ at each of $n$ distinct points
$z_1, \ldots, z_n \in X$. The locus of distinct points is the
open configuration space $C_n(X) = X^n \smallsetminus \Delta$,
and the extension to $\overline{C}_n(X)$ encodes the factorization
product at coincident points.

Part (ii): When $z_i \to z_j$, the factorization structure map
composes the two copies of $\overline{\cA}$ using the chiral
product. At the chain level, this composition is the OPE residue
$\mathrm{Res}_{z_i = z_j}(\mu^{\mathrm{ch}}(\bar{a}_i, \bar{a}_j)
\cdot \eta_{ij})$, which is exactly the bar differential $d_{ij}$.

Part (iii): The colimit over $n$ in the bar complex is the colimit
over the category $\mathrm{Fin}_{/X}$ of finite subsets of $X$,
which is the Ran space $\operatorname{Ran}(X) =
\operatorname{colim}_n X^n/\Sigma_n$. The factorization homology
$\int_X\overline{\cA}$ is this colimit by definition.
\end{proof}

\begin{proposition}[Cobar as factorization cohomology;
\ClaimStatusConditional]
\label{prop:cobar-fh}
\index{factorization homology!cobar construction}
\index{cobar construction!as factorization cohomology}
Assume $\cC$ is a conilpotent $\Eone$-chiral coalgebra whose
completed cobar filtration is Mittag--Leffler and whose Verdier dual
is constructible on the Ran strata. Then the cobar construction
$\Omega_X(\cC)$ is naturally equivalent, in this completed model, to
the
\emph{factorization cohomology}:
\[
\Omega_X(\cC) \;\simeq\;
\mathbb{R}\mathrm{Hom}_{\operatorname{Fact}(X)}
\bigl(k,\, \cC\bigr)
\;=:\;
\int^X \cC\,,
\]
the derived internal hom from the unit factorization algebra $k$
to~$\cC$ in the $\infty$-category of factorization coalgebras on~$X$.
Under Verdier duality $\mathbb{D}_{\mathrm{Ran}}$, this becomes:
\[
\Omega_X(\cC)
\;\simeq\;
\Bigl(\int_X \mathbb{D}_{\mathrm{Ran}}(\cC)\Bigr)^{\!\vee}\,,
\]
the linear dual of the factorization homology of the Verdier dual.
For $\cC = \Bbar_X(\cA)$, this gives
\[
\Omega_X\Bbar_X(\cA) \;\simeq\;
\Bigl(\int_X K_X(\cA)\Bigr)^{\!\vee}\,,
\]
where $K_X(\cA)=\mathbb D_{\Ran}\Bbar_X(\cA)$.  If a chiral Koszul
pair supplies $\nu\colon K_X(\cA)\xrightarrow{\sim}\cA^!$, the right
side becomes $(\int_X\cA^!)^\vee$.  The bar--cobar counit is the map
$\psi\colon\Omega_X\Bbar_X(\cA)\to\cA$.
\end{proposition}

\begin{proof}
The cobar construction $\Omega_X(\cC) = T(s^{-1}\,\overline{\cC})$ with
differential $d_{\mathrm{cobar}} = d_{\mathrm{coop}} +
d_{\mathrm{tw}}$ is the free algebra on $s^{-1}\,\overline{\cC}$, with
the differential encoding the coalgebra structure of~$\cC$.
The free algebra on a module $V$ is
$\mathbb{R}\mathrm{Hom}(k, T^c(V))$, and for
$V = s^{-1}\,\overline{\cC}$ supported on $\operatorname{Ran}(X)$, this
is the factorization cohomology $\int^X\cC$ by
\cite[Proposition~5.5.2.5]{HA}.

The Verdier duality identification follows from the definition
$K_X(\cA)=\mathbb D_{\mathrm{Ran}}\Bbar_X(\cA)$ and the general
adjunction
$\mathrm{Hom}(k, \cC) \simeq (\int_X \mathbb{D}\cC)^\vee$.
The final replacement uses functoriality of factorization homology
under the supplied equivalence~$\nu$.
\end{proof}

\begin{remark}[The bar-cobar adjunction as a duality of integration
theories]
\label{rem:bar-cobar-integration}
\index{factorization homology!duality of integration theories}
Under the compactness and constructibility hypotheses of
Propositions~\ref{prop:bar-fh} and~\ref{prop:cobar-fh}, the
bar and cobar functors admit factorization-homology realizations:
\[
\Omega_X \colon
\operatorname{CoAlg}_{\mathrm{conil}}(\operatorname{Fact}(X))
\;\rightleftarrows\;
\operatorname{Alg}_{\mathrm{aug}}(\operatorname{Fact}(X))
:\!\Bbar_X
\]
with the following typed outputs.
\begin{itemize}
\item $B_X(A)$ is assembled from configuration-space collision
operations and carries the conilpotent coalgebra structure.
\item $\Omega_X(C)$ is the free augmented algebra on the desuspension
of the coaugmentation coideal, with differential induced by the
coproduct and twisting datum.
\item The universal counit $\Omega_XB_X(A)\to A$ is the
reconstruction morphism of Theorem~A.
\item For $A=T_X(V)/(R)$, the morphism
$q_A\colon A^{\mathrm i}=C_X(s^{-1}V,s^{-2}R)\to B_X(A)$ is the quadratic
comparison of Theorem~B.
\item Under $H_{\mathbb D}^{\mathrm{bar}}(A)$, Verdier duality forms the algebra
$K_X(A)=\mathbb D_{\operatorname{Ran}}B_X(A)$.
\end{itemize}
The Fourier language of Remark~\ref{rem:koszul-fourier} is an
expository analogy for these five constructions.  In that analogy,
universal reconstruction corresponds to resolution, while quadratic
Koszulness corresponds to compression of the full bar coalgebra to
$A^{\mathrm i}$.
\end{remark}

\begin{remark}[The master diagram at genus~\texorpdfstring{$0$}{0}]
\label{rem:master-diagram-genus0}
\index{master diagram!genus zero}
At genus~$0$ the chapter uses three composable comparisons:
\[
 A^{\mathrm i}\xrightarrow{q_A}B_X(A),
 \qquad
 \Omega_XB_X(A)\xrightarrow{\varepsilon_A}A,
 \qquad
 K_X(A):=\mathbb D_{\operatorname{Ran}}B_X(A).
\]
The middle morphism is the universal reconstruction equivalence in
the enhanced associative Ran ambient.  The first morphism is an
equivalence precisely on the quadratic Koszul surface described by
$H_{\mathrm{CL}}$.  Verdier duality transports the coalgebra structure
of the full bar object to the algebra $K_X(A)$ under
$H_{\mathbb D}^{\mathrm{bar}}(A)$.  Applying Verdier duality to $q_A$ identifies
$K_X(A)$ with the continuous dual of $A^{\mathrm i}$ when the graded
pieces are dualizable.  These maps give the typed form of the master
diagram~\eqref{eq:master-square}.

A positive-genus modular enhancement supplies a curved coalgebra
$C_g=(C,d,h_g)$.  Positselski's fixed-coalgebra theorem then compares
$D^{\mathrm{co}}(C_g\text{-}\mathrm{CoFact})$ with
$D^{\mathrm{ctr}}(C_g\text{-}\mathrm{ContraFact})$ under its
second-kind package.  This modular enhancement is an additional construction
built from cyclic pairing and clutching data.
\end{remark}

\section{The categorical logarithm}
\label{sec:categorical-logarithm}
\index{categorical logarithm|textbf}
\index{nilpotence-periodicity correspondence!categorical logarithm}

The phrase \emph{categorical logarithm} is used here as notation for
the bar functor together with its bar-length filtration.  Its precise
mathematical content consists of the enhanced bar--cobar equivalence,
the quadratic comparison $q_A$, the Verdier algebra $K_X(A)$, and the
separately constructed modular extension.  The analogy organizes
these maps; each theorem below is stated in its native category.

\subsection{The exponential-logarithm correspondence}
\label{subsec:exp-log-correspondence}

\begin{definition}[Categorical logarithm and exponential;
\ClaimStatusDefinitional]
\label{def:categorical-log-exp}
\index{categorical logarithm!definition}
\index{categorical exponential!definition}
Let $\cA$ be an augmented $\Eone$-chiral algebra on a smooth
curve~$X$. The \emph{categorical logarithm} of~$\cA$ is the bar
construction
\[
\log^{\mathrm{cat}}(\cA) \;:=\; \barB_X(\cA)\,,
\]
viewed as a factorization coalgebra on $\operatorname{Ran}(X)$
(Proposition~\ref{prop:bar-fh}). The \emph{categorical exponential}
is the cobar construction
\[
\exp^{\mathrm{cat}}(\cC) \;:=\; \Omega_X(\cC)\,,
\]
viewed as a factorization algebra on $\operatorname{Ran}(X)$
(Proposition~\ref{prop:cobar-fh}).
\end{definition}

The following table separates the analogy from the maps it records.

\begin{center}
\small
\renewcommand{\arraystretch}{1.3}
\begin{tabular}{p{5.5cm}p{5.5cm}}
\textbf{Logarithmic analogy} & \textbf{Mathematical construction} \\
\hline
formal logarithm &
$B_X\colon\operatorname{Alg}^{\mathrm{aug}}\to
\operatorname{CoAlg}^{\mathrm{conil}}$ \\[2pt]
formal exponential &
$\Omega_X\colon\operatorname{CoAlg}^{\mathrm{conil}}\to
\operatorname{Alg}^{\mathrm{aug}}$ \\[2pt]
recovery of the input &
$\Omega_XB_X(A)\xrightarrow{\sim}A$ in the enhanced Ran ambient
(Theorem~A) \\[2pt]
Taylor filtration &
the bar-length/PBW filtration on $B_X(A)$ \\[2pt]
quadratic approximation &
$q_A\colon A^{\mathrm i}\to B_X(A)$
(Theorem~B) \\[2pt]
dual transform &
$K_X(A)=\mathbb D_{\operatorname{Ran}}B_X(A)$ under
$H_{\mathbb D}^{\mathrm{bar}}(A)$ \\[2pt]
curved extension &
a fixed chiral CDG-coalgebra $C_g$ and its second-kind module
categories under $H_{\mathrm{mod}}$
\end{tabular}
\end{center}

\begin{remark}[The logarithmic kernel]
\label{rem:logarithmic-kernel}
\index{logarithmic kernel}
On a coordinate disc the form
$\eta_{ij}=d\log(z_i-z_j)$ represents the logarithmic de Rham class
of the collision divisor.  Its residue pairs with the chiral
multiplication to define the binary collision term in the local bar
differential.  Arnold's three-point relation controls the geometric
codimension-two contribution, while associativity and the
Jacobi--Borcherds identity control the algebraic contribution.  This
local calculation is the precise content behind the logarithmic
terminology.
\end{remark}

\subsection{The Taylor expansion of the logarithm}
\label{subsec:taylor-expansion-log}
\index{PBW spectral sequence!as Taylor expansion}

The weight filtration on the bar complex
$\barB(\cA) = \bigoplus_{n \geq 0} (s^{-1}\overline{\cA})^{\otimes n}$
induces a spectral sequence
\[
E_1^{p,q} \;=\; H^q\bigl(\mathrm{gr}^p\, \barB(\cA)\bigr)
\;\Longrightarrow\;
H^{p+q}(\barB(\cA))\,.
\]
The filtration degree records tensor length and hence the number of
inputs before collision contractions.  The differential on the first
page is induced by the associated-graded binary operation.  Higher
differentials measure the failure of cycles to lift through successive
filtration windows.  Calling this spectral sequence a Taylor expansion
is an expository analogy for this exact filtered construction.

\begin{proposition}[PBW detection of the quadratic comparison;
\ClaimStatusConditional]
\label{prop:log-convergence}
\index{convergence!of categorical logarithm}
Under the split finite-window PBW hypotheses of
Proposition~\ref{prop:koszul-locus-criteria-package-bci}, assume
$H_{\mathrm{CL}}(A,A^{\mathrm i},\tau_{\mathrm i})$.  The following
conditions are equivalent:
\begin{enumerate}[label=(\roman*)]
\item $q_A\colon A^{\mathrm i}\to B_X(A)$ is a quasi-isomorphism;
\item $\Omega_X(A^{\mathrm i})\to A$ is a quasi-isomorphism;
\item the left and right twisted tensor products for
$\tau_{\mathrm i}$ are acyclic;
\item the PBW/bar spectral sequence has quadratic diagonal support and
strong convergence.
\end{enumerate}
\end{proposition}

\begin{proof}
This is the chiral Comparison Lemma package of
Theorem~\ref{thm:bar-cobar-inversion-qi}, transported through the
split finite windows and the strongly convergent PBW filtration.
\end{proof}

\begin{remark}[Analytic growth problem; \ClaimStatusHeuristic]
\label{rem:convergence-radius-ope}
\index{convergence!radius from OPE poles}
\index{Koszul locus!characterization via OPE}
An analytic refinement begins by choosing a complete locally convex
topology on every bar-weight space and continuous norms for the
collision residue operators.  A bound
\[
 \lVert R_n\rVert\leq CR^n
\]
would control convergence of the completed collision sum in that
topology.  A theorem relating this analytic convergence to quadratic
Koszul recognition would additionally identify the resulting
completion with the filtered comparison
$q_A\colon A^{\mathrm i}\to B_X(A)$ and prove strong convergence of
its comparison cone.  These are concrete family-level obligations.
\end{remark}

\subsection{Additivity and the branch structure}
\label{subsec:additivity-branch}
\index{categorical logarithm!additivity}
\index{complementarity!as additivity of logarithm}

The word \emph{additivity} refers here to a supplied direct-sum
decomposition of the genus complex.  Its construction belongs to the
complementarity package.

\begin{proposition}[Complementarity in logarithmic notation;
\ClaimStatusConditional]
\label{prop:log-additivity}
Let $(\cA,\cA^!)$ be a Verdier--Koszul pair satisfying the hypothesis
package of Theorem~\ref{thm:quantum-complementarity-main}.  Assume its
comparison map supplies an equivalence
\[
 \alpha_g\colon Q_g(\cA)\oplus Q_g(\cA^!)
 \xrightarrow{\ \sim\ }
 H^*(\overline{\mathcal{M}}_g,Z_\cA).
\]
If the same package supplies a perfect anti-invariant pairing for
which the two summands are isotropic, then $\alpha_g$ is a Lagrangian
decomposition.  In the terminology of this section, the two summands
are the two additive components of the logarithmic datum.  A formula
for $B_X(\cA\otimes\cA^!)$ is governed by its own monoidal K"unneth
comparison.
\end{proposition}

\begin{remark}[Filtered extensions and the BCH analogy;
\ClaimStatusConditional]
\label{rem:categorical-bch}
\index{Baker--Campbell--Hausdorff formula!categorical analogue}
Let
$0\to\cA'\to\cA\to\cA''\to0$ be an extension represented by a split
complete filtration in the chosen chiral category.  Suppose the
extension filtration is multiplicative, the induced bar filtration is
complete and strongly convergent, and a twisting cochain~$\tau$
realizes the extension data.  The filtered comparison problem is then
\[
 B_X(\cA)\longrightarrow
 B_X(\cA'')\widehat\otimes_\tau B_X(\cA').
\]
A proof that this map is a quasi-isomorphism gives the categorical
BCH statement for that extension.  The Maurer--Cartan recursion of
\S\ref{sec:maurer-cartan-curved} computes the successive twisting
components.  The classical BCH series is the motivating analogy.
\end{remark}

\subsection{From single-valued to multi-valued: the genus transition}
\label{subsec:genus-transition-log}
\index{categorical logarithm!genus transition}
\index{monodromy!of categorical logarithm}

Positive genus introduces cyclic contractions and stable-graph
sewing.  The transition is expressed by the following typed data.

\begin{center}
\small
\renewcommand{\arraystretch}{1.3}
\begin{tabular}{p{3.5cm}p{4.0cm}p{4.0cm}}
\textbf{Surface} & \textbf{Input} & \textbf{Output} \\
\hline
genus~$0$ & local collision package & square-zero bar differential
$d_0^2=0$ \\[2pt]
stable genus~$g$ & perfect cyclic pairing, signed clutching, complete
graph filtration & modular Maurer--Cartan element
$\Theta_{\cA,g}$ \\[2pt]
curved coalgebra & fixed $C_g=(C,d,h_g)$ and second-kind exactness &
$D^{\mathrm{co}}(C_g\text{-}\mathrm{CoFact})\simeq
D^{\mathrm{ctr}}(C_g\text{-}\mathrm{ContraFact})$ \\[2pt]
scalar projection & a specified trace map & cohomology classes on
$\overline{\mathcal M}_{g,n}$
\end{tabular}
\end{center}

\begin{remark}[Riemann--Hilbert comparison problem;
\ClaimStatusConditional]
\label{rem:chiral-rh-revisited}
\index{Riemann--Hilbert correspondence!chiral}
Suppose the modular bar construction forms a regular holonomic flat
family over a smooth base~$S$, with a specified connection
$\nabla_{\mathrm{bar}}$ and an identification of its de Rham complex
with the corrected total complex.  Riemann--Hilbert then associates a
local system and a monodromy representation to that family.  A
comparison with the scalar characteristic requires three maps:
\begin{enumerate}[label=(\roman*)]
\item a chain map from the corrected total differential to the de
Rham differential of $\nabla_{\mathrm{bar}}$;
\item a Chern--Weil or residue map from the fiber-curvature operator to
the chosen scalar class;
\item a theorem relating that scalar class to the logarithm of the
monodromy on a specified finite-rank subquotient.
\end{enumerate}
The differentials $d_0$, $d_{\mathrm{fib}}$, and $D_g$ of
Convention~\ref{conv:higher-genus-differentials} occupy the local,
fiberwise curved, and corrected total complexes respectively.
\end{remark}

\subsection{The characteristic hierarchy as logarithmic data}
\label{subsec:characteristic-log}
\index{modular characteristic!as logarithmic data}

The symbols $(\kappa,\Delta,\Theta,\Pi)$ denote four typed outputs.
Their relationships are comparison theorems rather than consequences
of notation.

\begin{enumerate}[label=\textbf{Level \arabic*}:, leftmargin=3.5em]
\item \textbf{Scalar datum.}  The modular characteristic
$\kappa(\cA)$ is obtained from the specified scalar trace of the
collision or curvature class under the theorem's stated hypotheses.

\item \textbf{Spectral datum.}  A finite-rank modular representation
has a characteristic polynomial or discriminant~$\Delta(\cA)$ and a
periodicity profile~$\Pi(\cA)$ when its eigenvalues have finite order.

\item \textbf{Chain datum.}  The modular Maurer--Cartan element
$\Theta_\cA$ is a complete stable-graph sum in the modular convolution
algebra under $H_{\mathrm{mod}}$.  A trace maps its components to
cohomology classes.
\end{enumerate}

\noindent
A hierarchy connecting these outputs consists of explicit maps from
the chain datum to the spectral representation and from the spectral
representation to the scalar trace.  Each family chapter must
construct those maps and verify their compatibility.

\subsection{Visible theorems}
\label{subsec:visible-theorems-log}
\index{categorical logarithm!predicted theorems}

The logarithmic analogy organizes five precise proof obligations:

\begin{enumerate}[label=(\arabic*)]
\item \textbf{Categorical Torelli; \ClaimStatusConjectured.}
Construct a monodromy functor on a specified category of modular bar
families and prove that it is conservative after adjoining the
quadratic comparison $q_A$ and the augmentation.  This is a new
reconstruction theorem.

\item \textbf{Verdier functional equation;
\ClaimStatusConditional.} Construct the comparison map of
Proposition~\ref{prop:log-additivity} and verify its perfect pairing.
The phrase \emph{functional equation} refers to that displayed
equivalence.

\item \textbf{Unit calculation; \ClaimStatusProvedHere.}  For the
augmented unit algebra $k$, the reduced augmentation ideal is zero and
$B(k)=k$ in bar degree~$0$.  Hence
$H^*(B(k))=k$ in degree~$0$.

\item \textbf{Analytic growth; \ClaimStatusHeuristic.}  Choose a
complete topology and prove residue-operator bounds of the form stated
in Remark~\ref{rem:convergence-radius-ope}; then compare the completed
sum with $\operatorname{Cone}(q_A)$.

\item \textbf{Screening comparison; \ClaimStatusConjectured.}  In a
lattice realization, construct a functor carrying the screening
complex to the chosen small-quantum-group module category and identify
the nilpotence relation on each side.  A monodromy interpretation then
uses the regular flat family of
Remark~\ref{rem:chiral-rh-revisited}.
\end{enumerate}

\begin{remark}[Three cofree coalgebra carriers;
\ClaimStatusDefinitional]
\label{rem:quantized-symmetric-power}
\index{bar complex!as quantized symmetric power}
\index{quantized symmetric power}
For a complex $W=s^{-1}\bar\cA$, three operadic choices produce three
cofree conilpotent coalgebras:
\begin{itemize}
\item $T^c(W)$ is the coassociative deconcatenation coalgebra used by
the ordered associative bar construction.
\item $\operatorname{Sym}^c(W)$ is the cocommutative coalgebra used by
the commutative bar construction.  When graded Euler characteristics
are finite, its symmetric-power series is the corresponding
plethystic exponential.
\item $\operatorname{coLie}(W)$ is the cofree Lie coalgebra used by
the Harrison or Lie-operadic bar construction.
\end{itemize}
Symmetrization, ordered descent, and PBW comparison are morphisms
between specified operadic models; each requires its equivariance and
completion hypotheses.  The bar-length spectral sequence belongs to
$T^c(W)$.  The stable-graph obstruction filtration belongs to the
modular convolution algebra.  A comparison between these filtrations
is a separate filtered chain map.
\end{remark}

\begin{conjecture}[Factorization finiteness criterion;
\ClaimStatusConjectured]\label{conj:factorization-finiteness-criterion}
\index{factorization bar-cobar!finiteness criterion}
\index{thick generation!bypassed by sectorwise finiteness}
Let $\cA=T_X(V)/(R)$ be a complete quadratically presented chiral
algebra, and suppose
$q_\cA\colon\cA^{\mathrm i}\to B_X(\cA)$ carries a complete
bigraded filtration with finite-dimensional pieces.  Assume:
\begin{enumerate}[label=\textup{(\roman*)}]
\item the spectral sequence of
$\operatorname{Cone}(q_\cA)$ converges strongly and has zero
$E_\infty$-page in every bidegree;
\item the finite-window dimensions satisfy a stated summability or
subexponential bound in the selected completion topology;
\item the collision restriction maps, disjoint-union products, and
inverse-limit transitions are uniformly compatible with this
filtration.
\end{enumerate}
Then the sectorwise quasi-isomorphism $q_\cA$ should assemble to a
factorization quasi-isomorphism on $\operatorname{Ran}(X)$, and the
finite-window generators should compactly generate its factorization
closure.  The open obligation is the passage from uniform sectorwise
control to compact generation across collision strata.

\emph{Candidate filtrations.}
Lattice charge, Yangian root/loop degree, and Drinfeld--Sokolov
filtrations provide concrete gradings on which the three clauses can
be tested.  Each application requires a direct calculation of the
comparison cone and its collision base-change maps.
\end{conjecture}

\begin{proposition}[Tensor-bar counterexample to automatic
subexponential growth; \ClaimStatusProvedHere]%
\label{prop:subexponential-growth-automatic}%
\index{sub-exponential growth!automatic for positive-energy}%
Let $A=k\oplus V$ be the augmented square-zero algebra with
$V$ concentrated in conformal weight~$1$ and
$d=\dim_kV\geq2$.  Its character has finite support, while the
weight-$p$ component of the reduced bar coalgebra has dimension
\begin{equation}\label{eq:bar-dim-bound}
 \dim_k B(A)_p=d^p.
\end{equation}
Consequently a Cardy--Hardy--Ramanujan upper bound on the input
character can coexist with exponential tensor-bar growth.
\end{proposition}

\begin{proof}
The product on $V$ is zero, so the multiplication part of the bar
differential vanishes.  Hence
\[
 B(A)=T^c(s^{-1}V)
 =\bigoplus_{p\geq0}(s^{-1}V)^{\otimes p}.
\]
Every length-$p$ word has conformal weight~$p$, and the corresponding
space has dimension $(\dim V)^p=d^p$.  Taking any
$1<c<d$ contradicts the proposed subexponential bound.
\end{proof}

\begin{corollary}[Independent growth and geometry obligations;
\ClaimStatusProvedHere]\label{cor:finiteness-criterion-reduction}
\index{factorization bar-cobar!finiteness criterion!reduction}
An application of
Conjecture~\textup{\ref{conj:factorization-finiteness-criterion}}
requires both a direct growth estimate for the relevant comparison
cone and a geometric proof of factorization descent and compact
generation.  Bounds for the input character and sectorwise spectral
sequence degeneration supply distinct pieces of this package.
\end{corollary}

\subsection{The categorical logarithm and its four properties}
\label{subsec:monograph-as-log}
\index{categorical logarithm!four properties}

The logarithmic terminology summarizes four mathematically distinct
constructions:
\begin{itemize}
\item \textbf{Theorem~A: universal reconstruction.} Enhanced bar and
cobar are inverse in the pro-nilpotent associative Ran ambient;
factorization and completed chain realizations use their named
packages.
\item \textbf{Theorem~B: quadratic recognition.} The map
$q_A\colon A^{\mathrm i}\to B_X(A)$ is a quasi-isomorphism exactly
when the equivalent quadratic twisting criteria hold under
$H_{\mathrm{CL}}$.
\item \textbf{Theorem~C: centre-local-system complementarity.} C0
supplies~$\mathcal Z(A)$ from the strict-flat degree-zero fibre; C1
splits
$R\Gamma(\overline{\mathcal M}_g,\mathcal Z(A))$ by the represented
Verdier involution; C2 supplies the paired shifted-symplectic upgrade
on its stated hypothesis package.
\item \textbf{Theorem~D: modular trace.} Cyclic pairing and
stable-graph sewing produce a modular Maurer--Cartan element; a
specified trace produces the scalar class under the theorem's stated
hypotheses.
\end{itemize}

The associated research programme has five separately testable
questions:
\begin{itemize}
\item \textbf{MC1}: compute the PBW/bar comparison cone.
\item \textbf{MC3}: construct the categorical
Clebsch--Gordan closure on its evaluation-generated core and specify
the extension domain.
\item \textbf{MC4}: prove continuity and Mittag--Leffler control for
the completion towers.
\item \textbf{MC5}: construct the BRST-to-bar comparison map and prove
it is a quasi-isomorphism on a named family surface.
\item \textbf{Periodicity}: construct the regular flat modular family
and determine its monodromy representation.
\end{itemize}
The bar object and its comparison maps provide the common algebraic
carrier for these questions.

\subsection{From tree-level to modular homotopy theory}
\label{subsec:tree-to-modular}
\index{modular homotopy theory!three-level ladder}

The constructions of this chapter are \emph{tree-level}: the bar
differential is indexed by collisions on a single smooth curve, and
the operadic compositions are indexed by planar rooted trees. Three
levels of homotopy-algebraic structure emerge as one opens the genus
filtration.

\begin{definition}[Tree-level operadic homotopy theory]
\label{def:tree-level-homotopy}
A \emph{tree-level operadic homotopy theory} for a cyclic operad~$P$
is an $\infty$-category~$\mathcal{H}$ with homotopy-coherent
$P$-action: operations $m_T$ indexed by $P$-decorated rooted
trees~$T$, satisfying the operadic master equation $\sum_T m_T = 0$
up to coherent homotopies. The bar-cobar constructions of this
chapter produce such a theory for $P = \chirCom$ or $P = \chirAss$.
\end{definition}

\begin{definition}[Modular homotopy theory, abstract]
\label{def:modular-homotopy-abstract}
A \emph{modular homotopy theory} for a modular operad~$M$ is an
$\infty$-category~$\mathcal{H}$ with homotopy-coherent $M$-action:
operations $m_\Gamma$ indexed by $M$-decorated \emph{stable
graphs}~$\Gamma$ (graphs with genus labels $g_v$ at each vertex,
satisfying $\sum_v g_v + b_1(\Gamma) = g$), compatible with
separating and nonseparating clutching
(Definition~\ref{def:modular-homotopy-theory-intro}).
\end{definition}

\begin{remark}[The formal step is small; the existence problem is hard]
\label{rem:trees-to-graphs}
Replacing rooted trees by stable graphs is the
\emph{only} combinatorial change from tree-level to modular homotopy
theory. The formal generalization is immediate. What is not
immediate is \emph{existence}: constructing the operations $m_\Gamma$
for graphs with loops requires integration over
$\overline{\mathcal{M}}_{g,n}$, which produces curvature
($\dfib^{\,2} \neq 0$) and demands the full Maurer--Cartan package
$\Theta_\cA$. This is the content of
Chapter~\ref{chap:higher-genus}.
\end{remark}

\medskip

In Chapter~\ref{chap:higher-genus}, the curve itself varies over
$\overline{\mathcal{M}}_g$, and the bar differential inherits
curvature from the open-sector trace and clutching data: the
genus-$g$ correction $m_0^{(g)}$ is a section of the obstruction
sheaf on $\overline{\mathcal{M}}_g$, controlled at the scalar
level by $\kappa(\cA)$ and at higher orders by the shadow
obstruction tower $\Theta_{\cA}^{\leq r}$
(Definition~\ref{def:shadow-postnikov-tower}). The passage
from this chapter to the next is the passage from operadic to
\emph{modular} Koszul duality, from the combinatorics of
associahedra to the geometry of the Deligne--Mumford
compactification, or, in the language of the categorical logarithm,
the passage from the single-valued to the multi-valued regime.

\subsection{Scalar rigidity and the quadratic cone}
\label{subsec:scalar-rigidity-quadratic-frontier}
\index{categorical logarithm!scalar rigidity}
\index{eta-Hessian@$\eta$-Hessian|textbf}
\index{quadratic cone|textbf}

The categorical logarithm encodes a chiral
algebra~$\cA$ as the bar complex $\barB(\cA)$. On the strict
genus-$0$ or associated-graded surface, its Taylor expansion is the PBW
spectral sequence, whose leading coefficient is the modular
characteristic~$\kappa(\cA)$. By contrast, once the higher-genus
fiber operator satisfies
$\dfib^{\,2} = \kappa(\cA) \cdot \omega_g$, the genus transition is
tracked either by the flat comparison differential $\Dg{g}$ or by the
coderived bar-degree filtration, not by ordinary cohomology of the
raw curved fiber complex. On a locus where the full scalar package is
known to be governed by~$\kappa(\cA)$, this curvature controls the
scalar genus tower.

Dropping the stronger scalar-package hypothesis gives a different
picture. The Whitehead decomposition
(Chapter~\ref{chap:higher-genus}) splits the cyclic deformation
cohomology as
\[
H^2_{\mathrm{cyc}}(\cA,\cA)
\;=\;
\mathbb{C}\cdot\eta
\;\oplus\;
H^2_{\mathrm{cyc,prim}}(\cA),
\]
where $\eta$ is the level class and the primitive summand consists of
$\mathfrak{g}$-equivariant cocycles vanishing on current-containing
pairs. Scalar saturation is the vanishing of the primitive summand.
The first non-scalar phenomenon might, na\"ively, be linear transport
by the level class: the operator $l_2^{\mathrm{tr}}(\eta,-)$ acting on
primitive deformations. The geometric surprise is that this operator
vanishes \emph{for formal reasons} (the level class is
cohomologically central), and the first genuine departure from
scalarity is \emph{quadratic}.

The invariant that governs this quadratic cone is a bilinear form
$H_{\eta,\cA}$ on the primitive sector, the \emph{$\eta$-Hessian}:
the transferred ternary $L_\infty$ bracket with one $\eta$-input. In
the singularity-theoretic language that organizes the deformation
theory, the Kodaira--Spencer centrality theorem kills the linear term
of the expansion around the scalar branch, and the $\eta$-Hessian is
the Hessian matrix of the resulting critical point. The quadratic cone
$B_{\kappa,\cA}(y,y) = 0$ in the primitive sector is the tangent cone
to the first non-scalar branch of the moduli space.

\medskip

\noindent\textbf{The deformation-theoretic setup.}\enspace
Throughout this subsection, the transferred $L_\infty$ operations
$\{l_n^{\mathrm{tr}}\}_{n \geq 1}$ are those of the minimal cyclic
$L_\infty$-model on the cohomology
$\mathcal{H} = H^*(\Defcyc(\cA))$, constructed by homotopy transfer
from a dg Lie model of the cyclic deformation complex. The model is
minimal ($l_1^{\mathrm{tr}} = 0$), and the binary bracket
$l_2^{\mathrm{tr}}$ is the induced Lie bracket on cohomology. The
distinguished class $\eta \in \mathcal{H}^2$ is the level class, and
$s^{-1}\eta$ has odd degree in the desuspension $s^{-1}\mathcal{H}$.

\begin{theorem}[Kodaira--Spencer centrality; \ClaimStatusProvedHere]
\label{thm:ks-centrality}
\index{Kodaira--Spencer class!centrality}
Let $(C, \partial, [-,-])$ be a dg Lie algebra equipped with a
one-parameter family of Maurer--Cartan elements
$m_k \in C^1$ depending formally on a scalar parameter~$k$, with the
underlying graded vector space and bracket independent of~$k$. Set
$d_k := \partial + [m_k, -]$ and let
$\eta_k := [\dot{m}_k] \in H^2(C, d_k)$ be the Kodaira--Spencer
class. Then:
\[
\mathrm{ad}_{\eta_k} = 0
\qquad\text{on } H^\bullet(C, d_k).
\]
Equivalently, after passing to any transferred minimal
$L_\infty$-model on the cohomology,
\[
l_2^{\mathrm{tr}}(\eta_k, -) = 0.
\]
\end{theorem}

\begin{proof}
Differentiate the Maurer--Cartan equation
$\partial m_k + \tfrac{1}{2}[m_k, m_k] = 0$ with respect to~$k$:
\[
\partial \dot{m}_k + [m_k, \dot{m}_k] = 0,
\qquad\text{i.e.,}\qquad
d_k \dot{m}_k = 0.
\]
Thus $\eta_k = [\dot{m}_k]$ is a well-defined cohomology class.

Now let $x \in C$ be $d_k$-closed. Differentiating $d_k x = 0$
with respect to~$k$ yields
\[
d_k\!\left(\frac{\partial x}{\partial k}\right) + [\dot{m}_k, x] = 0,
\]
since the bracket is $k$-independent. Hence $[\dot{m}_k, x]$ is
$d_k$-exact, so $\eta_k$ acts trivially on $H^\bullet(C, d_k)$.
The transferred binary bracket on the minimal model is the induced
bracket on cohomology, giving
$l_2^{\mathrm{tr}}(\eta_k, -) = 0$.
\end{proof}

\begin{lemma}[Two-$\eta$ vanishing; \ClaimStatusProvedHere]
\label{lem:two-eta-vanishing}
\index{two-eta vanishing}
Let $(\mathcal{H}, \{l_n\}_{n \geq 1})$ be a minimal $L_\infty$-model
for a deformation problem in which the distinguished level class
$\eta \in \mathcal{H}^2$ has odd suspension $s^{-1}\eta$. Then for
every $n \geq 2$ and all $u_3, \ldots, u_n \in \mathcal{H}$,
\[
l_n(\eta, \eta, u_3, \ldots, u_n) = 0.
\]
In particular, $l_n(\eta, \ldots, \eta) = 0$ for all $n \geq 2$.
\end{lemma}

\begin{proof}
In the desuspension $s^{-1}\mathcal{H}$, the element $s^{-1}\eta$ has
odd degree. The co-operations on the cofree cocommutative coalgebra
$\mathrm{Sym}^c(s^{-1}\mathcal{H})$ are graded symmetric, so
interchanging the two identical odd inputs $s^{-1}\eta$ reproduces the
expression with the Koszul sign $(-1)^{|s^{-1}\eta|^2} = -1$. A value
equal to its own negative is zero.
\end{proof}

\begin{remark}[The scalar direction is affine, not invisible]
\label{rem:scalar-direction-affine}
\index{affine normal form}
The two results above have an immediate structural consequence.
For any $t \in \mathbb{C}$ and $y \in \mathcal{H}$, the
Maurer--Cartan series satisfies
\[
\MC(t\eta + y)
\;=\;
\MC(y) + t\,\Xi_\eta(y),
\qquad
\Xi_\eta(y)
:= \sum_{m \geq 0} \frac{1}{m!}\, l_{m+1}(\eta, y^{\otimes m}).
\]
There are no quadratic or higher powers of~$t$: every term in the
multilinear expansion containing two or more copies of~$\eta$
vanishes by Lemma~\ref{lem:two-eta-vanishing}. The scalar direction
is not invisible (the insertion operator $\Xi_\eta$ carries
genuine structure), but it couples to the remaining directions only
\emph{affinely}. By Theorem~\ref{thm:ks-centrality}, the unary part
of the insertion vanishes on cohomology, so the first surviving
scalar/non-scalar coupling begins quadratically in the primitive
variables.
\end{remark}

The coalgebra-level consequence of the affine normal form is a
square-zero insertion differential, whose spectral sequence collapses
at~$E_2$.

\begin{proposition}[Square-zero insertion differential;
\ClaimStatusProvedHere]
\label{prop:square-zero-insertion}
\index{insertion differential!square-zero}
Let $Q$ denote the coderivation on the cofree conilpotent coalgebra
$T^c(s^{-1}\mathcal{H})$ determined by the Taylor coefficients
$\{l_n\}_{n \geq 1}$, and define a second coderivation $M_\eta$
by $(M_\eta)_n(u_1, \ldots, u_n) := l_{n+1}(\eta, u_1, \ldots, u_n)$.
Then the twist of~$Q$ by the scalar Maurer--Cartan point $t\eta$ is
\[
Q_t = Q + t M_\eta.
\]
Consequently,
\[
[Q, M_\eta] = 0,
\qquad
M_\eta^2 = 0.
\]
\end{proposition}

\begin{proof}
The Taylor coefficient of degree~$n$ in the twist by $t\eta$ is
\[
l_n^{t\eta}(u_1, \ldots, u_n)
= \sum_{r \geq 0} \frac{1}{r!}\,
 l_{n+r}\bigl((t\eta)^{\otimes r}, u_1, \ldots, u_n\bigr).
\]
By Lemma~\ref{lem:two-eta-vanishing}, all terms with $r \geq 2$
vanish, so $l_n^{t\eta} = l_n + t\,l_{n+1}(\eta, -)$, which is the
statement $Q_t = Q + tM_\eta$. Since $t\eta$ is Maurer--Cartan for
every~$t$, one has $Q_t^2 = 0$ for every~$t$. Expanding
$(Q + tM_\eta)^2 = Q^2 + t[Q, M_\eta] + t^2 M_\eta^2$ and comparing
coefficients of $1, t, t^2$ yields the claims.
\end{proof}

\begin{corollary}[Two-step scalar spectral sequence;
\ClaimStatusProvedHere]
\label{cor:two-step-scalar-sseq}
\index{spectral sequence!scalar direction}
Filter the family
$(T^c(s^{-1}\mathcal{H})[[t]], Q_t)$ by powers of~$t$. The
associated spectral sequence has
\[
E_1 \cong H^\bullet(T^c(s^{-1}\mathcal{H}), Q)
 \otimes \mathbb{C}[[t]],
\]
its first differential is induced by $M_\eta$, and all higher
differentials vanish: $d_r = 0$ for $r \geq 2$. The scalar-direction
spectral sequence collapses at~$E_2$.
\end{corollary}

\begin{proof}
The twisted differential $Q_t = Q + tM_\eta$ has exactly two
filtration pieces: $Q$ (preserving the $t$-adic filtration) and
$tM_\eta$ (raising it by one). There is no component raising
filtration by two or more, so the only nontrivial differential
beyond the $Q$-cohomology stage is the first one.
\end{proof}

\medskip

\noindent\textbf{The primitive quadratic cone.}\enspace
Theorems~\ref{thm:ks-centrality} and Lemma~\ref{lem:two-eta-vanishing}
have eliminated the linear obstruction. The first place where non-scalar
behavior can occur is the quadratic interaction between primitive
directions and the scalar branch.

\begin{notation}[Primitive sector and scalar basepoint]
\label{not:primitive-sector}
Write
\[
\mathcal{P}_\cA := H^2_{\mathrm{cyc,prim}}(\cA)
\]
for the primitive degree-two sector from the Whitehead decomposition.
Suppress, for notational simplicity, the inert Hodge factor
$\Lambda = \sum_{g \geq 1} \lambda_g$ and write
\[
x_\kappa := \kappa(\cA)\,\eta
\]
for the scalar Maurer--Cartan point in the minimal model.
\end{notation}

\begin{theorem}[Quadratic cone at the scalar basepoint; \ClaimStatusProvedHere]
\label{thm:quadratic-frontier}
\index{quadratic cone!theorem}
Let $\cA$ be any object for which the Whitehead decomposition holds and
for which the scalar branch $x_\kappa = \kappa(\cA)\eta$ is
Maurer--Cartan. Let $y \in \mathcal{P}_\cA$. Then:
\begin{enumerate}[label=\textup{(\alph*)}]
\item Every primitive direction is first-order unobstructed along the
scalar branch:
\[
l_1^{x_\kappa}(y) = 0.
\]
\item The first nontrivial obstruction is quadratic:
\[
\MC(x_\kappa + y)
= \tfrac{1}{2}\, B_{\kappa,\cA}(y,y) + O(y^3),
\]
where
\[
B_{\kappa,\cA}(u,v)
\;:=\;
l_2^{\mathrm{tr}}(u,v)
+ \kappa(\cA)\,l_3^{\mathrm{tr}}(\eta, u, v).
\]
\end{enumerate}
Consequently, the first non-scalar branch is governed by the quadratic
cone
\[
\mathcal{C}_{\kappa,\cA}
:= \bigl\{
 y \in \mathcal{P}_\cA
 \;:\;
 B_{\kappa,\cA}(y,y) = 0
\bigr\}.
\]
\end{theorem}

\begin{proof}
The twisted unary bracket is
\[
l_1^{x_\kappa}(y)
= \sum_{r \geq 0} \frac{1}{r!}\,
 l_{1+r}(x_\kappa^{\otimes r}, y).
\]
The $r = 0$ term vanishes because the model is minimal. The $r = 1$
term is $\kappa(\cA)\, l_2^{\mathrm{tr}}(\eta, y)$, which vanishes by
Theorem~\ref{thm:ks-centrality}. All terms with $r \geq 2$ contain at
least two copies of~$\eta$ and vanish by
Lemma~\ref{lem:two-eta-vanishing}. This proves~(a).

For the quadratic term, expand around $x_\kappa$:
\[
\MC(x_\kappa + y)
= \underbrace{\MC(x_\kappa)}_{=\,0}
 + \underbrace{l_1^{x_\kappa}(y)}_{=\,0}
 + \tfrac{1}{2}\, l_2^{x_\kappa}(y,y)
 + O(y^3).
\]
Compute
\[
l_2^{x_\kappa}(u,v)
= \sum_{r \geq 0} \frac{1}{r!}\,
 l_{2+r}(x_\kappa^{\otimes r}, u, v).
\]
The $r = 0$ term is $l_2^{\mathrm{tr}}(u,v)$; the $r = 1$ term is
$\kappa(\cA)\, l_3^{\mathrm{tr}}(\eta, u, v)$; all $r \geq 2$ terms
vanish by Lemma~\ref{lem:two-eta-vanishing}.
\end{proof}

\begin{definition}[$\eta$-Hessian]
\label{def:eta-hessian}
\index{eta-Hessian@$\eta$-Hessian!definition}
The \emph{$\eta$-Hessian} of the family at~$\cA$ is the bilinear form
on the primitive sector defined by
\[
H_{\eta,\cA}(u,v)
\;:=\;
l_3^{\mathrm{tr}}(\eta, u, v),
\qquad u, v \in \mathcal{P}_\cA.
\]
The quadratic-cone tensor decomposes as
\[
B_{\kappa,\cA}
\;=\;
l_2^{\mathrm{tr}}\big|_{\mathcal{P}_\cA \otimes \mathcal{P}_\cA}
\;+\;
\kappa(\cA)\, H_{\eta,\cA}.
\]
The first summand is the intrinsic binary bracket on primitive
cohomology; the second is the new invariant introduced by the scalar
curvature.
\end{definition}

\begin{proposition}[Homotopy-transfer construction of the
$\eta$-Hessian; \ClaimStatusProvedHere]
\label{prop:eta-hessian-transfer}
\index{eta-Hessian@$\eta$-Hessian!transfer formula}
\index{homotopy transfer!ternary bracket}
Let $(C, d, [-,-])$ be a dg Lie model for the cyclic deformation
complex of~$\cA$, and choose a contraction
\[
(\mathcal{P}_\cA \oplus \mathbb{C}\eta)
 \xrightarrow{\;i\;} C
 \xrightarrow{\;p\;} (\mathcal{P}_\cA \oplus \mathbb{C}\eta),
\qquad
dh + hd = \mathrm{id}_C - ip.
\]
Then $H_{\eta,\cA}$ is represented by the standard transferred
rooted-tree formula:
\begin{align*}
H_{\eta,\cA}(u,v)
&= p\bigl([h[i(\eta), i(u)],\, i(v)]\bigr)
 + p\bigl([h[i(u), i(v)],\, i(\eta)]\bigr) \\
&\qquad
 + p\bigl([h[i(v), i(\eta)],\, i(u)]\bigr),
\end{align*}
with Koszul signs in the general graded case. The formula is
independent of the choice of contraction, by the standard
homotopy-transfer invariance of $L_\infty$-structures.
\end{proposition}

\begin{proof}
For a dg Lie algebra, the transferred ternary bracket on the minimal
$L_\infty$-model is the sum over rooted binary trees with three leaves
and one internal homotopy edge. Evaluating at the ordered triple
$(\eta, u, v)$ yields the displayed expression. Independence of the
contraction is homotopy-transfer invariance.
\end{proof}

\begin{lemma}[Shifted symmetry on degree-two primitives;
\ClaimStatusProvedHere]
\label{lem:shifted-symmetry-H}
\index{eta-Hessian@$\eta$-Hessian!symmetry}
On the primitive degree-two sector
$\mathcal{P}_\cA \subset H^2_{\mathrm{cyc}}(\cA, \cA)$, both
$l_2^{\mathrm{tr}}(u,v)$ and $H_{\eta,\cA}(u,v)$ are symmetric in the
ordinary sense. Consequently, the quadratic-cone tensor
$B_{\kappa,\cA}(u,v)$ is symmetric on~$\mathcal{P}_\cA$.
\end{lemma}

\begin{proof}
After desuspension, degree-two classes become odd.
The transferred bracket is graded antisymmetric in the shifted sense,
so exchanging two odd inputs contributes the extra minus sign coming
from their parity.
That sign cancels the shifted antisymmetry sign, leaving ordinary
symmetry on the primitive degree-two sector.
The same argument applies to both
$l_2^{\mathrm{tr}}(u,v)$ and the Hessian $H_{\eta,\cA}(u,v)$, so
$B_{\kappa,\cA}(u,v)$ is symmetric as well.
\end{proof}

\begin{remark}[The first obstruction is quadratic]
\label{rem:frontier-moved}
\index{quadratic cone!interpretation}
The natural question was to compute $l_2^{\mathrm{tr}}(\eta, -)$ and
its induced quotient differential.
Theorem~\ref{thm:ks-centrality} and
Theorem~\ref{thm:quadratic-frontier} show that this is the wrong
linear invariant. Linear transport by the level class vanishes formally. The
first genuine deformation-theoretic information lives instead in the
$\eta$-Hessian $H_{\eta,\cA}$ and the resulting quadratic cone
$\mathcal{C}_{\kappa,\cA}$. In the singularity-theoretic dictionary:
one has killed the gradient of the Maurer--Cartan potential along the
scalar branch; the $\eta$-Hessian is the Hessian at the resulting
critical point; and the quadratic cone is the tangent cone to the
critical locus.
\end{remark}

\medskip

\noindent\textbf{Cyclic rigidity in the admissible sector.}\enspace
The one-channel locus extends well beyond the generic regime. Generic
rigidity (Chapter~\ref{chap:higher-genus}) gives
$H^2_{\mathrm{cyc,prim}}(\cA) = 0$ whenever the relevant module
category is semisimple. The first non-scalar quadratic cone therefore
lies beyond the semisimple regime. Even admissible levels, where
semisimplicity can fail, remain one-channel at the minimal
cyclic level in their semisimple highest-weight sector.

\begin{theorem}[Admissible cyclic rigidity; \ClaimStatusProvedHere]
\label{thm:admissible-scalar-rigidity}
\index{admissible level!scalar rigidity}
\index{scalar saturation!admissible sector}
Let $k$ be an admissible level for a simple Lie algebra~$\mathfrak{g}$,
and let $\cA = L_k(\mathfrak{g})$ be the corresponding simple admissible
affine vertex algebra. Consider the cyclic first-order deformation
problem internal to the semisimple admissible highest-weight module
category of~\cite{Arakawa2016RationalAdmissible}. Then
\[
H^2_{\mathrm{cyc}}(\cA, \cA) = \mathbb{C}\eta,
\qquad
H^2_{\mathrm{cyc,prim}}(\cA) = 0.
\]
In particular, the semisimple admissible highest-weight sector of the
vacuum admissible affine algebra has one-dimensional cyclic
line concentration.
\end{theorem}

\begin{proof}
The Whitehead decomposition applies at all non-critical levels, so
only the vanishing of the primitive part must be addressed. The
generic rigidity argument kills primitive cyclic deformations whenever
the relevant module category is semisimple: the algebra object then
has vanishing topological Hochschild cohomology $\mathrm{Ext}^2$ in that semisimple tensor
category. For admissible~$k$,
Arakawa~\cite{Arakawa2016RationalAdmissible} proves that the ordinary
admissible highest-weight category is semisimple and that
$L_k(\mathfrak{g})$ belongs to it. Running the generic rigidity
argument inside that semisimple sector gives
$H^2_{\mathrm{cyc,prim}}(\cA) = 0$.
\end{proof}

\begin{corollary}[Drinfeld--Sokolov reductions remain one-channel in
the semisimple admissible sector;
\ClaimStatusConditional]
\label{cor:ds-not-first-frontier}
\index{Drinfeld--Sokolov reduction!scalar rigidity}
Let $k$ be admissible and $f \in \mathfrak{g}$ nilpotent. Then the
$\mathcal{W}$-algebra
$\mathcal{W}^k(\mathfrak{g}, f)
 = \mathrm{DS}_f(V_k(\mathfrak{g}))$
has one-dimensional cyclic direction in that semisimple sector.
In particular, neither the vacuum admissible
affine algebra nor its ordinary Drinfeld--Sokolov reductions provide
the first admissible example with genuinely nonzero primitive sector.
\end{corollary}

\begin{proof}
Apply Theorem~\ref{thm:admissible-scalar-rigidity} and the
BRST pullback argument of
Proposition~\ref{prop:saturation-functorial}(a), which gives
one-dimensional cyclic direction for the Drinfeld--Sokolov
reduction.
\end{proof}

\begin{remark}[The first admissible non-scalar source]
\label{rem:true-admissible-frontier}
\index{admissible level!nonsemisimple source}
The first admissible family with nonzero primitive sector must lie
beyond the semisimple ordinary highest-weight regime. Natural targets
include: logarithmic extensions of admissible affine algebras,
admissible conformal cosets with Jordan blocks, relaxed highest-weight
or nonordinary sectors, and nonrational $\mathcal{W}$-algebras for
which the ordinary category ceases to be semisimple. In every case,
the route to the first non-scalar quadratic cone passes through the
explicit computation described below.
\end{remark}

\medskip

\noindent\textbf{The scalar genus as a classified multiplicative
genus.}\enspace
On a locus where the scalar genus tower is known to depend only on the
single number~$\kappa(\cA)$, the following classification theorem
makes this uniqueness precise: any
multiplicative genus factoring through the modular characteristic must
take the universal form.

\begin{theorem}[Classification of scalar genera \textup{(}uniform-weight\textup{)};
\ClaimStatusConditional]
\label{thm:classification-scalar-genera}
\index{scalar logarithm!classification}
\index{multiplicative genus!scalar}
Let $\mathcal{U} := 1 + x^2\mathbb{C}[[x^2]]$ with multiplication of
formal series. Suppose that to each object~$\cA$ in a one-scalar
class one assigns an element $Z(\cA; x) \in \mathcal{U}$ such that:
\begin{enumerate}[label=\textup{(\alph*)}]
\item $Z(\cA \star \mathcal{B}; x)
 = Z(\cA; x)\, Z(\mathcal{B}; x)$ whenever
 $\kappa(\cA \star \mathcal{B})
 = \kappa(\cA) + \kappa(\mathcal{B})$;
\item $Z(\cA; x)$ depends on~$\cA$ only through $\kappa(\cA)$;
\item the dependence on~$\kappa$ is formal power-series dependence.
\end{enumerate}
Then there exists a unique series
$\Psi(x) \in x^2\mathbb{C}[[x^2]]$ such that
\[
Z(\cA; x) = \exp\bigl(\kappa(\cA)\,\Psi(x)\bigr)
\]
for every~$\cA$. If the genus coefficients are normalized to the
scalar genus formula
\[
\sum_{g \geq 1} F_g(\cA)\, x^{2g}
= \kappa(\cA)\left(\frac{x/2}{\sin(x/2)} - 1\right),
\]
then $\Psi(x) = \Phi_{\log}(x) := \frac{x/2}{\sin(x/2)} - 1$ and
\[
Z(\cA; x) = \exp\bigl(\kappa(\cA)\,\Phi_{\log}(x)\bigr).
\]
\end{theorem}

\begin{proof}
Write $Z_z(x)$ for the value of~$Z$ at an object of scalar
characteristic~$z$. Multiplicativity gives
$Z_{z+w}(x) = Z_z(x)\, Z_w(x)$. Setting
$f_z(x) := \log Z_z(x) \in x^2\mathbb{C}[[x^2]]$ yields
$f_{z+w}(x) = f_z(x) + f_w(x)$. Expanding formally,
$f_z(x) = \sum_{n \geq 0} c_n(x)\, z^n$; the additive functional
equation forces $c_n = 0$ for $n \neq 1$, giving
$f_z(x) = z\,\Psi(x)$ for a unique~$\Psi$. The final statement
follows from the scalar genus formula
(Theorem~\ref{thm:modular-characteristic}).
\end{proof}

\begin{remark}[Self-dual anomaly neutrality]
\label{rem:self-dual-anomaly-neutrality}
\index{self-duality!anomaly neutrality}
For algebras with $\kappa(\cA^!) = -\kappa(\cA)$ (the exact
anti-symmetry that holds for affine Kac--Moody and free fields;
Theorem~\ref{thm:modular-characteristic}(iii)), Koszul self-duality
$\cA \simeq \cA^!$ forces $\kappa(\cA) = 0$ and hence $F_g(\cA) = 0$
for all $g \geq 1$ \textup{(}uniform-weight\textup{)}: the scalar logarithm of such a Koszul-self-dual chiral
algebra vanishes identically. For $\mathcal{W}$-algebras, the
complementarity sum $\kappa + \kappa'$ is a nonzero constant
(e.g.\ $\kappa + \kappa' = 13$ for Virasoro), so Koszul
self-duality $\mathrm{Vir}_{13} \simeq \mathrm{Vir}_{26-13}
= \mathrm{Vir}_{13}$ (Koszul duality at $c = 13$)
gives $\kappa(\mathrm{Vir}_{13}) = 13/2 \neq 0$: the
Virasoro Koszul self-dual point ($c=13$) is \emph{not} at critical level.
See Theorem~\ref{thm:quantum-complementarity-main} for the full
complementarity.
\end{remark}

\medskip

\noindent\textbf{The platonic hierarchy.}\enspace
The results of this subsection organize the categorical logarithm at
the deformation-theoretic level into five steps, each resolving a
na\"ive expectation.

\begin{theorem}[Five-step hierarchy of the categorical logarithm;
\ClaimStatusConditional]
\label{thm:platonic-hierarchy-log}
\index{categorical logarithm!five-step hierarchy}
Assume the Whitehead decomposition and the scalar genus formula
(Chapter~\ref{chap:higher-genus}). Then the deformation theory of the
categorical logarithm organizes itself as follows.
\begin{enumerate}[label=\textup{(\roman*)}]
\item \emph{Logarithmic linearization.} The bar construction is the
categorical logarithm
$\log^{\mathrm{cat}}(\cA) = \barB_X(\cA)$ and the cobar construction
is the categorical exponential
$\exp^{\mathrm{cat}}(\mathcal{C}) = \Omega_X(\mathcal{C})$. The four
main theorems are existence, local inversion, branch structure, and
leading coefficient for this logarithm.

\item \emph{Scalar package.} On a locus where the full scalar genus
tower is known to depend only on~$\kappa(\cA)$, it is governed by the
scalar logarithm
\[
\kappa(\cA)\left(\frac{x/2}{\sin(x/2)} - 1\right).
\]
The exponential
$\exp\bigl(\kappa(\cA)\, \Phi_{\log}(x)\bigr)$ is the unique
multiplicative genus factoring through~$\kappa$
\textup{(}Theorem~\textup{\ref{thm:classification-scalar-genera})}.

\item \emph{Vanishing of the linear obstruction.} For every genuine
one-parameter level family, the Kodaira--Spencer class is
cohomologically central:
$l_2^{\mathrm{tr}}(\eta, -) = 0$
\textup{(}Theorem~\textup{\ref{thm:ks-centrality})}. There is no
first-order non-scalar obstruction along the scalar branch.

\item \emph{First non-scalar invariant: the $\eta$-Hessian.}
The first genuine departure from scalarity is the $\eta$-Hessian
$H_{\eta,\cA}(u,v) = l_3^{\mathrm{tr}}(\eta, u, v)$, and the first
non-scalar branch is governed by the quadratic cone
$B_{\kappa,\cA}(y,y) = 0$
\textup{(}Theorem~\textup{\ref{thm:quadratic-frontier}},
Definition~\textup{\ref{def:eta-hessian})}.

\item \emph{First admissible non-scalar source.} In semisimple admissible highest-weight
sectors the primitive cohomology still vanishes
\textup{(}Theorem~\textup{\ref{thm:admissible-scalar-rigidity})}, so
the first family carrying genuine non-scalar data must be
nonsemisimple admissible.
\end{enumerate}
\end{theorem}

\begin{proof}
Parts~(i)--(ii) are the scalar genus formula repackaged by
Definition~\ref{def:categorical-log-exp} and
Theorem~\ref{thm:classification-scalar-genera}. Part~(iii) is
Theorem~\ref{thm:ks-centrality}. Part~(iv) is
Theorem~\ref{thm:quadratic-frontier} and
Definition~\ref{def:eta-hessian}. Part~(v) is
Theorem~\ref{thm:admissible-scalar-rigidity} and
Corollary~\ref{cor:ds-not-first-frontier}.
\end{proof}

\begin{remark}[Connection to the shadow obstruction tower]
\label{rem:eta-hessian-shadow-tower}
\index{shadow obstruction tower!$\eta$-Hessian}
The $\eta$-Hessian $H_{\eta,\cA}$ is the first non-scalar shadow in
the obstruction tower $\Theta_\cA^{\leq r}$
(Definition~\ref{def:shadow-postnikov-tower}). At the scalar level
($r = 2$), the tower is determined by $\kappa(\cA)$. The quadratic
cone tensor $B_{\kappa,\cA}$ records exactly the degree-$3$
correction: the transferred ternary bracket with one $\eta$-input.
The hierarchy $\kappa \to H_{\eta,\cA} \to \cdots$ is the
leading-term expansion of the full shadow obstruction tower, and the
$\eta$-Hessian is the first invariant that the scalar characteristic
does not see. Beyond the quadratic level, the tower continues with
the quartic resonance class $Q^{\mathrm{contact}}_\cA$
and higher shadows, as developed in
Chapter~\ref{chap:higher-genus}.
\end{remark}

\begin{remark}[Computing the first non-scalar branch]
\label{rem:computing-first-frontier}
\index{eta-Hessian@$\eta$-Hessian!computation procedure}
The practical recipe for detecting the first non-scalar branch is:
\begin{enumerate}[label=\textup{(\arabic*)}]
\item Build a dg Lie or cyclic $L_\infty$ model $C_k$ for
 $\Defcyc(\cA_k)$ with fixed underlying graded space and explicit
 $k$-dependence in the Maurer--Cartan element.
\item Compute
 $\mathcal{P}_{\cA_k} = H^2_{\mathrm{cyc,prim}}(\cA_k)$; if it
 vanishes, the family remains scalar-saturated.
\item Choose a contraction onto
 $\mathcal{P}_{\cA_k} \oplus \mathbb{C}\eta_k$ and evaluate the
 transferred ternary tree formula
 (Proposition~\ref{prop:eta-hessian-transfer}) to obtain
 $H_{\eta_k, \cA_k}$.
\item Form $B_{\kappa,\cA_k}
 = l_2^{\mathrm{tr}} + \kappa(\cA_k)\, H_{\eta_k, \cA_k}$ and solve
 the quadratic cone $B_{\kappa,\cA_k}(y,y) = 0$.
\end{enumerate}
This is the exact analogue of computing the Hessian of a singularity
once the linear term has been killed. One does not need the full
nonlinear theory; the branch is detected by a single transferred
ternary bracket.
\end{remark}

% ================================================================
% ČECH RESOLUTION AND HOMOTOPY CHIRAL ALGEBRAS
% ================================================================

\section{\v{C}ech resolution and homotopy chiral algebras}
\label{sec:cech-hca}
\index{homotopy chiral algebra}
\index{Cech complex@\v{C}ech complex!HCA structure}

The preceding sections established bar-cobar inversion via
configuration-space techniques on $\operatorname{Ran}(X)$.
A complementary approach works on the Zariski site of the
curve itself: the \v{C}ech complex of a sheaf of chiral
algebras.
Malikov--Schechtman~\cite{MS24} prove that this \v{C}ech complex
carries the structure of a \emph{homotopy chiral algebra}
(HCA), a chiral algebra up to coherent higher homotopies, with
operadic control by the Lie Eilenberg--Zilber operad.

The HCA operations $\{F_n\}$
are the \emph{components} of
a single Maurer--Cartan element in an appropriate convolution
dg~Lie algebra, the \v{C}ech convolution algebra
$\mathfrak{g}^{\check{C}}(\cA,\mathcal{U})$
(\S\ref{sec:categorical-logarithm},
Chapter~\ref{chap:higher-genus}). The passage from local data (sections on affine opens) to
global data (cohomology on~$X$) carries algebraic structure through the Moore
complex functor: an $L_\infty$
enhancement of the Lie bracket in Beilinson--Drinfeld's
pseudo-tensor category~\cite{BD04}.

\subsection{The cosimplicial Lie algebra of sections}
\label{subsec:cosimplicial-cech}

Let $\cA$ be a sheaf of chiral algebras on a smooth curve~$X$
(or more generally on a smooth variety), and let
$\mathcal{U} = \{U_i\}_{i \in I}$ be an affine open cover.
The sections on iterated intersections assemble into a
\emph{cosimplicial} object in the pseudo-tensor category of
$\cD_X$-modules:

\begin{definition}[Cosimplicial \v{C}ech object]
\label{def:cosimplicial-cech}
Define
\[
\check{C}^p(\mathcal{U};\cA)
 = \bigoplus_{i_0 < \cdots < i_p}
 \Gamma(U_{i_0} \cap \cdots \cap U_{i_p},\;\cA)
\]
with coface maps
$\partial^j\colon \check{C}^p \to \check{C}^{p+1}$
given by restriction (omitting the $j$-th index) and
codegeneracy maps
$\sigma^j\colon \check{C}^{p+1} \to \check{C}^p$ given by
diagonal insertion. The alternating sum
$d_{\check{C}} = \sum_{j=0}^{p+1}(-1)^j\partial^j$
is the \v{C}ech differential, and $d_{\check{C}}^2 = 0$ by the
standard cosimplicial identities.
\end{definition}

The chiral product on $\cA$ endows each
$\Gamma(U_{i_0\cdots i_p},\cA)$ with a Lie bracket in the
pseudo-tensor sense of Beilinson--Drinfeld~\cite{BD04}:
\[
\mu_{i_0\cdots i_p}\colon
 j_*j^*\bigl(\cA|_{U_{i_0\cdots i_p}}
 \boxtimes \cA|_{U_{i_0\cdots i_p}}\bigr)
 \longrightarrow
 \Delta_*\cA|_{U_{i_0\cdots i_p}}.
\]
The restriction maps are compatible with these brackets, so
$\check{C}^\bullet(\mathcal{U};\cA)$ is a cosimplicial Lie algebra
in the pseudo-tensor category of $\cD_X$-modules.

\subsection{The \v{C}ech convolution dg~Lie algebra}
\label{subsec:cech-convolution}
\index{convolution dg Lie algebra!\v{C}ech}

The Moore complex $M\check{C}^\bullet(\mathcal{U};\cA)$ is a cochain
complex of $\cD_X$-modules whose cohomology computes the sheaf
cohomology $H^\bullet(X,\cA)$. As in
Proposition~\ref{prop:chriss-ginzburg-structure}, we
package the HCA operations into a single convolution algebra.

\begin{definition}[Moore complex]
\label{def:moore-complex}
For a cosimplicial object $A^\bullet$ in an abelian category, the
\emph{Moore complex} (or normalized cochain complex) is
\[
(MA)^p = \bigcap_{j=0}^{p-1}\ker(\sigma^j\colon A^p \to A^{p-1}),
\qquad
d_M = \sum_{j=0}^{p+1}(-1)^j\partial^j\big|_{(MA)^p}.
\]
The Dold--Kan correspondence identifies cosimplicial objects with
cochain complexes via this functor.
\end{definition}

\begin{definition}[\v{C}ech convolution dg~Lie algebra]
\label{def:cech-convolution}
\index{gcheckC@$\mathfrak{g}^{\check{C}}$}
Let $V = M\check{C}^\bullet(\mathcal{U};\cA)$ denote the Moore
complex. Define the \emph{\v{C}ech convolution dg~Lie algebra}
\[
\mathfrak{g}^{\check{C}}(\cA,\mathcal{U})
 \;=\;
 \prod_{n \geq 1}
 \operatorname{Hom}_{\Sigma_n}\!\bigl(
 \mathrm{Lie}(n) \otimes \mathcal{Z}(n),\;
 \operatorname{Hom}(V^{\otimes n},V)
 \bigr),
\]
where $\mathcal{Z}(n) =
\operatorname{Tot}\operatorname{Hom}_{\Delta(k)}(Z,Z^{\otimes n})$
is the Eilenberg--Zilber chain complex and $\mathrm{Lie}(n)$ is
the degree-$n$ part of the Lie operad.
The grading is
\[
\bigl|\mathfrak{g}^{\check{C}}\bigr|^k
 = \prod_{n \geq 1}
 \operatorname{Hom}_{\Sigma_n}\!\bigl(
 \mathrm{Lie}(n) \otimes \mathcal{Z}(n),\;
 \operatorname{Hom}(V^{\otimes n}, V[k+1-n])
 \bigr).
\]
The differential $D$ on $\mathfrak{g}^{\check{C}}$ is the
commutator $D(\alpha) = [d_M, \alpha]$ with the Moore differential.
The Lie bracket on $\mathfrak{g}^{\check{C}}$ is the operadic
composition bracket:
\[
[\alpha,\beta](a_1,\ldots,a_{m+n-1})
 = \sum_{i=1}^m \pm\,
 \alpha(a_1,\ldots,\beta(a_i,\ldots,a_{i+n-1}),\ldots,a_{m+n-1})
 \;-\; (\alpha \leftrightarrow \beta),
\]
where $\alpha \in \mathfrak{g}^{\check{C}}$ has degree~$m$ and
$\beta$ has degree~$n$.
\end{definition}

\begin{remark}[\v{C}ech and modular convolution algebras]
\label{rem:cech-modular-parallel}
The algebra $\mathfrak{g}^{\check{C}}(\cA,\mathcal{U})$ is the
\v{C}ech convolution dg~Lie algebra encoding HCA operations as
MC data. Its modular counterpart is
$\mathfrak{g}^{\mathrm{mod}}_\cA$
(Definition~\ref{def:modular-convolution-dg-lie}), and the two share
the same formal architecture:
\begin{enumerate}[label=\textup{(\roman*)}]
\item both encode all algebraic operations on a resolution
 of~$\cA$ as components of a single MC element
 ($\Phi_\cA$ and $\Theta_\cA$ respectively);
\item both are \emph{strict} dg~Lie models of underlying shifted
 $L_\infty$~algebras
 (Convention~\ref{rem:two-level-convention}): the strict bracket
 suffices for Maurer--Cartan theory, while the full $L_\infty$
 structure governs homotopy transfer and gauge equivalence;
\item both produce degree-by-degree truncations
 ($\Phi_\cA^{\le r}$, $\Theta_\cA^{\le r}$) that satisfy the
 MC equation modulo degree~$r{+}1$, forming parallel Postnikov
 towers.
\end{enumerate}
The structural difference is the indexing combinatorics:
$\mathfrak{g}^{\mathrm{mod}}_\cA$ is indexed by stable graphs
(capturing all genera), while $\mathfrak{g}^{\check{C}}$ is indexed
by the simplicial combinatorics of the cover (capturing all \v{C}ech
degrees). The comparison between them is an $L_\infty$ morphism
that intertwines the two MC elements.
\end{remark}

\subsection{HCA operations as a Maurer--Cartan element}
\label{subsec:hca-mc}

The theorem of Malikov--Schechtman is now a statement about
existence of an MC element:

\begin{theorem}[\v{C}ech complex as homotopy chiral algebra;
\ClaimStatusProvedElsewhere]
\label{thm:cech-hca}
\index{Malikov--Schechtman theorem}
\textup{(Malikov--Schechtman \cite{MS24}, Theorem~3.1.)}
Let $\cA$ be a sheaf of chiral algebras on~$X$ and
$\mathcal{U}$ an affine open cover. The HCA structure on the Moore
complex $V = M\check{C}^\bullet(\mathcal{U};\cA)$ is encoded by a
Maurer--Cartan element
\begin{equation}\label{eq:cech-mc-element}
\Phi_\cA
 \;=\; \sum_{n \geq 1} F_n
 \;\in\;
 \mathfrak{g}^{\check{C}}(\cA,\mathcal{U})^1
\end{equation}
satisfying the master equation
\begin{equation}\label{eq:cech-master-equation}
D\,\Phi_\cA + \tfrac{1}{2}[\Phi_\cA,\Phi_\cA] = 0
\end{equation}
in $\mathfrak{g}^{\check{C}}(\cA,\mathcal{U})$.

Concretely, each component $F_n$ is a multilinear operation
\[
F_n\colon
 \mathrm{Lie}(n) \otimes \mathcal{Z}(n)
 \otimes V^{\otimes n}
 \longrightarrow V,
 \qquad n \geq 1,
\]
where $\mathcal{Z}(n)$ is the Eilenberg--Zilber chain complex, and
the collection $\{F_n\}_{n \geq 1}$ defines a pseudo-tensor
functor from the Lie Eilenberg--Zilber operad
$\mathcal{Y}(n) = \mathcal{Z}(n) \otimes_k \mathrm{Lie}(n)$
to the endomorphism operad of~$V$ in the pseudo-tensor category.
\end{theorem}

\begin{remark}[Degree-by-degree expansion of the MC equation]
\label{rem:cech-mc-degree}
Expanding~\eqref{eq:cech-master-equation} by degree yields the
familiar $L_\infty$ relations.
Write $\Phi_\cA = F_1 + F_2 + F_3 + \cdots$ where $F_n$ has
degree~$n$.

\smallskip\noindent
\textbf{Degree 1} ($n=1$): $D(F_1) = 0$. Since
$F_1 = d_M$ is the Moore differential, this says $d_M^2 = 0$.

\smallskip\noindent
\textbf{Degree 2} ($n=2$):
$D(F_2) + \tfrac{1}{2}[F_1,F_1] = 0$.
This says $[d_M, F_2] = 0$: the binary bracket $F_2$ is a chain map.

\smallskip\noindent
\textbf{Degree 3} ($n=3$):
\begin{equation}\label{eq:cech-degree-3}
D(F_3) + [F_1,F_2] + \tfrac{1}{6}[F_2,[F_2,F_2]] = 0.
\end{equation}
The term $[F_2,[F_2,F_2]]$ is the Jacobiator of~$F_2$. So $F_3$
is the nullhomotopy of the Jacobiator: $[d_M, F_3] =
-\operatorname{Jac}(F_2)$.

\smallskip\noindent
\textbf{Degree $n$} ($n \geq 4$): each relation has the form
$D(F_n) + (\text{compositions of } F_1,\ldots,F_{n-1}) = 0$,
asserting that $F_n$ is the nullhomotopy of the obstruction
assembled from lower-degree data. The tower
$\Phi_\cA^{\leq r} = \sum_{n=1}^r F_n$ satisfies the MC equation
modulo degree~$r{+}1$, precisely paralleling the shadow obstruction
tower $\Theta_\cA^{\leq r}$
(Definition~\ref{def:shadow-postnikov-tower}).
\end{remark}

\begin{remark}[Operadic action map]
\label{rem:operadic-action-map}
The explicit formula for $F_n$ is given by
Malikov--Schechtman~\cite{MS24}, equation~(3.1.2.1):
for $\psi \in \mathrm{Lie}(n)$,
$\varphi \in \operatorname{Hom}_\Delta(Z, Z^{\otimes n})$,
and elements
$a_{m_1} \in V^{m_1}, \ldots,
 a_{m_n} \in V^{m_n}$,
\[
F_n(\psi \otimes \varphi)_{m_1,\ldots,m_n;\,m}(a_{m_1},\ldots,a_{m_n})
=
\begin{cases}
 \psi_m(\bar\varphi)(a_{m_1},\ldots,a_{m_n})
 & \text{if } m = m_1 + \cdots + m_n, \\
 0 & \text{otherwise},
\end{cases}
\]
where $\psi_m$ is the chiral Lie action on
$\check{C}^m(\mathcal{U};\cA)$ and
$\bar\varphi \in \bigotimes_{i=1}^n
\operatorname{Hom}(\check{C}^{m_i},\check{C}^m)$
is the component of~$\varphi$ at the specified multi-degree.
\end{remark}

\subsection{\v{C}ech--bar comparison as $L_\infty$ morphism}
\label{subsec:cech-bar-comparison}
\index{Cech--bar comparison@\v{C}ech--bar comparison}

The bar complex and the \v{C}ech complex are two resolutions of the
same object (the chiral homology of~$\cA$ on~$X$), and the
convolution formulation makes their compatibility a statement
about MC elements in two convolution algebras.

\begin{conjecture}[\v{C}ech--bar intertwining;
\ClaimStatusConjectured]
\label{conj:cech-bar-intertwining}
Let $\cA$ be a chiral algebra on a smooth projective curve~$X$.
There exists an $L_\infty$ morphism
\[
\Phi_*\colon
 \bigl(\mathfrak{g}^{\check{C}}(\cA,\mathcal{U}),\;
 \Phi_\cA\bigr)
 \xrightarrow{\;\sim\;}
 \bigl(\Convinf(\barBgeom(\cA),\cA),\;
 \Theta^{(0)}_\cA\bigr)
\]
intertwining the \v{C}ech MC element $\Phi_\cA = \sum F_n$ with the
genus-$0$ bar MC element $\Theta^{(0)}_\cA = \sum m_n$. On
cohomology, $\Phi_*$ induces the canonical isomorphism
$\check{H}^\bullet(X,\cA) \cong H_\bullet^{\mathrm{ch}}(X,\cA)$
between \v{C}ech cohomology and chiral homology.

The $L_\infty$ morphism $\Phi_*$ is constructed by the standard
comparison between \v{C}ech and de~Rham/Dolbeault resolutions,
enhanced to an $L_\infty$ morphism by homological perturbation.
It is generically \emph{non-strict}: the component
$(\Phi_*)_n\colon F_{i_1} \otimes \cdots \otimes F_{i_n}
\to m_{i_1 + \cdots + i_n - n + 1}$
need not vanish for $n \geq 2$. Strictness holds precisely when
the \v{C}ech complex has the same truncation depth as the bar
filtration, a condition guaranteed for two-element covers
(Proposition~\ref{prop:cech-two-element-strict}).
\end{conjecture}

\begin{remark}[\v{C}ech--bar convolution dictionary]
\label{rem:cech-bar-dictionary}
The MC formulation yields a parallel between the two resolutions:
\begin{center}
\renewcommand{\arraystretch}{1.3}
\begin{tabular}{lll}
& \textbf{\v{C}ech side} & \textbf{Bar side} \\
\hline
\textit{Ambient algebra}
 & $\mathfrak{g}^{\check{C}}(\cA,\mathcal{U})$
 & $\Convinf(\barBgeom(\cA),\cA)$ \\
\textit{MC element}
 & $\Phi_\cA = \sum_{n \geq 1} F_n$
 & $\Theta^{(0)}_\cA = \sum_{n \geq 1} m_n$ \\
\textit{Master equation}
 & $D\,\Phi + \tfrac{1}{2}[\Phi,\Phi] = 0$
 & $D\,\Theta + \tfrac{1}{2}[\Theta,\Theta] = 0$ \\
\textit{Underlying complex}
 & Moore $M\check{C}^\bullet(\mathcal{U};\cA)$
 & Bar $\barBgeom(\cA)$ on $\mathrm{Ran}(X)$ \\
\textit{Combinatorial index}
 & Nerve of cover (simplicial trees)
 & Rooted trees / $\mathrm{FCom}$ \\
\textit{Degree-$n$ component}
 & $F_n$ (Eilenberg--Zilber)
 & $m_n$ (associahedra) \\
\textit{Degree-$3$ obstruction}
 & Jacobiator of $F_2$
 & $m_1 m_2 + m_2(m_2 \otimes 1 + 1 \otimes m_2)$ \\
\textit{Truncation depth}
 & $|I| - 1$ (\v{C}ech degree)
 & $\infty$ (no bound) \\
\textit{Genus $\geq 1$}
 & Higher \v{C}ech on modular covers
 & Stable graphs in $\mathrm{FCom}$ \\
\end{tabular}
\end{center}
At genus~$0$, both sides are governed by trees: the
\v{C}ech side by simplicial trees (nerve of the cover),
the bar side by rooted planar trees (bar-length
filtration). At genus $g \geq 1$, the bar side sees stable
graphs (the Feynman transform $\mathrm{FCom}$), while the
\v{C}ech side sees the cosimplicial structure of the modular
operad's boundary strata. The comparison
$\Phi_*$ intertwines these combinatorics via an $L_\infty$
morphism that is strict only when the cover has few enough
opens.
\end{remark}

\begin{remark}[Log-geometric path to the \v{C}ech--bar comparison]
\label{rem:cech-bar-log-path}
\index{Cech--bar comparison@\v{C}ech--bar comparison!log-geometric path}
The foundational triangle
(Malikov--Schechtman~\cite{MS24} +
Robert-Nicoud--Wierstra~\cite{RNW19}/Vallette~\cite{Val16}
+ Mok~\cite{Mok25}) gives a conjectural geometric
construction of the $L_\infty$ morphism $\Phi_*$ of
Conjecture~\ref{conj:cech-bar-intertwining}. On Mok's semistable
degeneration $W \to B$ \cite[Theorem~5.2.3]{Mok25}:
\begin{enumerate}[label=\textup{(\roman*)}]
\item Each smooth fibre $W_b$ ($b \neq b_0$) carries the bar
 resolution on its Ran space $\mathrm{Ran}(W_b)$, with MC element
 $\Theta^{(0)}_\cA = \sum m_n$.
\item The total space $W$ carries a \v{C}ech resolution with respect
 to the cover induced by the irreducible components of the special
 fibre $Y_1 \cup_D Y_2$, with MC element
 $\Phi_\cA = \sum F_n$ by \cite{MS24} applied fiberwise.
\item The degeneration formula \cite[Corollary~5.3.4]{Mok25}
 decomposes the special fibre of
 $\operatorname{FM}_n(W/B)$ into products of log-local
 FM spaces $\prod_v \operatorname{FM}_{I_v}(Y_v \mathbin{|} D_v)$.
 On each factor, the \v{C}ech and bar resolutions agree
 \textup{(}both reduce to affine data on a single component\textup{)}.
\item Gluing the log-local comparisons via the cutting map
 $\kappa$ \cite[Lemma~5.3.1]{Mok25} produces a candidate
 for~$\Phi_*$ that intertwines the two MC elements by
 construction: the defining property is
 $(\Phi_*)_*(\Theta^{(0)}_\cA) = \Phi_\cA$, i.e., the
 pushforward of the bar MC element under $\Phi_*$ equals the
 \v{C}ech MC element.
\end{enumerate}
Establishing that this candidate is an $L_\infty$
quasi-isomorphism reduces to the cohomological comparison
$\check{H}^\bullet(X,\cA) \cong H^{\mathrm{ch}}_\bullet(X,\cA)$
enhanced to the $L_\infty$ level by the one-slot functoriality
of $\operatorname{Conv}$~\cite[Theorem~5.1]{RNW19}.
\end{remark}

\subsection{Strictness and the cover}
\label{subsec:strictness-cover}
\index{two-element cover!strictness}

The depth of the \v{C}ech complex is controlled by the
combinatorics of the cover. For a cover by $|I|$ opens,
$\check{C}^p = 0$ for $p \geq |I|$, so the MC element
$\Phi_\cA$ can have at most $|I|$ nonzero degree components.
When $|I| = 2$ (the case of $\mathbb{P}^1$ with its standard
two-element cover, and more generally any curve equipped with a
decomposition into two affine opens), the MC element simplifies.

\begin{proposition}[Two-element covers are strict;
\ClaimStatusProvedHere]
\label{prop:cech-two-element-strict}
Let $\mathcal{U} = \{U_0, U_1\}$ be a two-element affine open
cover of~$X$. Then:
\begin{enumerate}[label=\textup{(\roman*)}]
\item The Moore complex $V = M\check{C}^\bullet(\mathcal{U};\cA)$
 is concentrated in degrees $0$ and $1$:
 $V^0 = \cA(U_0)$,\;
 $V^1 = \cA(U_0 \cap U_1)/\cA(U_1)$,\;
 $V^p = 0$ for $p \geq 2$.
\item For any degree $n \geq 3$, the Eilenberg--Zilber component
 $F_n\colon V^{\otimes n} \to V$ vanishes for degree reasons:
 since $F_n$ raises \v{C}ech degree by $n{-}1$ and $V^p = 0$
 for $p \geq 2$, the target $V^{m_1 + \cdots + m_n}$ is zero
 whenever $n \geq 3$ and $m_i \geq 0$.
\item The MC element reduces to
 $\Phi_\cA = (d_M, F_2, 0, 0, \ldots)$, and the master
 equation~\eqref{eq:cech-master-equation} collapses to
 two conditions:
 \begin{align}
 d_M^2 &= 0, \label{eq:two-element-dm} \\
 [d_M, F_2] &= 0,\quad
 \operatorname{Jac}(F_2) = 0.
 \label{eq:two-element-jacobi}
 \end{align}
 That is, $F_2$ is a \emph{strict} dg~Lie bracket: the Jacobi
 identity holds on the nose, not merely up to homotopy.
\item The comparison $L_\infty$ morphism
 $\Phi_*\colon \Phi_\cA \to \Theta^{(0)}_\cA$ is strict
 at the chain level (all higher Taylor components vanish),
 though the bar MC element $\Theta^{(0)}_\cA$ may still have
 $m_n \neq 0$ for $n \geq 3$; the non-strictness is
 absorbed entirely by the \emph{bar} side via homotopy transfer.
\end{enumerate}
\end{proposition}

\begin{proof}
Part~(i) is the standard Moore normalization for a two-element
cover. For part~(ii), observe that $F_n$ has \v{C}ech output
degree $m_1 + \cdots + m_n - (n-1)$ (it maps $n$ inputs to a
single output, combining $n$ \v{C}ech cochains via $n{-}1$
face-map contractions). When all inputs lie in $V^0 \oplus V^1$
and $n \geq 3$, at least $\lceil n/2\rceil \geq 2$ inputs must
be in degree~$0$ for the output to land in $V^0 \cup V^1$; but
the Eilenberg--Zilber diagonal forces the output degree to be
$\geq n - 1 \geq 2$, which lies in $V^{\geq 2} = 0$.

For part~(iii), with $F_n = 0$ for all $n \geq 3$, the MC
equation $D\,\Phi + \tfrac{1}{2}[\Phi,\Phi] = 0$ at degree~$3$
reads $0 + [F_2,[F_2,F_2]]/6 = 0$, which is the strict Jacobi
identity $\operatorname{Jac}(F_2) = 0$.

Part~(iv) follows because the $L_\infty$ morphism $\Phi_*$ is
constructed by homological perturbation of the \v{C}ech-to-bar
comparison. The perturbation series terminates at order~$1$ on
the \v{C}ech side (since $\Phi_\cA$ is strict), but the
transferred structure on the bar side acquires all $m_n$ via the
standard homological perturbation lemma.
\end{proof}

\begin{remark}[When does the Jacobiator appear?]
\label{rem:jacobiator-appearance}
\index{Jacobiator!cover dependence}
The vanishing of the Jacobiator for two-element covers is a
\emph{combinatorial} accident, not an algebraic one: the
underlying chiral algebra may be highly nonabelian (as in the
$\widehat{\mathfrak{sl}}_2$ example below), but the \v{C}ech
resolution is too short to see the ternary obstruction. Genuine
higher homotopies $F_n \neq 0$ for $n \geq 3$ appear precisely
when:
\begin{enumerate}[label=\textup{(\alph*)}]
\item The cover $\mathcal{U}$ has \emph{three or more opens} with
 nontrivial triple intersections, so that $\check{C}^2 \neq 0$
 and $F_3$ has room to land.
\item The genus is $g \geq 1$ and the cover is refined past two
 opens. Two affine opens \textup{(}the complements
 $X\setminus\{p\}$, $X\setminus\{q\}$ of two distinct points\textup{)}
 cover any smooth projective curve and keep the \v{C}ech side
 strict, so genus alone produces no higher operations: the
 curvature $\kappa(\cA)\cdot\omega_g$ is carried by the bar side.
 It becomes \v{C}ech-visible as a nonzero $F_3$ only on a cover
 with nontrivial triple overlaps.
\item The chiral algebra is nonabelian: for abelian $\cA$
 (e.g., the Heisenberg algebra), $F_2 = 0$ in the Lie sense
 (only the central extension survives), so
 $\operatorname{Jac}(F_2) = 0$ regardless of the cover.
\end{enumerate}
The sharpest case is an affine Kac--Moody algebra on an elliptic
curve with a three-element cover: here $F_3 \neq 0$ and reflects
the genus-$1$ curvature $\kappa(\cA)\cdot\omega_1$
(Example~\ref{ex:cech-hca-genus1}).
\end{remark}

\subsection{Examples}
\label{subsec:cech-hca-examples}

\begin{example}[\v{C}ech HCA for the Heisenberg algebra on
$\mathbb{P}^1$]
\label{ex:cech-hca-heisenberg}
\index{Heisenberg algebra!\v{C}ech HCA}
Take $X = \mathbb{P}^1$ with the standard cover
$U_0 = \mathbb{A}^1 = \operatorname{Spec}k[z]$ and
$U_\infty = \mathbb{P}^1 \setminus\{0\}
= \operatorname{Spec}k[z^{-1}]$,
so $U_0 \cap U_\infty = \mathbb{G}_m
= \operatorname{Spec}k[z,z^{-1}]$.
This is a two-element cover, so
Proposition~\ref{prop:cech-two-element-strict} applies. The \v{C}ech
complex for the Heisenberg sheaf $\cH_k$ is the two-term complex
\[
\check{C}^0 = \cH_k(U_0) \oplus \cH_k(U_\infty)
\xrightarrow{\;d_{\check{C}}\;}
\check{C}^1 = \cH_k(U_0 \cap U_\infty),
\]
where $d_{\check{C}}(f_0,f_\infty) = f_0|_{\mathbb{G}_m}
- f_\infty|_{\mathbb{G}_m}$.
Moore normalization gives $V^0 = \cH_k(U_0)$ and
$V^1 = \cH_k(\mathbb{G}_m)/\cH_k(U_\infty)$.

The MC element is
\begin{equation}\label{eq:heisenberg-cech-mc}
\Phi_{\cH_k} = (d_M, F_2, 0, 0, \ldots)
 \;\in\; \mathfrak{g}^{\check{C}}(\cH_k,\mathcal{U})^1,
\end{equation}
and the master equation reduces to $d_M^2 = 0$ (immediate) and
the strict Jacobi identity for~$F_2$. For the Heisenberg algebra,
$F_2$ is the transferred chiral bracket: for
$J_n = z^n\,dz \in \cH_k(U_0)$,
\[
F_2(J_m, J_n) = k\,m\,\delta_{m+n,0}\cdot\mathbf{1}.
\]
Since the Heisenberg bracket is abelian (only the central
extension is nonzero), the Jacobi identity $[F_2,[F_2,F_2]] = 0$
is automatic: the double commutator vanishes by centrality.

The comparison morphism $\Phi_*\colon \Phi_{\cH_k} \to
\Theta^{(0)}_{\cH_k}$ is strict because \emph{both} sides are
strict: the bar $A_\infty$ structure has $m_n = 0$ for $n \geq 3$
(\S\ref{sec:explicit-low-degree}), matching the \v{C}ech
strictness.

The \v{C}ech cohomology
$H^0(\mathbb{P}^1,\cH_k) = k\cdot\mathbf{1}$ (global constants) and
$H^1(\mathbb{P}^1,\cH_k) = \bigoplus_{n<0} k\cdot J_n$
(negative modes modulo positive modes) reproduce
the chiral homology of~$\cH_k$ on~$\mathbb{P}^1$.
\end{example}

\begin{example}[\v{C}ech HCA for $\widehat{\mathfrak{sl}}_2$ on
$\mathbb{P}^1$: strict \v{C}ech, non-strict bar]
\label{ex:cech-hca-sl2}
\index{sl2hat@$\widehat{\mathfrak{sl}}_2$!\v{C}ech HCA}
For $\cA = \widehat{\mathfrak{sl}}_2$ at level~$k$ on
$\mathbb{P}^1$ with the same two-element cover, the \v{C}ech
MC element is again strict:
\[
\Phi_{\widehat{\mathfrak{sl}}_2}
 = (d_M, F_2, 0, 0, \ldots),
\]
by Proposition~\ref{prop:cech-two-element-strict}. The binary bracket
$F_2$ encodes the full affine OPE on the overlap
$\mathbb{G}_m$, and the Jacobi identity holds \emph{on the nose}
for the transferred bracket on the Moore complex, not because the
algebra is abelian (it is not), but because the \v{C}ech complex
is too short for the Jacobiator to have a nonzero target.

The bar side is \emph{not} strict. The
bar $A_\infty$ structure on $\barBgeom(\widehat{\mathfrak{sl}}_2)$
has $m_3 \neq 0$: the ternary operation is nonzero and is
responsible for the nontrivial $\bar{H}^2$ (which equals~$5$ for
$\mathfrak{sl}_2$; cf.\ \S\ref{sec:non-renormalization-tree}).
The comparison $L_\infty$ morphism $\Phi_*$ is therefore
non-strict at order~$2$:
\[
(\Phi_*)_2(F_2 \otimes F_2) = m_3 \neq 0.
\]
That is, $m_3$ on the bar side is \emph{created} from $F_2$ on the
\v{C}ech side by homotopy transfer: the standard homological
perturbation lemma applied to the \v{C}ech-to-bar comparison
produces $m_3 = h \circ F_2(F_2 \otimes \operatorname{id}) +
\cdots$, where $h$ is the contracting homotopy of the comparison
quasi-isomorphism.

This example cleanly separates the two notions of strictness:
the \v{C}ech resolution is strict (a property of the cover), while
the bar resolution is non-strict (a property of the algebra).
The $L_\infty$ morphism bridges the two.
\end{example}

\begin{example}[\v{C}ech HCA at genus~$1$: the Jacobiator and
curvature]
\label{ex:cech-hca-genus1}
\index{elliptic curve!\v{C}ech HCA}
\index{Jacobiator!genus-1 appearance}
Let $E_\tau$ be an elliptic curve (genus~$1$), and let
$\mathcal{U} = \{U_0, U_1, U_2\}$ be a three-element affine open
cover chosen with nonempty triple overlap \textup{(}two affine opens
already cover $E_\tau$ and compute $H^1(E_\tau,\cO)=k$, but on a
two-element cover the \v{C}ech side is strict by
Proposition~\ref{prop:cech-two-element-strict}; the third open makes
$\check{C}^2 \neq 0$, giving the Jacobiator a nonzero target\textup{)}.
The \v{C}ech complex for a chiral
algebra $\cA$ on $E_\tau$ now extends to degree~$2$:
\[
\check{C}^0 \xrightarrow{\;d_{\check{C}}\;} \check{C}^1
\xrightarrow{\;d_{\check{C}}\;} \check{C}^2,
\]
where $\check{C}^2 = \Gamma(U_0 \cap U_1 \cap U_2,\;\cA)$ is the
sections on the (generically nonempty) triple overlap. The Moore
complex has $V^2 \neq 0$, so
Proposition~\ref{prop:cech-two-element-strict} does \emph{not} apply.

The MC element now has a genuine ternary component:
\[
\Phi_\cA = (d_M, F_2, F_3, 0, 0, \ldots).
\]
The truncation at $F_3$ follows because $V^p = 0$ for $p \geq 3$
(a three-element cover has \v{C}ech degree at most~$2$). The
master equation at degree~$3$
(equation~\eqref{eq:cech-degree-3}) reads
\[
[d_M, F_3] + \operatorname{Jac}(F_2) = 0,
\]
asserting that $F_3$ is the explicit nullhomotopy of the Jacobiator
of the binary bracket.

Under the \v{C}ech--bar comparison, $F_3$ maps to the bar-side
curvature. Specifically, for
$\cA = \widehat{\mathfrak{g}}_k$ (affine Kac--Moody at level~$k$),
the class $[F_3] \in H^2(\mathfrak{g}^{\check{C}})$ is the descent
shadow of the bar curvature:
\[
[\Phi_*\!(F_3)] = \kappa(\widehat{\mathfrak{g}}_k)\cdot\omega_1
 \;\in\; H^1(E_\tau,\;\omega_{E_\tau}),
\]
where $\kappa(\widehat{\mathfrak{g}}_k) =
\dim\mathfrak{g}\,(k + h^\vee) / (2h^\vee)$ is the modular
characteristic
(Theorem~\ref{thm:genus-universality}) and
$\omega_1 = dz$ is the period form on $E_\tau$. This is the
first appearance of the genus-$1$ obstruction in \v{C}ech language:
the bar curvature $m_0^{(1)} = \kappa \cdot \omega_1$ on
$\barBgeom(\cA)|_{E_\tau}$, which records the failure of $d_B^2 = 0$
at genus~$1$
(cf.\ Chapter~\ref{chap:higher-genus},
\S\ref{sec:genus-1-elliptic-complete}), is \emph{witnessed} on the
\v{C}ech side by the nonvanishing of $F_3$.

The pattern generalizes: on a genus-$g$ curve with a $(g{+}1)$-element
cover, the MC element has
$\Phi_\cA = (d_M, F_2, F_3, \ldots, F_{g+1}, 0, 0, \ldots)$,
and the degree-$(g{+}1)$ component $F_{g+1}$ descends from the
genus-$g$ bar curvature $\kappa(\cA) \cdot \omega_g$. The
\v{C}ech resolution thus provides an alternative, explicitly
combinatorial window into the modular tower of obstructions
$\Theta_\cA^{\leq r}$.
\end{example}

%% ===================================================================
%% Casimir invariants and the divisor-core lattice
%% ===================================================================

\section{Casimir invariants and the divisor-core lattice}
\label{sec:casimir-divisor-core}
\index{Casimir invariants!divisor-core lattice}
\index{divisor-core lattice}

The bar cohomology $H^\bullet(\barBch(\cA))$ carries a natural
action of the Casimir algebra~$Z_\cA$, the center of the
universal enveloping algebra, generated by the rank-$r$ Casimir
operators $C_1,\ldots,C_r$ together with the level
operator~$K$. This action is visible already in the explicit
tables of Chapter~\ref{chap:detailed-computations}. On the
finite-character recurrence loci isolated below, the Casimir character
sequence generates a cyclic $k[\sigma]$-module with minimal polynomial
$p_\cA(t)$; when $p_\cA(t)=t^{d_\cA}\Delta_\cA(t^{-1})$, the
spectral discriminant records the characteristic polynomial of this
module. The eigenvalue structure of the Casimir action organizes the
bar cohomology into a \emph{divisor lattice}, the same lattice that
classifies submodules of cyclic $k[t]$-modules. Transport between
algebras related by exact functors (Drinfeld--Sokolov reduction,
quantum Hamiltonian reduction, level-rank duality) factors through
the greatest common divisor of their recurrence polynomials exactly
when the functor induces a shift-compatible map of the corresponding
recurrence modules. This is where the algebraic formalism of Koszul
duality makes contact with the classical Jordan normal form.

\medskip

The section is organized as follows. We first develop the pure
algebra of cyclic transport objects and divisor cores
(\S\ref{subsec:cyclic-transport-objects}--\ref{subsec:hom-gcd-factorization}),
then apply the theory to bar cohomology
(\S\ref{subsec:casimir-recurrence-module}--\ref{subsec:sector-determinant}),
and close with the common-core transport theorem and its
implications for the Drinfeld--Sokolov correspondence
(\S\ref{subsec:common-core-transport}--\ref{subsec:casimir-frontier}).

\subsection{Cyclic transport objects and divisor cores}
\label{subsec:cyclic-transport-objects}

Throughout this section, $k$ denotes a field of
characteristic~$0$.

\begin{definition}[Universal cyclic transport object]
\label{def:cyclic-transport-object}
\index{cyclic transport object}
Let $p(t) \in k[t]$ be monic of degree~$d$. The \emph{cyclic
transport object} attached to~$p$ is the $k[t]$-module
\[
M(p) \;:=\; k[t]/(p).
\]
Multiplication by~$t$ defines the shift endomorphism
$T_p \colon M(p) \to M(p)$, $T_p([f]) = [tf]$. The module
$M(p)$ satisfies the evident universal property: every
finite-dimensional cyclic $k[t]$-module annihilated by~$p$ is
a quotient of~$M(p)$, and the quotient is an isomorphism
precisely when $p$~is the minimal polynomial.
\end{definition}

\begin{definition}[Divisor core]
\label{def:divisor-core}
\index{divisor core}
Let $p(t) \in k[t]$ be monic and $g \mid p$ a monic divisor.
The \emph{divisor core} associated to~$g$ is the submodule
\[
\mathcal{C}_g(p)
 \;:=\;
 (p/g)\,M(p)
 \;\subset\; M(p).
\]
\end{definition}

The divisor core is the canonical embedding of the
$g$-factor into the ambient module. The key theorem is that
divisor cores account for \emph{all} submodules, and that
their lattice structure mirrors the divisor lattice of~$p$.

\begin{theorem}[Divisor-core calculus; \ClaimStatusProvedHere]
\label{thm:divisor-core-calculus-inv}
\index{divisor-core calculus}
Let $p \in k[t]$ be monic, and let $g,g_1,g_2 \mid p$ be monic
divisors.
\begin{enumerate}[label=\textup{(\roman*)}]
\item \textup{(Canonical identification.)} There is a canonical
 isomorphism of $k[t]$-modules
 \[
 \mathcal{C}_g(p) \;\cong\; k[t]/(g).
 \]
\item \textup{(Short exact sequence.)} The inclusion
 $\mathcal{C}_g(p) \hookrightarrow M(p)$ fits into a canonical
 short exact sequence
 \[
 0 \longrightarrow \mathcal{C}_g(p)
 \longrightarrow M(p)
 \longrightarrow k[t]/(p/g)
 \longrightarrow 0.
 \]
\item \textup{(Lattice identities.)} The divisor cores satisfy
 \begin{align*}
 \mathcal{C}_{g_1}(p) \cap \mathcal{C}_{g_2}(p)
 &= \mathcal{C}_{\gcd(g_1,g_2)}(p), \\
 \mathcal{C}_{g_1}(p) + \mathcal{C}_{g_2}(p)
 &= \mathcal{C}_{\operatorname{lcm}(g_1,g_2)}(p).
 \end{align*}
\item \textup{(Idempotent splitting criterion.)} The following
 are equivalent:
 \begin{enumerate}[label=\textup{(\alph*)}]
 \item $\mathcal{C}_g(p)$ is the image of an idempotent
 polynomial in~$T_p$;
 \item $M(p) \cong \mathcal{C}_g(p) \oplus k[t]/(p/g)$;
 \item $\gcd(g, p/g) = 1$.
 \end{enumerate}
\end{enumerate}
\end{theorem}

\begin{proof}
For~(i), the map $k[t] \to \mathcal{C}_g(p)$ defined by
$f \mapsto f \cdot (p/g) \bmod p$ descends to a surjection
$k[t]/(g) \twoheadrightarrow \mathcal{C}_g(p)$. If
$f \cdot (p/g) \equiv 0 \pmod{p}$, then $p \mid f(p/g)$;
since $p = g \cdot (p/g)$ and $k[t]$ is a domain,
$g \mid f$. Hence the kernel is~$(g)$ and the map is an
isomorphism.

For~(ii), the quotient $M(p) \twoheadrightarrow k[t]/(p/g)$
has kernel $\{[f] : (p/g) \mid f \bmod p\} = \mathcal{C}_g(p)$.

For~(iii), the intersection and sum of ideals in a PID yield
$\operatorname{lcm}$ and $\gcd$ respectively; the stated
lattice identities follow by the correspondence
$\mathcal{C}_g(p) = (p/g)/(p)$.

For~(iv), a direct-sum splitting of~(ii) exists if and only if
the Chinese remainder theorem applies, which holds precisely
when $(g)$ and~$(p/g)$ are comaximal, i.e.\ $\gcd(g,p/g)=1$.
\end{proof}

\begin{corollary}[Divisors classify submodules; \ClaimStatusProvedHere]
\label{cor:divisors-classify-submodules-inv}
Every submodule of $M(p)$ is of the form $\mathcal{C}_g(p)$
for a unique monic divisor $g \mid p$. The lattice of
submodules of $M(p)$ is isomorphic to the divisor lattice
of~$p$.
\end{corollary}

\begin{proof}
Submodules of $k[t]/(p)$ correspond to ideals $(a)$ with
$(p) \subseteq (a) \subseteq k[t]$. Since $k[t]$ is a PID,
$a$~is a monic divisor of~$p$. Setting $g = p/a$ gives
$\mathcal{C}_g(p)$.
Uniqueness of $g$ is equivalent to uniqueness of the intermediate ideal
$(a)$.
The lattice isomorphism is immediate from
Theorem~\ref{thm:divisor-core-calculus-inv}(iii): intersection and
sum of submodules correspond respectively to $\gcd$ and
$\operatorname{lcm}$ of divisors.
\end{proof}

\begin{remark}[The first structural correction]
\label{rem:first-structural-correction}
A factor $g \mid p$ always defines a canonical submodule of
$M(p)$. It defines a canonical \emph{projector} only when
$\gcd(g,p/g) = 1$. When this coprimality fails (that is,
when the divisor shares roots with its complement), the
submodule is canonically embedded but not idempotently split.
The extension class $\operatorname{Ext}^1(k[t]/(p/g),\,
\mathcal{C}_g(p))$ that prevents the splitting is the
algebraic shadow of a genuine Jordan block.
\end{remark}

\subsection{Hom~\texorpdfstring{$=$}{=}~gcd and factorization
 through common cores}
\label{subsec:hom-gcd-factorization}

\begin{theorem}[\texorpdfstring{$\operatorname{Hom} =
 \gcd$}{Hom = gcd}; \ClaimStatusProvedHere]
\label{thm:hom-equals-gcd-inv}
\index{Hom = gcd}
Let $p,q \in k[t]$ be monic, and set $g = \gcd(p,q)$. Then
\[
\operatorname{Hom}_{k[t]}\!\bigl(M(p),\,M(q)\bigr)
 \;\cong\; k[t]/(g).
\]
In particular,
$\dim_k\operatorname{Hom}_{k[t]}(M(p),M(q)) = \deg g$.
\end{theorem}

\begin{proof}
A $k[t]$-linear map $f \colon M(p) \to M(q)$ is determined by
$x := f(1) \in M(q)$, subject to $p \cdot x = 0$ in~$M(q)$,
i.e.\ $q \mid px$. Write $p = gp'$, $q = gq'$ with
$\gcd(p',q') = 1$. Then $q \mid px$ becomes $q' \mid p'x$,
and coprimality forces $q' \mid x$. Hence the space of
allowed images is $(q')/(gq') \cong k[t]/(g)$.
\end{proof}

\begin{theorem}[Universal factorization through the common
 core; \ClaimStatusProvedHere]
\label{thm:factorization-through-common-core-inv}
\index{common core!factorization}
Let $p,q \in k[t]$ be monic, $g = \gcd(p,q)$, and write
$\pi_p \colon M(p) \twoheadrightarrow M(g)$ for the canonical
quotient and $\iota_q \colon M(g) \hookrightarrow M(q)$ for
the inclusion onto the divisor core $\mathcal{C}_g(q)$. Then
every $k[t]$-linear map $f \colon M(p) \to M(q)$ factors
uniquely as
\[
f \;=\; \iota_q \circ m_u \circ \pi_p
\]
for a unique class $u \in M(g)$, where
$m_u \colon M(g) \to M(g)$ denotes multiplication by~$u$.
\end{theorem}

\begin{proof}
By Theorem~\ref{thm:hom-equals-gcd-inv},
$\operatorname{Hom}_{k[t]}(M(p),M(q)) \cong M(g)$. The map
$u \mapsto \iota_q \circ m_u \circ \pi_p$ is $k$-linear,
injective (since $\pi_p$ is surjective and $\iota_q$ is
injective), and both source and target have dimension
$\deg g$; hence it is an isomorphism.
\end{proof}

\begin{remark}
The theorem says that transport maps between cyclic spectral
objects are \emph{never mysterious}: they are exactly
multiplication operators on the maximal common core
$M(\gcd(p,q))$. The common divisor controls everything.
\end{remark}

\subsection{The Casimir recurrence module}
\label{subsec:casimir-recurrence-module}
\index{recurrence module}

The abstract divisor-core theory connects to bar cohomology as
follows.

\begin{hypothesis}[Finite-character recurrence]
\label{hyp:finite-character-recurrence-inv}
Let $\cA$ be a modular Koszul chiral algebra in the
Lie-theoretic world, and write $Z_\cA = k[C_1,\ldots,C_r,K]$
for the commutative Casimir algebra. We assume:
\begin{enumerate}[label=\textup{(\roman*)}]
\item Each $H^n(\barBch(\cA))$ has finite generalized
 $Z_\cA$-support.
\item The class sequence
 $u_\cA(n) := [H^n(\barBch(\cA))] \in K_0(Z_\cA\text{-mod}_{\mathrm{fl}})
 \otimes_{\mathbb{Z}} k$
 satisfies a minimal monic shift recurrence
 $p_\cA(\sigma)\,u_\cA = 0$
 for some monic $p_\cA(t) \in k[t]$.
\item The recurrence polynomial is the reversal of the
 spectral discriminant:
 $p_\cA(t) = t^{d_\cA}\Delta_\cA(t^{-1})$,
 where $d_\cA = \deg\Delta_\cA$.
\end{enumerate}
\end{hypothesis}

\noindent
This hypothesis holds for all standard families: the
affine Kac--Moody algebras
(\S\ref{sec:non-renormalization-tree}),
the $\mathcal{W}$-algebras
(Theorem~\ref{thm:ds-bar-gf-discriminant}),
and the Virasoro algebra
(Theorem~\ref{thm:discriminant-spectral}).

\begin{definition}[Casimir recurrence module and
 Kodaira--Spencer operator]
\label{def:casimir-recurrence-module}
\index{Kodaira--Spencer operator!Casimir}
Under Hypothesis~\ref{hyp:finite-character-recurrence-inv},
define the \emph{Casimir recurrence module}
\[
\mathcal{R}_\cA
 \;:=\;
 k[\sigma]\cdot u_\cA
 \;\subset\;
 \bigl(K_0(Z_\cA\text{-mod}_{\mathrm{fl}})
 \otimes k\bigr)^{\mathbb{N}},
\]
and the \emph{algebraic reduced first cohomology}
$H^{1,\mathrm{alg}}_{\mathrm{red}}(\cA)
 := \mathcal{R}_\cA^\vee$.
The \emph{Kodaira--Spencer operator} is the dual of the
shift:
\[
KS_\cA \;:=\; \sigma^\vee
 \;\in\;
 \operatorname{End}\!\bigl(
 H^{1,\mathrm{alg}}_{\mathrm{red}}(\cA)
 \bigr).
\]
\end{definition}

\begin{theorem}[Minimal intrinsic realization;
\ClaimStatusProvedHere]
\label{thm:minimal-intrinsic-realization-inv}
Under Hypothesis~\ref{hyp:finite-character-recurrence-inv},
$\mathcal{R}_\cA \cong k[t]/(p_\cA)$ as $k[t]$-modules.
Consequently
$\dim_k H^{1,\mathrm{alg}}_{\mathrm{red}}(\cA)
 = d_\cA$,
and the Kodaira--Spencer operator recovers the spectral
discriminant:
\begin{equation}\label{eq:ks-recovers-discriminant}
\det(1 - x\,KS_\cA)
 \;=\;
 \Delta_\cA(x).
\end{equation}
\end{theorem}

\begin{proof}
The $k[t]$-linear surjection
$k[t] \twoheadrightarrow \mathcal{R}_\cA$ sending
$f(t) \mapsto f(\sigma)\,u_\cA$ has kernel $(p_\cA)$ by
minimality of the recurrence, so
$\mathcal{R}_\cA \cong M(p_\cA)$. The characteristic
polynomial of~$T_{p_\cA}$ is~$p_\cA$, so
$\det(1-x\,KS_\cA)
 = x^{d_\cA}\,p_\cA(x^{-1})
 = \Delta_\cA(x)$.
\end{proof}

\subsection{Sector determinant and Casimir reconstruction}
\label{subsec:sector-determinant}
\index{sector determinant}

Every divisor $S(x) \mid \Delta_\cA(x)$ of the spectral
discriminant produces a canonical finite-dimensional sector
inside the Casimir recurrence module.

\begin{definition}[Reduced spectral sector]
\label{def:reduced-spectral-sector}
Let $S(x) \mid \Delta_\cA(x)$ be a monic divisor, and set
$s(t) := t^{\deg S}\,S(t^{-1})$. The \emph{spectral sector
of type~$S$} is the canonical subquotient
\[
H^{1,\mathrm{alg}}_S(\cA)
 \;:=\;
 \bigl(\mathcal{C}_s(p_\cA)\bigr)^\vee.
\]
\end{definition}

\begin{theorem}[Sector determinant; \ClaimStatusProvedHere]
\label{thm:sector-determinant-inv}
For every divisor $S \mid \Delta_\cA$, the dual of the
$S$-sector satisfies
$(\,H^{1,\mathrm{alg}}_S(\cA)\,)^\vee
 \cong k[t]/(s)$,
and the induced endomorphism $KS_{\cA,S}$ satisfies
\[
\det(1 - x\,KS_{\cA,S}) \;=\; S(x).
\]
\end{theorem}

\begin{proof}
By Theorem~\ref{thm:divisor-core-calculus-inv}(i),
$\mathcal{C}_s(p_\cA) \cong k[t]/(s)$; the determinant
identity follows exactly as in
Theorem~\ref{thm:minimal-intrinsic-realization-inv}.
Namely, the endomorphism on the $S$-sector is multiplication by~$t$ on
$k[t]/(s)$, whose characteristic polynomial is~$s$.
Reversing the polynomial back to the $x$-variable gives $S(x)$.
\end{proof}

\begin{theorem}[Casimir moment reconstruction;
\ClaimStatusProvedHere]
\label{thm:casimir-moment-reconstruction-inv}
\index{Casimir moment reconstruction}
Suppose the generalized $Z_\cA$-spectrum occurring in
$H^n(\barBch(\cA))$ is eventually contained in a finite set
$\Xi_\cA = \{\chi_1,\ldots,\chi_s\}$. Choose
$Q \in Z_\cA$ separating~$\Xi_\cA$ \textup{(}a generic
linear combination of Casimirs suffices over an infinite
field\textup{)}, and write $q_i = \chi_i(Q)$. Then the
scalar moment sequences
\[
\tau_m(n) \;:=\;
 \operatorname{Tr}\!\bigl(
 Q^m \mid H^n(\barBch(\cA))
 \bigr),
\qquad 0 \leq m \leq s-1,
\]
determine the multiplicity vector of generalized characters
at each degree~$n$ via the Vandermonde inversion
$m(n) = V^{-1}\tau(n)$. In particular, the full recurrence
module $\mathcal{R}_\cA$ is recoverable from finitely many
scalar Casimir traces.
\end{theorem}

\begin{proof}
On the generalized $\chi_i$-block, $Q$ acts as
$q_i\,I + N_i$ with $N_i$ nilpotent, so
$\operatorname{Tr}(Q^m \mid V_{\chi_i})
 = q_i^m\dim V_{\chi_i}$.
The moment vector $\tau(n)$ is related to the multiplicity
vector $m(n)$ by the Vandermonde matrix
$(q_i^m)_{m,i}$, which is invertible because the~$q_i$ are
distinct.
\end{proof}

\subsection{Primary Jordan filtration}
\label{subsec:primary-jordan}
\index{Jordan filtration!primary}

The divisor-core lattice refines to a Jordan filtration at
each irreducible factor of the recurrence polynomial. This
is where repeated roots of the discriminant reveal their
geometric content.

\begin{theorem}[Primary Jordan filtration;
\ClaimStatusProvedHere]
\label{thm:primary-jordan-filtration-inv}
Let $q$ be an irreducible factor of $p_\cA$ and set
$n = \nu_q(p_\cA)$. Then the primary divisor cores
$\mathcal{J}_q^m(p_\cA) := \mathcal{C}_{q^m}(p_\cA)$,
$1 \leq m \leq n$, form a canonical chain
\[
0 \;\subset\;
 \mathcal{J}_q^1(p_\cA) \;\subset\;
 \mathcal{J}_q^2(p_\cA) \;\subset\;
 \cdots \;\subset\;
 \mathcal{J}_q^n(p_\cA)
 \;\subset\; M(p_\cA),
\]
with successive quotients
$\mathcal{J}_q^m / \mathcal{J}_q^{m-1} \cong k[t]/(q)$.
If $q = t - \lambda$ splits linearly over~$k$, then
$T_{p_\cA} - \lambda$ acts on $\mathcal{J}_{t-\lambda}^m$
by a single Jordan block of size~$m$.
\end{theorem}

\begin{proof}
By Theorem~\ref{thm:divisor-core-calculus-inv}(i),
$\mathcal{J}_q^m(p_\cA) \cong k[t]/(q^m)$. The successive
quotient $k[t]/(q^m) \big/ q^{m-1}k[t]/(q^m) \cong k[t]/(q)$.
For $q = t-\lambda$, set $\tau = t-\lambda$; then $t$ acts as
$\lambda + \tau$ on $k[\tau]/(\tau^m)$, and $\tau$ is the
standard nilpotent Jordan shift.
\end{proof}

\begin{corollary}[Repeated roots are extension classes;
\ClaimStatusProvedHere]
\label{cor:repeated-roots-extension-data}
Let $p \in k[t]$ be monic, let $g \mid p$, set $h=p/g$, and set
$d_g=\gcd(g,h)$. The canonical sequence
\[
0 \longrightarrow \mathcal C_g(p)
\longrightarrow M(p)
\longrightarrow k[t]/(h)
\longrightarrow 0
\]
defines an obstruction class
\[
\varepsilon_g(p)\in
\operatorname{Ext}^1_{k[t]}\!\bigl(k[t]/(h),k[t]/(g)\bigr)
\cong k[t]/(d_g).
\]
Under the standard resolution
\[
0 \longrightarrow k[t]\xrightarrow{\;h\;}k[t]
\longrightarrow k[t]/(h)\longrightarrow 0,
\]
the class $\varepsilon_g(p)$ is represented by $1\bmod d_g$.
Thus $\varepsilon_g(p)=0$ if and only if $\gcd(g,h)=1$.
If $d_g$ is nonconstant, the divisor core is canonical but the
extension is not split.
\end{corollary}

\begin{proof}
Identify $\mathcal C_g(p)$ with $k[t]/(g)$ by
Theorem~\ref{thm:divisor-core-calculus-inv}(i). The middle term
$M(p)=k[t]/(gh)$ is generated by $e=1\bmod gh$; the submodule is
generated by $n=h\bmod gh$. The relations are
\[
g n=0,\qquad h e=n.
\]
The displayed resolution computes
\[
\operatorname{Ext}^1_{k[t]}(k[t]/(h),k[t]/(g))
\cong
\operatorname{coker}\!\bigl(h:k[t]/(g)\to k[t]/(g)\bigr)
\cong k[t]/(g,h).
\]
In the presentation above, the relation $h e=n$ is exactly the class
$1\bmod (g,h)$. This class vanishes precisely when $d_g=1$,
equivalently when the Chinese-remainder splitting of
Theorem~\ref{thm:divisor-core-calculus-inv}(iv) exists.
\end{proof}

\subsection{Common-core transport between algebras}
\label{subsec:common-core-transport}
\index{common core!transport}

The common-core theorem is the structural heart of the
section: it controls what exact functors between chiral
algebras can see at the level of the spectral discriminant.

Let $\cA$ and~$\mathcal{B}$ both satisfy
Hypothesis~\ref{hyp:finite-character-recurrence-inv}, with
recurrence polynomials $p_\cA$ and~$p_\mathcal{B}$. Write
$g_{\cA,\mathcal{B}} = \gcd(p_\cA, p_\mathcal{B})$.

\begin{definition}[Maximal common transport core]
\label{def:maximal-common-core}
\index{common core!maximal}
The \emph{maximal common transport core} is the module
$\mathcal{C}^{\max}_{\cA,\mathcal{B}}
 := M(g_{\cA,\mathcal{B}})
 = k[t]/(g_{\cA,\mathcal{B}})$.
\end{definition}

\begin{theorem}[Common-core exact sequences;
\ClaimStatusProvedHere]
\label{thm:common-core-exact-sequences-inv}
There are canonical short exact sequences
\begin{align*}
0 \to \mathcal{C}^{\max}_{\cA,\mathcal{B}}
 &\longrightarrow \mathcal{R}_\cA
 \longrightarrow k[t]/(p_\cA/g_{\cA,\mathcal{B}})
 \longrightarrow 0, \\[4pt]
0 \to \mathcal{C}^{\max}_{\cA,\mathcal{B}}
 &\longrightarrow \mathcal{R}_\mathcal{B}
 \longrightarrow k[t]/(p_\mathcal{B}/g_{\cA,\mathcal{B}})
 \longrightarrow 0.
\end{align*}
The quotient modules
$k[t]/(p_\cA/g_{\cA,\mathcal{B}})$ and
$k[t]/(p_\mathcal{B}/g_{\cA,\mathcal{B}})$
are the \emph{spectral defects}: they measure what each
algebra carries that the other cannot see.
\end{theorem}

\begin{proof}
Apply Theorem~\ref{thm:divisor-core-calculus-inv}(ii) with
$g = g_{\cA,\mathcal{B}}$ to each of the recurrence modules
$\mathcal{R}_\cA \cong M(p_\cA)$ and
$\mathcal{R}_\mathcal{B} \cong M(p_\mathcal{B})$.
This identifies the common core with the divisor core cut out by
$g_{\cA,\mathcal{B}}$ on both sides.
The two quotient modules are therefore exactly the complementary
defect factors recorded in the statement.
\end{proof}

\begin{proposition}[Transport factors through the common core;
\ClaimStatusProvedHere]
\label{prop:transport-factors-inv}
Every $k[t]$-linear map
$F \colon \mathcal{R}_\cA \to \mathcal{R}_\mathcal{B}$
factors uniquely as
\[
\mathcal{R}_\cA
 \twoheadrightarrow
 \mathcal{C}^{\max}_{\cA,\mathcal{B}}
 \xrightarrow{\;m_u\;}
 \mathcal{C}^{\max}_{\cA,\mathcal{B}}
 \hookrightarrow
 \mathcal{R}_\mathcal{B}
\]
for a unique $u \in \mathcal{C}^{\max}_{\cA,\mathcal{B}}$.
Equivalently, every shift-compatible transport between the
algebraic reduced cohomologies is governed by multiplication
on the maximal common core.
\end{proposition}

\begin{proof}
Theorem~\ref{thm:minimal-intrinsic-realization-inv} identifies each
recurrence module with its minimal polynomial model
$M(p_\cA)$ or $M(p_\mathcal{B})$.
Theorem~\ref{thm:factorization-through-common-core-inv} then says that
every $k[t]$-linear map between these models factors uniquely through
the common core $M(g_{\cA,\mathcal{B}})$ by multiplication by an
element of that core.
Transporting this factorization back through the intrinsic
realizations gives exactly the stated factorization for
$\mathcal{R}_\cA$ and $\mathcal{R}_\mathcal{B}$.
\end{proof}

\subsection{The \texorpdfstring{$\widehat{\mathfrak{sl}}_3$
 versus $W_3$}{sl3-hat vs W3} common core}
\label{subsec:sl3-w3-example}
\index{sl3hat versus W3@$\widehat{\mathfrak{sl}}_3$ versus $W_3$!common core}

The first nontrivial example is the Drinfeld--Sokolov pair
$(\widehat{\mathfrak{sl}}_3,\,W_3)$, where the common core
is strictly smaller than either recurrence module. For the
rank-one pair $(\widehat{\mathfrak{sl}}_2,\,
\mathrm{Vir})$, both algebras share
$\Delta(x) = (1-3x)(1+x)$
(cf.~Theorem~\ref{thm:discriminant-spectral}), so the
common core is the entire transport module and the theory
degenerates. The first rank where nontrivial defects appear
is rank~$2$.

Assume the finite-character denominator chart recorded in
Theorem~\ref{thm:ds-bar-gf-discriminant} supplies
\begin{align*}
\Delta_{\widehat{\mathfrak{sl}}_3}(x)
 &= (1-8x)(1-3x-x^2), \\
\Delta_{W_3}(x)
 &= (1-x)(1-3x-x^2),
\end{align*}
so that the reversed recurrence polynomials are
\begin{align*}
p_{\widehat{\mathfrak{sl}}_3}(t)
 &= (t-8)(t^2-3t-1), \\
p_{W_3}(t)
 &= (t-1)(t^2-3t-1).
\end{align*}
The maximal common core is the two-dimensional module
\[
\mathcal{C}^{\max}_{\widehat{\mathfrak{sl}}_3,\,W_3}
 \;\cong\;
 k[t]/(t^2-3t-1),
\]
with eigenvalues
$(3 \pm \sqrt{13})/2$.

\begin{proposition}[Exact defect decomposition;
\ClaimStatusProvedHere]
\label{prop:sl3-w3-defect-inv}
The common-core exact sequences of
Theorem~\ref{thm:common-core-exact-sequences-inv} take the
form
\begin{align*}
0 \to k[t]/(t^2{-}3t{-}1)
 &\longrightarrow k[t]/\bigl((t{-}8)(t^2{-}3t{-}1)\bigr)
 \longrightarrow k[t]/(t{-}8)
 \longrightarrow 0, \\[2pt]
0 \to k[t]/(t^2{-}3t{-}1)
 &\longrightarrow k[t]/\bigl((t{-}1)(t^2{-}3t{-}1)\bigr)
 \longrightarrow k[t]/(t{-}1)
 \longrightarrow 0.
\end{align*}
The one-dimensional defect modules $k[t]/(t-8)$ and
$k[t]/(t-1)$ record the non-preserved dominant modes:
eigenvalue~$8$ for
$\widehat{\mathfrak{sl}}_3$ and eigenvalue~$1$ for~$W_3$.
\end{proposition}

\begin{proof}
Theorem~\ref{thm:divisor-core-calculus-inv}(ii) gives the short exact
sequences
\[
0\to M(g_{\cA,\mathcal{B}})\to M(p_\cA)\to
M(p_\cA/g_{\cA,\mathcal{B}})\to 0
\]
and the analogous sequence for $M(p_\mathcal{B})$.
Under the explicit polynomials in this example,
$g_{\cA,\mathcal{B}}=t^2-3t-1$,
$p_{\widehat{\mathfrak{sl}}_3}/g_{\cA,\mathcal{B}}=t-8$, and
$p_{W_3}/g_{\cA,\mathcal{B}}=t-1$.
Substituting these values into the general exact sequences gives the
two displayed decompositions.
The quotient factors are one-dimensional, so they record exactly the
dominant modes that are not seen by the maximal common core.
\end{proof}

\begin{proposition}[Explicit coprime-locus projectors;
\ClaimStatusProvedHere]
\label{prop:sl3-w3-projectors}
Since $\gcd(t^2-3t-1,\, t-8) = \gcd(t^2-3t-1,\, t-1) = 1$,
the common quadratic sector splits idempotently on both sides.
\begin{enumerate}[label=\textup{(\roman*)}]
\item For $\widehat{\mathfrak{sl}}_3$: the B\'{e}zout identity
 \[
 \tfrac{1}{39}(t^2-3t-1) - \tfrac{t+5}{39}(t-8) = 1
 \]
 gives the projector onto the common core
 \[
 P^{\mathrm{core}}_{\widehat{\mathfrak{sl}}_3}
 \;=\; -\tfrac{1}{39}(T+5I)(T-8I).
 \]
\item For $W_3$: the B\'{e}zout identity
 \[
 -\tfrac{1}{3}(t^2-3t-1) + \tfrac{t-2}{3}(t-1) = 1
 \]
 gives the projector
 \[
 P^{\mathrm{core}}_{W_3}
 \;=\; \tfrac{1}{3}(T-2I)(T-I).
 \]
\end{enumerate}
Each projected image has characteristic polynomial
$t^2-3t-1$, hence eigenvalues $(3 \pm \sqrt{13})/2$. The
projectors intertwine with any Drinfeld--Sokolov map whose induced
map on recurrence modules is shift-compatible and whose image
preserves the common quadratic sector.
\end{proposition}

\begin{proof}
Both B\'{e}zout identities follow by direct expansion.
The projectors are instances of the coprime-locus functional
calculus: if $\gcd(g,p/g)=1$ and $a(t)g(t)+b(t)(p/g)(t)=1$,
then $e_g(T)=b(T)(p/g)(T)$ is idempotent with image the
$g$-summand in the Chinese remainder decomposition
$k[t]/(p) \cong k[t]/(g) \oplus k[t]/(p/g)$.
\end{proof}

\begin{remark}[Why rank one is deceptive]
\label{rem:rank-one-deceptive}
For the $\widehat{\mathfrak{sl}}_2/\mathrm{Vir}$
correspondence, the recurrence polynomials coincide:
$p_{\widehat{\mathfrak{sl}}_2} = p_{\mathrm{Vir}}
 = (t-3)(t+1)$.
The common core is the entire module, both defect modules
vanish, and every divisor is coprime to its complement.
\emph{Nothing} in this rank-one picture foreshadows
the Jordan and defect phenomena that appear as soon as
the common factor becomes proper. The
$\widehat{\mathfrak{sl}}_3/W_3$ pair is the first genuine
example.
\end{remark}

\subsection{Geometric divisor-core lift data}
\label{subsec:casimir-frontier}
\index{divisor-core lattice!geometric lifts}

The algebra developed above corrects four overstatements
natural on first encounter:
\begin{enumerate}[label=\textup{(\arabic*)}]
\item The discriminant is \emph{not} the final transport
 object; it is the characteristic polynomial of the final
 object.
\item A divisor of the discriminant does \emph{not} in general
 produce a projector; it produces a canonical subquotient.
\item Repeated roots are \emph{not} a pathology to be deformed
 away; they are the source of genuine indecomposable
 transport sectors.
\item Drinfeld--Sokolov transport in rank $\geq 2$ is
 \emph{not} full operator similarity; the corrected
 invariant is the common divisor-core.
\end{enumerate}

The geometric lift is not an automatic consequence of the
discriminant. It is extra data.

\begin{definition}[Geometric divisor-core lift datum; \ClaimStatusDefinitional]
\label{def:geometric-divisor-core-lift-datum-inv}
\index{geometric divisor-core lift}
Let $\cA$ satisfy
Hypothesis~\ref{hyp:finite-character-recurrence-inv}. A
\emph{geometric divisor-core lift datum} for $\cA$ consists of a
finite-dimensional subquotient
\[
H^1_{\mathrm{red}}(\cA)
\]
of $H^1(\barBch(\cA))$, an endomorphism
\[
KS_\cA \in \operatorname{End}(H^1_{\mathrm{red}}(\cA)),
\]
and a $k[t]$-module isomorphism
\[
\eta_\cA\colon
\bigl(H^1_{\mathrm{red}}(\cA)\bigr)^\vee
\xrightarrow{\sim}
\mathcal R_\cA
\]
which intertwines $KS_\cA^\vee$ with the shift operator
$\sigma$ on $\mathcal R_\cA$.
\end{definition}

\begin{theorem}[Divisor-core consequences of lift data;
\ClaimStatusProvedHere]
\label{thm:geometric-lift-datum-consequences-inv}
Let $\cA$ satisfy
Hypothesis~\ref{hyp:finite-character-recurrence-inv}, and fix a
geometric divisor-core lift datum in the sense of
Definition~\ref{def:geometric-divisor-core-lift-datum-inv}. Then:
\begin{enumerate}[label=\textup{(\roman*)}]
\item for every divisor $S(x)\mid \Delta_\cA(x)$ normalized by
 $S(0)=1$, writing $s(t):=t^{\deg S}S(t^{-1})$, there is a
 canonical subquotient
 \[
 H^1_S(\cA)
 \]
 of $H^1_{\mathrm{red}}(\cA)$ whose dual is the divisor core
 \[
 \bigl(H^1_S(\cA)\bigr)^\vee \cong \mathcal C_s(p_\cA);
 \]
\item these sectors satisfy the divisor-lattice identities of
 Theorem~\ref{thm:divisor-core-calculus-inv}, up to canonical
 isomorphism of subquotients;
\item the obstruction to an idempotent splitting of $H^1_S(\cA)$
 is the class
 \[
 \varepsilon_S(\cA):=\varepsilon_s(p_\cA)\in
 \operatorname{Ext}^1_{k[t]}\!\bigl(k[t]/(p_\cA/s),k[t]/(s)\bigr)
 \cong k[t]/\gcd(s,p_\cA/s),
 \]
 and this class vanishes if and only if
 \[
 \gcd\!\left(S,\frac{\Delta_\cA}{S}\right)=1.
 \]
\end{enumerate}
\end{theorem}

\begin{proof}
The isomorphism $\eta_\cA$ identifies
$\bigl(H^1_{\mathrm{red}}(\cA)\bigr)^\vee$ with $M(p_\cA)$.
For $S\mid \Delta_\cA$, the submodule
$\mathcal C_s(p_\cA)\subset M(p_\cA)$ therefore defines the quotient
\[
H^1_{\mathrm{red}}(\cA)\twoheadrightarrow
H^1_S(\cA):=\bigl(\mathcal C_s(p_\cA)\bigr)^\vee.
\]
The divisor-lattice identities are dual to
Theorem~\ref{thm:divisor-core-calculus-inv}. The obstruction class
and its vanishing criterion are
Corollary~\ref{cor:repeated-roots-extension-data}, applied to
$g=s$ and $p=p_\cA$; reversal identifies
$\gcd(s,p_\cA/s)=1$ with
$\gcd(S,\Delta_\cA/S)=1$.
\end{proof}

\begin{proposition}[Primary quotient filtration from lift data;
\ClaimStatusProvedHere]
\label{prop:primary-quotient-filtration-lift-inv}
\index{primary quotient filtration}
Let $\cA$ carry a geometric divisor-core lift datum, let $q$ be an
irreducible factor of $p_\cA$, and let
$1\leq m\leq \nu_q(p_\cA)$. The primary divisor core
\[
\mathcal J_q^m(p_\cA)\subset \mathcal R_\cA
\]
defines a canonical quotient
\[
H^1_{\mathrm{red}}(\cA)\twoheadrightarrow
H^1_{q,m}(\cA):=
\bigl(\mathcal J_q^m(p_\cA)\bigr)^\vee.
\]
The quotients form an inverse filtration dual to
\[
0\subset
\mathcal J_q^1(p_\cA)\subset
\cdots\subset
\mathcal J_q^m(p_\cA).
\]
If $q=t-\lambda$ splits, the induced operator $KS_\cA-\lambda$ on
$H^1_{q,m}(\cA)$ is nilpotent of order $m$ and has one Jordan block
of size~$m$.
\end{proposition}

\begin{proof}
This is Theorem~\ref{thm:primary-jordan-filtration-inv} transported
through the isomorphism $\eta_\cA$ and then dualized.
\end{proof}

\begin{theorem}[Geometric common-core factorization;
\ClaimStatusProvedHere]
\label{thm:geometric-common-core-factorization-inv}
Let $\cA$ and $\mathcal B$ carry geometric divisor-core lift data.
Let $F$ be an exact factorization functor or reduction procedure
carrying $\cA$ to $\mathcal B$, and assume that the induced map on
the lift data is shift-compatible, equivalently that its dual is a
$k[t]$-linear map
\[
F^\vee_{\mathrm{red}}\colon
\mathcal R_\cA\longrightarrow \mathcal R_\mathcal B.
\]
Then $F^\vee_{\mathrm{red}}$ factors uniquely through the maximal
common core
\[
\mathcal R_\cA
\twoheadrightarrow
\mathcal C^{\max}_{\cA,\mathcal B}
\xrightarrow{\;m_u\;}
\mathcal C^{\max}_{\cA,\mathcal B}
\hookrightarrow
\mathcal R_\mathcal B
\]
for a unique $u\in\mathcal C^{\max}_{\cA,\mathcal B}$. Dually, the
reduced geometric transport factors through
\[
H^1_{\cA\wedge\mathcal B}:=
\bigl(\mathcal C^{\max}_{\cA,\mathcal B}\bigr)^\vee.
\]
\end{theorem}

\begin{proof}
Apply Proposition~\ref{prop:transport-factors-inv} to the
$k[t]$-linear map $F^\vee_{\mathrm{red}}$, then dualize the resulting
factorization through $\mathcal C^{\max}_{\cA,\mathcal B}$.
\end{proof}

\begin{theorem}[Drinfeld--Sokolov common-core transport under lift data;
\ClaimStatusProvedHere]
\label{thm:geometric-ds-common-core-inv}
Let
\[
\cA=\widehat{\mathfrak g}_k,\qquad
W=W_k(\mathfrak g),\qquad
S_{\mathfrak g,k}(x):=\gcd(\Delta_\cA(x),\Delta_W(x)),
\qquad
s_{\mathfrak g,k}(t):=
t^{\deg S_{\mathfrak g,k}}S_{\mathfrak g,k}(t^{-1}),
\]
with the gcd in $k[x]$ normalized by $S_{\mathfrak g,k}(0)=1$.
Assume $\cA$ and $W$ carry geometric divisor-core lift data and that
Drinfeld--Sokolov reduction induces a shift-compatible map on the
corresponding Casimir recurrence modules. Then the induced transport
factors through the reduced common-core sector
\[
H^1_{S_{\mathfrak g,k}}(-).
\]
For $\mathfrak{sl}_2$,
$S_{\mathfrak g,k}=(1-3x)(1+x)$, so the common core is the whole
discriminant-core object. On the $\mathfrak{sl}_3/W_3$ denominator
chart above,
\[
S_{\mathfrak g,k}(x)=1-3x-x^2,\qquad
s_{\mathfrak g,k}(t)=t^2-3t-1,
\]
and the two linear factors $t-8$ and $t-1$ are the defect modules.
\end{theorem}

\begin{proof}
The reversal of the common discriminant factor is
$\gcd(p_\cA,p_W)=s_{\mathfrak g,k}$. The factorization is
Theorem~\ref{thm:geometric-common-core-factorization-inv} with this
gcd. The $\mathfrak{sl}_2$ and $\mathfrak{sl}_3/W_3$ statements are
the displayed factorizations in the two examples above.
\end{proof}

\begin{remark}[Place in the characteristic hierarchy]
\label{rem:casimir-position}
The Casimir recurrence module $\mathcal{R}_\cA$ sits at the
$L_2$~level of the characteristic hierarchy
(cf.~the concordance, \S\ref{subsec:completion-kinematics-theory}): the
scalar level~$L_1$ is the modular
characteristic~$\kappa(\cA)$
(Theorem~\ref{thm:genus-universality}), while $L_2$ is the
spectral discriminant~$\Delta_\cA$ together with its module
structure. The passage from $\Delta_\cA$ as a polynomial to
$\mathcal{R}_\cA$ as a module is exactly the categorification
step that the shadow obstruction tower
(Definition~\ref{def:shadow-postnikov-tower}) anticipates at
the scalar level: the scalar shadow $\kappa$ is the decategorified
trace of $\Theta_\cA^{\leq 2}$, and the spectral sheet is the
decategorified trace of $\Theta_\cA^{\leq 3}$. The divisor-core
lattice lifts the trace to the module. The next layer, the
periodic table~$\Pi_\cA$ at $L_3$, should admit a similar
refinement from generating function to multilinear transport
object, but this lies beyond the divisor-core construction.
\end{remark}

\section{K3 reconstruction and the Borcherds trace}
\label{sec:bcinv-supplement}

\begin{remark}[K3 reconstruction and quadratic recognition;
\ClaimStatusConditional]
\label{rem:bcinv-inversion-K3}
Let $A_{\mathrm{K3}}$ be a constructed augmented chiral algebra
obtained from a specified K3-to-chiral functor.  The enhanced Ran
counit
\[
 \Omega_XB_X(A_{\mathrm{K3}})\xrightarrow{\sim}A_{\mathrm{K3}}
\]
is the universal reconstruction map of Theorem~A.  A quadratic K3
presentation supplies
$q_{A_{\mathrm{K3}}}\colon A_{\mathrm{K3}}^{\mathrm i}\to
B_X(A_{\mathrm{K3}})$; its comparison cone is the Theorem~B
obligation.  Under $H_{\mathbb D}^{\mathrm{bar}}(A_{\mathrm{K3}})$, Verdier duality
forms
$K_X(A_{\mathrm{K3}})=\mathbb D_{\Ran}B_X(A_{\mathrm{K3}})$.
An identification of this algebra with a Borcherds or Hall object
requires the corresponding Vol.~III comparison functor and
quasi-isomorphism.
\end{remark}

\begin{remark}[Borcherds and chiral scalar comparison;
\ClaimStatusConditional]
\label{rem:bcinv-phi-trace}
The K3 comparison has two independently defined scalars.  The
Borcherds product has
$\kappa_{\mathrm{BKM}}(\Delta_5)=c_\Lambda(0)/2=5$.  The proposed
Mukai chiral conductor has value $K^\kappa_{\mathcal B}=8$ after the
specified scalar trace.  A bridge between them consists of a
constructed pushforward
\[
 \Phi^{\mathrm{Borch}}_*(\mathfrak K_{A_{\mathrm{K3}}})
 \longrightarrow \mathfrak B_{\Delta_5}
\]
together with a theorem computing its effect on the two trace
functionals.  Bruinier's Heegner-divisor formula supplies the
automorphic Chern-class input; the chiral-to-automorphic comparison
map supplies the remaining obligation.
\end{remark}

\section{Lusztig specialization and K3 Hall inversion}
\label{sec:bcai-deep}

\begin{remark}[The formal branch $\hbar^2=-1/8$;
\ClaimStatusConditional]
\label{rem:bcai-deep-1}
Assume a recognized Hall source
$A_{\mathrm{Hall}}\simeq\mathbf H_{\Delta_5}$, a chosen scalar
normalization $K^\kappa_{\mathcal B}=8$, and the equation
$\hbar^2K^\kappa_{\mathcal B}=-1$.  These data select the formal branch
$\hbar^2=-1/8$.  Universal reconstruction gives
\[
 \Omega_XB_X(A_{\mathrm{Hall}})
 \xrightarrow{\sim}A_{\mathrm{Hall}}
\]
in the enhanced Ran ambient.  A strict Hall-complex representative at
this branch requires a chain model for the recognition map, continuous
bar and cobar functors, and a direct contraction of the comparison
cone.  Quadratic Koszulness requires the separate calculation of
$\operatorname{Cone}(q_{A_{\mathrm{Hall}}})$.
\end{remark}

\begin{remark}[Derived Hall realization problem;
\ClaimStatusOpen]
\label{rem:bcai-deep-2}
Construct a functor from the recognized compact Hall source to the
chosen complete chiral Ran category, together with a natural
identification of its bar construction and a strict representative of
the universal counit.  To obtain a Koszul algebra object, construct
the quadratic presentation and prove
$q_{A_{\mathrm{Hall}}}$ is a quasi-isomorphism.  The Hall and CoHA
results of Schiffmann--Vasserot, Davison, and
Kontsevich--Soibelman provide source-side structures for this
comparison.
\end{remark}

\begin{remark}[Normalization equation
$\hbar^2K^\kappa_{\mathcal B}=-1$; \ClaimStatusConditional]
\label{rem:bcai-deep-3}
Given the scalar normalization
$K^\kappa_{\mathcal B}=8$, the displayed equation selects
\[
 \hbar^2=-\frac18.
\]
A theorem identifying this formal branch with a Lusztig parameter,
Humbert monodromy order, or curvature cancellation must specify the
normalization map for each parameter and prove the three
identifications.  The cross-volume sources currently use distinct
normalizations, so this remark records the equation as a conditional
choice of branch.
\end{remark}
